\documentclass[10pt]{article}
% AUTO-GENERATED -- do not hand-edit.  generated by scripts/build_tables_document.py from the run started 2026-10-01T22:16:39 on code d82fd26 (full; 29,256 s; 7 workers)
% Every fragment of this folder, input unchanged, in name order; compile twice (list of tables).
\usepackage[T1]{fontenc}
\usepackage[utf8]{inputenc}
\usepackage[a4paper,margin=15mm]{geometry}
\usepackage{booktabs}
\usepackage{array}
\usepackage{amsmath}
\usepackage{amssymb}
\usepackage{adjustbox}
\usepackage{pdflscape}
\providecommand{\meth}[1]{\texttt{#1}}
\providecommand{\diag}[1]{\textsf{#1}}
\providecommand{\rhoV}{\rho_{\mathrm{v}}}
\providecommand{\rhoC}{\rho_{\mathrm{c}}}
\providecommand{\TE}{\textsc{te}}
\providecommand{\LE}{\textsc{le}}
\providecommand{\nobs}{n_{\mathrm{obs}}}
\providecommand{\nf}{n_f}
\setlength{\tabcolsep}{4pt}
\begin{document}
\section*{Generated tables of the sequence-acceleration benchmark}
Generated by scripts/build\_tables\_document.py from the run started 2026-10-01T22:16:39 on code d82fd26 (full; 29,256 s; 7 workers).  Each table is one fragment of \texttt{paper\_fragments/}, input unchanged and scaled to the page; the caption is the fragment's own description line, and the fragment's comment header names its sources and filters.  The index \texttt{README.md} maps every fragment to its sources.
\listoftables
\clearpage
\begin{landscape}
\begin{table}[p]
\centering
\caption{\texttt{f01\_trivial\_baseline.tex}: Trivial baseline (the paper's first result): the deployable trivials, the oracle constant as a labelled reference, the rank-1 and median accelerator, and how many accelerators beat predict-L\_hat / last\_value on more than half of their cells, with the hindsight (strict) skill alongside}
\label{tab:f01_trivial_baseline}
\begin{adjustbox}{max width=\linewidth, max totalheight=0.86\textheight}
% AUTO-GENERATED -- do not hand-edit.  generated by scripts/make_paper_tables.py at code d82fd26-dirty, results as of commit 2232d3d
% Trivial baseline (the paper's first result): the deployable trivials, the oracle constant as a labelled reference, the rank-1 and median accelerator, and how many accelerators beat predict-L_hat / last_value on more than half of their cells, with the hindsight (strict) skill alongside
% source : results/phase1/phase1_global.csv (core regimes, all three strata, capped cells excluded)
% source : results/phase1/phase1_global_holdout.csv (held-out regimes)
% filter : per stratum, core | held-out; win vs X = per-cell win rate (fraction of seeds with error below X's, mean over cells); skill vs X = median over cells of the per-cell median error ratio vs X; strict skill = hindsight best-of-four (err / best of the four deployable trivials on the same cell); the oracle is excluded from every count
% note   : the Prompt-4 'pooled median error below the best single trivial' count was dropped (Report-5A review): it compared pooled medians of different cells; the win rates are per-cell statements
\begin{tabular}{lrrrrrrrrrrrr}
\toprule
Stratum / row & \multicolumn{6}{c}{core (18 regimes)} & \multicolumn{6}{c}{held-out (6 regimes)} \\
\cmidrule(lr){2-7}\cmidrule(lr){8-13}
 & med.\ err & win vs $\hat L$ & skill vs $\hat L$ & win vs last & skill vs last & strict skill & med.\ err & win vs $\hat L$ & skill vs $\hat L$ & win vs last & skill vs last & strict skill \\
\midrule
\multicolumn{13}{l}{\textit{$g = 0.5$}} \\
predict $\hat L$ (\meth{constant\_assumed}) & 0.1094 & -- & -- & 0.040 & 1.973 & 2.313 & 0.0951 & -- & -- & 0.000 & 1.961 & 1.989 \\
\meth{last\_value} & 0.0690 & 0.961 & 0.507 & -- & -- & 1.021 & 0.0391 & 1.000 & 0.510 & -- & -- & 1.000 \\
\meth{window\_mean} & 0.1066 & 0.611 & 0.838 & 0.197 & 1.706 & 1.792 & 0.0982 & 0.561 & 0.974 & 0.000 & 1.766 & 1.905 \\
\meth{window\_min} & 0.0484 & 0.947 & 0.438 & 0.391 & 1.000 & 1.000 & 0.0386 & 1.000 & 0.503 & 0.354 & 1.000 & 1.000 \\
\meth{constant\_oracle} \textit{(reference, not deployable)} & 0.0579 & 1.000 & 0.494 & 0.728 & 0.995 & 1.000 & 0.0372 & 1.000 & 0.498 & 0.698 & 0.995 & 0.999 \\
\addlinespace[2pt]
best accelerator, core rank 1: \meth{log\_linear} & 0.0111 & 0.904 & 0.092 & 0.825 & 0.334 & 0.362 & 0.0121 & 0.972 & 0.119 & 0.787 & 0.298 & 0.323 \\
best accelerator, held-out rank 1: \meth{rational\_fit} & 0.0159 & 0.952 & 0.145 & 0.849 & 0.310 & 0.495 & 0.0062 & 0.900 & 0.062 & 0.830 & 0.196 & 0.206 \\
median accelerator (52) & 0.0494 & 0.891 & 0.452 & 0.557 & 0.999 & 1.043 & 0.0296 & 0.926 & 0.415 & 0.585 & 0.998 & 1.023 \\
accelerators with win rate $> 0.5$ / strict skill $< 1$ & & 46/52 & & 30/52 & & 18/52 & & 47/52 & & 35/52 & & 22/52 \\
\midrule
\multicolumn{13}{l}{\textit{$g = 0.1$} (headline)} \\
predict $\hat L$ (\meth{constant\_assumed}) & 0.0620 & -- & -- & 0.524 & 0.682 & 1.000 & 0.0561 & -- & -- & 0.569 & 0.627 & 1.000 \\
\meth{last\_value} & 0.1030 & 0.476 & 1.488 & -- & -- & 1.570 & 0.0862 & 0.431 & 1.616 & -- & -- & 1.632 \\
\meth{window\_mean} & 0.1463 & 0.381 & 2.071 & 0.158 & 1.427 & 2.733 & 0.1952 & 0.287 & 2.671 & 0.000 & 1.455 & 2.671 \\
\meth{window\_min} & 0.0681 & 0.503 & 1.304 & 0.348 & 1.000 & 1.493 & 0.0862 & 0.451 & 1.616 & 0.329 & 1.000 & 1.616 \\
\meth{constant\_oracle} \textit{(reference, not deployable)} & 0.0077 & 1.000 & 0.151 & 1.000 & 0.111 & 0.164 & 0.0092 & 1.000 & 0.173 & 1.000 & 0.111 & 0.174 \\
\addlinespace[2pt]
best accelerator, core rank 1: \meth{log\_linear} & 0.0202 & 0.809 & 0.299 & 0.756 & 0.332 & 0.601 & 0.0305 & 0.760 & 0.469 & 0.847 & 0.265 & 0.526 \\
best accelerator, held-out rank 1: \meth{rational\_fit} & 0.0296 & 0.787 & 0.328 & 0.877 & 0.299 & 0.670 & 0.0074 & 0.780 & 0.092 & 0.853 & 0.227 & 0.262 \\
median accelerator (52) & 0.0477 & 0.548 & 0.924 & 0.602 & 1.000 & 1.529 & 0.0476 & 0.537 & 1.353 & 0.657 & 0.994 & 1.458 \\
accelerators with win rate $> 0.5$ / strict skill $< 1$ & & 39/52 & & 32/52 & & 10/52 & & 35/52 & & 36/52 & & 15/52 \\
\midrule
\multicolumn{13}{l}{\textit{$g = 0.02$}} \\
predict $\hat L$ (\meth{constant\_assumed}) & 0.0541 & -- & -- & 0.544 & 0.558 & 1.000 & 0.0545 & -- & -- & 0.611 & 0.499 & 1.000 \\
\meth{last\_value} & 0.0970 & 0.456 & 1.871 & -- & -- & 1.873 & 0.0935 & 0.389 & 2.042 & -- & -- & 2.042 \\
\meth{window\_mean} & 0.1749 & 0.353 & 2.682 & 0.102 & 1.398 & 2.782 & 0.2077 & 0.267 & 3.468 & 0.000 & 1.418 & 3.468 \\
\meth{window\_min} & 0.0970 & 0.440 & 1.751 & 0.304 & 1.000 & 1.962 & 0.0935 & 0.407 & 2.042 & 0.329 & 1.000 & 2.042 \\
\meth{constant\_oracle} \textit{(reference, not deployable)} & 0.0011 & 1.000 & 0.031 & 1.000 & 0.020 & 0.031 & 0.0019 & 1.000 & 0.041 & 1.000 & 0.020 & 0.041 \\
\addlinespace[2pt]
best accelerator, core rank 1: \meth{richardson\_2} & 0.0218 & 0.777 & 0.319 & 0.847 & 0.394 & 0.883 & 0.0289 & 0.700 & 0.362 & 0.847 & 0.275 & 0.669 \\
best accelerator, held-out rank 1: \meth{rational\_fit} & 0.0228 & 0.761 & 0.398 & 0.843 & 0.278 & 0.605 & 0.0061 & 0.755 & 0.120 & 0.900 & 0.267 & 0.311 \\
median accelerator (52) & 0.0514 & 0.526 & 1.060 & 0.586 & 0.999 & 1.671 & 0.0522 & 0.479 & 1.720 & 0.640 & 0.994 & 1.764 \\
accelerators with win rate $> 0.5$ / strict skill $< 1$ & & 37/52 & & 31/52 & & 9/52 & & 23/52 & & 36/52 & & 17/52 \\
\bottomrule
\end{tabular}

\end{adjustbox}
\end{table}
\end{landscape}
\begin{landscape}
\begin{table}[p]
\centering
\caption{\texttt{f02\_ranking\_g0.02.tex}: Main ranking at g = 0.02: rank pool = accelerators above the validity floor, core and out-of-design side by side, unranked block appended}
\label{tab:f02_ranking_g0.02}
\begin{adjustbox}{max width=\linewidth, max totalheight=0.86\textheight}
% AUTO-GENERATED -- do not hand-edit.  generated by scripts/make_paper_tables.py from the run started 2026-10-01T22:16:39 on code d82fd26 (full; 29,256 s; 7 workers)
% Main ranking at g = 0.02: rank pool = accelerators above the validity floor, core and held-out side by side, unranked block appended
% source : results/phase1/phase1_global.csv (core regimes, all three strata, capped cells excluded)
% source : results/phase1/phase1_global_holdout.csv (held-out regimes)
% filter : target_g == 0.02; rank r = position by med_error among the 52 accelerators with rank_eligible == 1 (valid_rate >= 0.9), core and held-out ranked separately; trivial comparators not ranked (see f01); dagger = in the dangerous artifact
\begin{tabular}{rlllrrrrrrrrr}
\toprule
$r$ & Method & Family & T & \multicolumn{4}{c}{core} & \multicolumn{5}{c}{held-out} \\
\cmidrule(lr){5-8}\cmidrule(lr){9-13}
 & & & & med.\ err & med.\ skill & $\rhoC$ & $\rhoV$ & $r$ & med.\ err & med.\ skill & $\rhoC$ & $\rhoV$ \\
\midrule
1 & \meth{richardson\_2} & richardson & \TE & 0.0218 & 0.883 & 0.028 & 0.992 & 6 & 0.0289 & 0.669 & 0.016 & 0.987 \\
2 & \meth{richardson\_3} & richardson & \TE & 0.0228 & 0.860 & 0.045 & 0.986 & 3 & 0.0224 & 0.562 & 0.038 & 0.973 \\
3 & \meth{rational\_fit} & parametric & \TE & 0.0228 & 0.605 & 0.012 & 1.000 & 1 & 0.0061 & 0.311 & 0.002 & 1.000 \\
4 & \meth{log\_linear} & baseline & \TE & 0.0251 & 0.718 & 0.062 & 0.979 & 12 & 0.0360 & 0.734 & 0.009 & 1.000 \\
5 & \meth{richardson\_1} & richardson & \TE & 0.0264 & 0.801 & 0.033 & 1.000 & 5 & 0.0284 & 0.499 & 0.000 & 1.000 \\
6 & \meth{pade\_12} & pade & \TE & 0.0290 & 1.041 & 0.096 & 0.999 & 8 & 0.0331 & 0.753 & 0.093 & 0.993 \\
7 & \meth{pade\_34} & pade & \TE & 0.0299 & 1.060 & 0.108 & 0.997 & 7 & 0.0328 & 0.713 & 0.096 & 0.996 \\
8 & \meth{pade\_23} & pade & \TE & 0.0308 & 0.973 & 0.106 & 0.996 & 10 & 0.0332 & 0.809 & 0.082 & 0.996 \\
9 & \meth{pade\_45} & pade & \TE & 0.0308 & 0.992 & 0.098 & 0.998 & 9 & 0.0332 & 0.766 & 0.091 & 0.991 \\
10 & \meth{pade\_13} & pade & \TE & 0.0312 & 1.073 & 0.165 & 0.980 & 36 & 0.0677 & 1.751 & 0.122 & 0.991 \\
11 & \meth{pade\_33} & pade & \TE & 0.0321 & 1.314 & 0.167 & 0.973 & 2 & 0.0192 & 0.571 & 0.082 & 0.953 \\
12 & \meth{pade\_22} & pade & \TE & 0.0324 & 1.893 & 0.172 & 0.927 & 4 & 0.0272 & 0.699 & 0.071 & 0.947 \\
13 & \meth{double\_exp\_fit} & parametric & \TE & 0.0327 & 0.932 & 0.032 & 0.978 & 13 & 0.0366 & 0.903 & 0.002 & 1.000 \\
14 & \meth{single\_exp\_fit} & parametric & \TE & 0.0368 & 1.026 & 0.009 & 1.000 & 14 & 0.0383 & 1.382 & 0.002 & 1.000 \\
15 & \meth{richardson\_a05} & richardson & \TE & 0.0374 & 0.690 & 0.014 & 1.000 & 11 & 0.0339 & 0.757 & 0.000 & 1.000 \\
16 & \meth{richardson\_a10} & richardson & \TE & 0.0392 & 1.276 & 0.007 & 1.000 & 16 & 0.0411 & 0.907 & 0.000 & 1.000 \\
17 & \meth{wynn\_eps\_3} & wynn\_eps & \LE & 0.0456 & 1.606 & 0.060 & 0.954 & 24 & 0.0521 & 1.701 & 0.016 & 0.998 \\
18 & \meth{stability\_weighted} & ensemble & \TE & 0.0457 & 1.461 & 0.024 & 1.000 & 17 & 0.0474 & 1.424 & 0.002 & 1.000 \\
19 & \meth{pade\_44} & pade & \TE & 0.0463 & 1.810 & 0.119 & 0.975 & -- & 0.0241 & 0.614 & 0.156 & 0.891 \\
20 & \meth{brezinski\_theta2} & brezinski & \LE & 0.0477 & 2.067 & 0.092 & 0.975 & 27 & 0.0566 & 2.410 & 0.011 & 0.998 \\
21 & \meth{median\_ensemble} & ensemble & \TE & 0.0482 & 1.521 & 0.007 & 1.000 & 18 & 0.0477 & 1.685 & 0.000 & 1.000 \\
22 & \meth{pade\_11} & pade & \TE & 0.0487 & 1.718 & 0.050 & 0.999 & 15 & 0.0383 & 0.828 & 0.169 & 0.996 \\
23 & \meth{shanks\_4} & shanks & \LE & 0.0508 & 1.686 & 0.042 & 0.978 & 26 & 0.0530 & 1.847 & 0.004 & 1.000 \\
24 & \meth{wynn\_eps\_2} & wynn\_eps & \LE & 0.0513 & 1.595 & 0.050 & 0.965 & 22 & 0.0516 & 2.007 & 0.004 & 1.000 \\
25 & \meth{levin\_t2} & levin & \LE & 0.0513 & 1.602 & 0.042 & 0.977 & 20 & 0.0512 & 1.867 & 0.004 & 1.000 \\
26 & \meth{levin\_u2} & levin & \LE & 0.0513 & 1.868 & 0.040 & 0.978 & 33 & 0.0617 & 1.864 & 0.004 & 1.000 \\
27 & \meth{brezinski\_theta1} & brezinski & \LE & 0.0514 & 1.900 & 0.042 & 1.000 & 28 & 0.0588 & 1.566 & 0.002 & 1.000 \\
28 & \meth{shanks\_2} & shanks & \LE & 0.0514 & 1.928 & 0.020 & 0.999 & 23 & 0.0518 & 1.984 & 0.000 & 1.000 \\
29 & \meth{shanks\_3} & shanks & \LE & 0.0514 & 1.585 & 0.019 & 1.000 & 21 & 0.0514 & 1.847 & 0.002 & 0.998 \\
30 & \meth{levin\_v1} & levin & \LE & 0.0515 & 1.766 & 0.045 & 0.978 & 35 & 0.0634 & 1.951 & 0.000 & 1.000 \\
31 & \meth{levin\_t1} & levin & \LE & 0.0523 & 1.654 & 0.037 & 0.978 & 31 & 0.0601 & 1.764 & 0.002 & 1.000 \\
32 & \meth{wynn\_eps\_1} & wynn\_eps & \LE & 0.0523 & 1.654 & 0.037 & 0.978 & 32 & 0.0601 & 1.764 & 0.002 & 1.000 \\
33 & \meth{wynn\_rho\_1} & wynn\_rho & \LE & 0.0527 & 1.657 & 0.037 & 0.977 & 19 & 0.0510 & 1.845 & 0.002 & 1.000 \\
34 & \meth{best\_shanks\_wynn} & ensemble & \LE & 0.0531 & 1.720 & 0.015 & 1.000 & 29 & 0.0601 & 1.764 & 0.002 & 1.000 \\
35 & \meth{shanks\_1} & shanks & \LE & 0.0531 & 1.720 & 0.015 & 1.000 & 30 & 0.0601 & 1.764 & 0.002 & 1.000 \\
36 & \meth{levin\_v2} & levin & \LE & 0.0531 & 1.898 & 0.046 & 0.977 & 25 & 0.0524 & 1.886 & 0.000 & 1.000 \\
37 & \meth{wynn\_rho\_3} & wynn\_rho & \LE & 0.0545 & 1.653 & 0.073 & 0.971 & 42 & 0.1099 & 1.853 & 0.009 & 0.996 \\
38 & \meth{wynn\_rho\_2} & wynn\_rho & \LE & 0.0549 & 1.871 & 0.056 & 0.976 & 34 & 0.0629 & 1.929 & 0.084 & 0.931 \\
39 & \meth{log\_fit} & parametric & \TE & 0.0709 & 1.634 & 0.014 & 1.000 & 41 & 0.0962 & 1.827 & 0.000 & 1.000 \\
40 & \meth{levin\_u1} & levin & \LE & 0.0839 & 1.924 & 0.081 & 0.978 & 37 & 0.0925 & 1.885 & 0.000 & 1.000 \\
41 & \meth{anderson\_2} & anderson & \LE & 0.0958 & 1.904 & 0.012 & 1.000 & 40 & 0.0950 & 2.110 & 0.000 & 1.000 \\
42 & \meth{anderson\_3} & anderson & \LE & 0.0963 & 1.895 & 0.012 & 1.000 & 39 & 0.0943 & 2.026 & 0.000 & 1.000 \\
43 & \meth{anderson\_1} & anderson & \LE & 0.0969 & 1.921 & 0.015 & 1.000 & 38 & 0.0926 & 1.976 & 0.000 & 1.000 \\
44 & \meth{richardson\_a20} & richardson & \TE & 0.0970 & 1.854 & 0.005 & 1.000 & 43 & 0.1365 & 2.571 & 0.000 & 1.000 \\
\midrule
\multicolumn{13}{l}{\textit{Unranked: below the validity floor $\rhoV < 0.9$ on the core regimes (8 methods; shown, never ranked)}} \\
-- & $\dagger$\meth{neville\_2} & neville & \TE & 6.1997 & 256.914 & 0.726 & 0.646 & -- & 5.7956 & 161.967 & 0.653 & 0.720 \\
-- & $\dagger$\meth{pade\_31} & pade & \TE & 1.8859 & 37.209 & 0.839 & 0.626 & -- & 4.7879 & 70.359 & 1.000 & 0.736 \\
-- & $\dagger$\meth{geom\_avg\_diff} & baseline & \LE & 0.0533 & 1.974 & 0.450 & 0.596 & -- & 0.0615 & 1.878 & 0.331 & 0.678 \\
-- & $\dagger$\meth{pade\_32} & pade & \TE & 0.2729 & 8.577 & 0.687 & 0.590 & -- & 0.1400 & 2.802 & 0.436 & 0.667 \\
-- & $\dagger$\meth{neville\_3} & neville & \TE & 160.8795 & 4146.730 & 0.778 & 0.453 & -- & 167.5300 & 4246.912 & 0.667 & 0.531 \\
-- & $\dagger$\meth{pade\_21} & pade & \TE & 0.2148 & 8.100 & 0.716 & 0.450 & -- & 0.4278 & 5.814 & 0.471 & 0.609 \\
-- & $\dagger$\meth{linear} & baseline & \TE & 0.4055 & 6.227 & 0.888 & 0.296 & -- & 0.6270 & 10.306 & 0.904 & 0.293 \\
-- & $\dagger$\meth{neville\_4} & neville & \TE & 0.2840 & 6.281 & 0.756 & 0.286 & -- & 0.0266 & 0.773 & 0.667 & 0.360 \\
\bottomrule
\end{tabular}

\end{adjustbox}
\end{table}
\end{landscape}
\begin{landscape}
\begin{table}[p]
\centering
\caption{\texttt{f02\_ranking\_g0.1.tex}: Main ranking at g = 0.1: rank pool = accelerators above the validity floor, core and out-of-design side by side, unranked block appended}
\label{tab:f02_ranking_g0.1}
\begin{adjustbox}{max width=\linewidth, max totalheight=0.86\textheight}
% AUTO-GENERATED -- do not hand-edit.  generated by scripts/make_paper_tables.py at code d82fd26-dirty, results as of commit 2232d3d
% Main ranking at g = 0.1: rank pool = accelerators above the validity floor, core and held-out side by side, unranked block appended
% source : results/phase1/phase1_global.csv (core regimes, all three strata, capped cells excluded)
% source : results/phase1/phase1_global_holdout.csv (held-out regimes)
% filter : target_g == 0.1; rank r = position by med_error among the 52 accelerators with rank_eligible == 1 (valid_rate >= 0.9), core and held-out ranked separately; trivial comparators not ranked (see f01); dagger = in the dangerous artifact
\begin{tabular}{rlllrrrrrrrrr}
\toprule
$r$ & Method & Family & T & \multicolumn{4}{c}{core} & \multicolumn{5}{c}{held-out} \\
\cmidrule(lr){5-8}\cmidrule(lr){9-13}
 & & & & med.\ err & med.\ skill & $\rhoC$ & $\rhoV$ & $r$ & med.\ err & med.\ skill & $\rhoC$ & $\rhoV$ \\
\midrule
1 & \meth{log\_linear} & baseline & \TE & 0.0202 & 0.601 & 0.112 & 0.953 & 10 & 0.0305 & 0.526 & 0.000 & 1.000 \\
2 & \meth{richardson\_3} & richardson & \TE & 0.0222 & 0.827 & 0.097 & 0.980 & 8 & 0.0260 & 0.426 & 0.035 & 0.973 \\
3 & \meth{richardson\_2} & richardson & \TE & 0.0257 & 0.696 & 0.085 & 0.992 & 11 & 0.0316 & 0.479 & 0.013 & 0.987 \\
4 & \meth{double\_exp\_fit} & parametric & \TE & 0.0287 & 0.764 & 0.029 & 0.979 & 13 & 0.0322 & 0.804 & 0.002 & 1.000 \\
5 & \meth{pade\_34} & pade & \TE & 0.0291 & 1.008 & 0.095 & 0.992 & 3 & 0.0211 & 0.517 & 0.087 & 0.991 \\
6 & \meth{rational\_fit} & parametric & \TE & 0.0296 & 0.670 & 0.009 & 1.000 & 1 & 0.0074 & 0.262 & 0.002 & 1.000 \\
7 & \meth{richardson\_1} & richardson & \TE & 0.0301 & 0.879 & 0.088 & 1.000 & 12 & 0.0318 & 0.310 & 0.000 & 1.000 \\
8 & \meth{pade\_23} & pade & \TE & 0.0302 & 0.889 & 0.093 & 0.996 & 6 & 0.0243 & 0.482 & 0.082 & 0.993 \\
9 & \meth{pade\_12} & pade & \TE & 0.0314 & 0.903 & 0.087 & 0.995 & 5 & 0.0225 & 0.544 & 0.084 & 0.991 \\
10 & \meth{pade\_45} & pade & \TE & 0.0318 & 0.973 & 0.085 & 0.994 & 7 & 0.0253 & 0.502 & 0.089 & 0.998 \\
11 & \meth{pade\_33} & pade & \TE & 0.0329 & 1.334 & 0.139 & 0.995 & 2 & 0.0153 & 0.415 & 0.080 & 0.960 \\
12 & \meth{pade\_22} & pade & \TE & 0.0333 & 1.904 & 0.163 & 0.924 & 4 & 0.0217 & 0.478 & 0.071 & 0.947 \\
13 & \meth{single\_exp\_fit} & parametric & \TE & 0.0361 & 1.001 & 0.008 & 1.000 & 15 & 0.0342 & 1.080 & 0.002 & 1.000 \\
14 & \meth{richardson\_a05} & richardson & \TE & 0.0384 & 0.674 & 0.009 & 1.000 & 9 & 0.0267 & 0.564 & 0.000 & 1.000 \\
15 & \meth{pade\_13} & pade & \TE & 0.0390 & 1.303 & 0.151 & 0.972 & 33 & 0.0563 & 1.700 & 0.158 & 0.987 \\
16 & \meth{pade\_11} & pade & \TE & 0.0397 & 1.558 & 0.049 & 0.999 & 16 & 0.0356 & 0.578 & 0.000 & 1.000 \\
17 & \meth{pade\_44} & pade & \TE & 0.0410 & 1.947 & 0.111 & 0.976 & 14 & 0.0338 & 0.472 & 0.111 & 0.944 \\
18 & \meth{wynn\_eps\_3} & wynn\_eps & \LE & 0.0438 & 1.518 & 0.057 & 0.957 & 29 & 0.0480 & 1.569 & 0.016 & 0.998 \\
19 & \meth{stability\_weighted} & ensemble & \TE & 0.0456 & 1.362 & 0.019 & 1.000 & 17 & 0.0423 & 1.453 & 0.002 & 1.000 \\
20 & \meth{median\_ensemble} & ensemble & \TE & 0.0473 & 1.397 & 0.006 & 1.000 & 18 & 0.0436 & 1.370 & 0.000 & 1.000 \\
21 & \meth{wynn\_eps\_2} & wynn\_eps & \LE & 0.0473 & 1.521 & 0.048 & 0.967 & 23 & 0.0476 & 1.679 & 0.004 & 1.000 \\
22 & \meth{levin\_t2} & levin & \LE & 0.0475 & 1.520 & 0.040 & 0.979 & 21 & 0.0471 & 1.434 & 0.004 & 1.000 \\
23 & \meth{shanks\_4} & shanks & \LE & 0.0475 & 1.581 & 0.041 & 0.979 & 31 & 0.0490 & 1.588 & 0.004 & 1.000 \\
24 & \meth{levin\_u2} & levin & \LE & 0.0475 & 1.592 & 0.038 & 0.979 & 36 & 0.0742 & 1.663 & 0.004 & 1.000 \\
25 & \meth{levin\_v1} & levin & \LE & 0.0476 & 1.571 & 0.042 & 0.979 & 37 & 0.0759 & 1.696 & 0.000 & 1.000 \\
26 & \meth{shanks\_3} & shanks & \LE & 0.0476 & 1.516 & 0.018 & 1.000 & 22 & 0.0473 & 1.554 & 0.002 & 0.998 \\
27 & \meth{brezinski\_theta2} & brezinski & \LE & 0.0478 & 1.806 & 0.087 & 0.976 & 32 & 0.0526 & 1.835 & 0.011 & 0.998 \\
28 & \meth{shanks\_2} & shanks & \LE & 0.0479 & 1.672 & 0.019 & 0.999 & 28 & 0.0478 & 1.861 & 0.002 & 1.000 \\
29 & \meth{wynn\_rho\_1} & wynn\_rho & \LE & 0.0489 & 1.528 & 0.035 & 0.979 & 20 & 0.0469 & 1.719 & 0.002 & 1.000 \\
30 & \meth{wynn\_rho\_3} & wynn\_rho & \LE & 0.0502 & 1.546 & 0.068 & 0.973 & 43 & 0.1044 & 1.937 & 0.011 & 0.996 \\
31 & \meth{levin\_v2} & levin & \LE & 0.0503 & 1.680 & 0.043 & 0.979 & 30 & 0.0483 & 1.666 & 0.002 & 1.000 \\
32 & \meth{levin\_t1} & levin & \LE & 0.0505 & 1.523 & 0.035 & 0.979 & 26 & 0.0476 & 1.458 & 0.002 & 1.000 \\
33 & \meth{wynn\_eps\_1} & wynn\_eps & \LE & 0.0505 & 1.523 & 0.035 & 0.979 & 27 & 0.0476 & 1.458 & 0.002 & 1.000 \\
34 & \meth{wynn\_rho\_2} & wynn\_rho & \LE & 0.0510 & 1.651 & 0.053 & 0.978 & 35 & 0.0628 & 1.552 & 0.087 & 0.931 \\
35 & \meth{richardson\_a10} & richardson & \TE & 0.0561 & 1.064 & 0.006 & 1.000 & 34 & 0.0614 & 1.164 & 0.000 & 1.000 \\
36 & \meth{brezinski\_theta1} & brezinski & \LE & 0.0584 & 1.518 & 0.040 & 1.000 & 19 & 0.0465 & 1.412 & 0.002 & 1.000 \\
37 & \meth{best\_shanks\_wynn} & ensemble & \LE & 0.0589 & 1.538 & 0.014 & 1.000 & 24 & 0.0476 & 1.458 & 0.002 & 1.000 \\
38 & \meth{shanks\_1} & shanks & \LE & 0.0589 & 1.538 & 0.014 & 1.000 & 25 & 0.0476 & 1.458 & 0.002 & 1.000 \\
39 & \meth{log\_fit} & parametric & \TE & 0.0759 & 1.566 & 0.010 & 1.000 & 42 & 0.1008 & 1.499 & 0.000 & 1.000 \\
40 & \meth{levin\_u1} & levin & \LE & 0.0863 & 1.603 & 0.076 & 0.979 & 41 & 0.0889 & 1.769 & 0.000 & 1.000 \\
41 & \meth{anderson\_2} & anderson & \LE & 0.1036 & 1.530 & 0.011 & 1.000 & 40 & 0.0877 & 1.754 & 0.000 & 1.000 \\
42 & \meth{anderson\_3} & anderson & \LE & 0.1038 & 1.559 & 0.011 & 1.000 & 39 & 0.0870 & 1.681 & 0.000 & 1.000 \\
43 & \meth{anderson\_1} & anderson & \LE & 0.1039 & 1.573 & 0.014 & 1.000 & 38 & 0.0853 & 1.785 & 0.000 & 1.000 \\
44 & \meth{richardson\_a20} & richardson & \TE & 0.1113 & 1.477 & 0.004 & 1.000 & 44 & 0.1378 & 1.939 & 0.000 & 1.000 \\
\midrule
\multicolumn{13}{l}{\textit{Unranked: below the validity floor $\rhoV < 0.9$ on the core regimes (8 methods; shown, never ranked)}} \\
-- & $\dagger$\meth{pade\_31} & pade & \TE & 1.9535 & 31.176 & 0.773 & 0.867 & 45 & 2.8890 & 44.868 & 0.760 & 0.929 \\
-- & $\dagger$\meth{pade\_32} & pade & \TE & 0.2574 & 5.547 & 0.582 & 0.688 & -- & 0.0680 & 1.359 & 0.336 & 0.669 \\
-- & $\dagger$\meth{neville\_2} & neville & \TE & 4.5269 & 192.499 & 0.719 & 0.661 & -- & 4.9545 & 64.141 & 0.644 & 0.727 \\
-- & $\dagger$\meth{pade\_21} & pade & \TE & 0.3364 & 5.161 & 0.605 & 0.602 & -- & 0.1909 & 4.083 & 0.358 & 0.682 \\
-- & $\dagger$\meth{geom\_avg\_diff} & baseline & \LE & 0.0505 & 1.672 & 0.442 & 0.601 & -- & 0.0503 & 1.430 & 0.331 & 0.680 \\
-- & $\dagger$\meth{neville\_3} & neville & \TE & 111.0798 & 3262.867 & 0.771 & 0.480 & -- & 137.1949 & 2636.988 & 0.667 & 0.553 \\
-- & $\dagger$\meth{linear} & baseline & \TE & 0.3057 & 6.048 & 0.765 & 0.444 & -- & 0.5665 & 6.706 & 0.700 & 0.524 \\
-- & $\dagger$\meth{neville\_4} & neville & \TE & 96.0674 & 1742.476 & 0.750 & 0.335 & -- & 0.1069 & 1.564 & 0.667 & 0.382 \\
\bottomrule
\end{tabular}

\end{adjustbox}
\end{table}
\end{landscape}
\begin{landscape}
\begin{table}[p]
\centering
\caption{\texttt{f02\_ranking\_g0.5.tex}: Main ranking at g = 0.5: rank pool = accelerators above the validity floor, core and out-of-design side by side, unranked block appended}
\label{tab:f02_ranking_g0.5}
\begin{adjustbox}{max width=\linewidth, max totalheight=0.86\textheight}
% AUTO-GENERATED -- do not hand-edit.  generated by scripts/make_paper_tables.py at code d82fd26-dirty, results as of commit 2232d3d
% Main ranking at g = 0.5: rank pool = accelerators above the validity floor, core and held-out side by side, unranked block appended
% source : results/phase1/phase1_global.csv (core regimes, all three strata, capped cells excluded)
% source : results/phase1/phase1_global_holdout.csv (held-out regimes)
% filter : target_g == 0.5; rank r = position by med_error among the 52 accelerators with rank_eligible == 1 (valid_rate >= 0.9), core and held-out ranked separately; trivial comparators not ranked (see f01); dagger = in the dangerous artifact
\begin{tabular}{rlllrrrrrrrrr}
\toprule
$r$ & Method & Family & T & \multicolumn{4}{c}{core} & \multicolumn{5}{c}{held-out} \\
\cmidrule(lr){5-8}\cmidrule(lr){9-13}
 & & & & med.\ err & med.\ skill & $\rhoC$ & $\rhoV$ & $r$ & med.\ err & med.\ skill & $\rhoC$ & $\rhoV$ \\
\midrule
1 & \meth{log\_linear} & baseline & \TE & 0.0111 & 0.362 & 0.091 & 0.982 & 8 & 0.0121 & 0.323 & 0.004 & 1.000 \\
2 & \meth{richardson\_3} & richardson & \TE & 0.0116 & 0.464 & 0.072 & 0.982 & 5 & 0.0100 & 0.185 & 0.030 & 0.972 \\
3 & \meth{pade\_11} & pade & \TE & 0.0132 & 0.469 & 0.113 & 0.991 & 12 & 0.0137 & 0.185 & 0.033 & 0.998 \\
4 & \meth{pade\_22} & pade & \TE & 0.0138 & 0.541 & 0.152 & 0.948 & 2 & 0.0067 & 0.203 & 0.065 & 0.987 \\
5 & \meth{pade\_33} & pade & \TE & 0.0143 & 0.397 & 0.151 & 0.994 & 3 & 0.0086 & 0.205 & 0.065 & 0.994 \\
6 & \meth{pade\_44} & pade & \TE & 0.0153 & 0.706 & 0.084 & 0.990 & 4 & 0.0092 & 0.199 & 0.054 & 0.996 \\
7 & \meth{richardson\_2} & richardson & \TE & 0.0155 & 0.444 & 0.075 & 0.993 & 7 & 0.0119 & 0.201 & 0.011 & 0.989 \\
8 & \meth{rational\_fit} & parametric & \TE & 0.0159 & 0.495 & 0.012 & 1.000 & 1 & 0.0062 & 0.206 & 0.000 & 1.000 \\
9 & \meth{double\_exp\_fit} & parametric & \TE & 0.0165 & 0.447 & 0.024 & 0.982 & 14 & 0.0159 & 0.606 & 0.004 & 1.000 \\
10 & \meth{richardson\_a05} & richardson & \TE & 0.0191 & 0.347 & 0.011 & 1.000 & 17 & 0.0205 & 0.346 & 0.004 & 1.000 \\
11 & \meth{single\_exp\_fit} & parametric & \TE & 0.0207 & 0.629 & 0.005 & 1.000 & 16 & 0.0204 & 0.654 & 0.004 & 1.000 \\
12 & \meth{pade\_23} & pade & \TE & 0.0211 & 0.693 & 0.119 & 0.980 & 10 & 0.0128 & 0.297 & 0.061 & 0.994 \\
13 & \meth{pade\_12} & pade & \TE & 0.0214 & 0.833 & 0.130 & 0.963 & 11 & 0.0132 & 0.292 & 0.056 & 0.998 \\
14 & \meth{richardson\_1} & richardson & \TE & 0.0218 & 0.456 & 0.030 & 1.000 & 6 & 0.0111 & 0.218 & 0.004 & 1.000 \\
15 & \meth{pade\_45} & pade & \TE & 0.0229 & 0.975 & 0.099 & 0.986 & 9 & 0.0125 & 0.280 & 0.069 & 0.989 \\
16 & \meth{pade\_34} & pade & \TE & 0.0244 & 0.888 & 0.126 & 0.980 & 13 & 0.0144 & 0.281 & 0.067 & 0.998 \\
17 & \meth{pade\_13} & pade & \TE & 0.0365 & 1.207 & 0.096 & 0.990 & 15 & 0.0173 & 0.361 & 0.189 & 0.989 \\
18 & \meth{richardson\_a10} & richardson & \TE & 0.0409 & 0.614 & 0.005 & 1.000 & 25 & 0.0292 & 0.604 & 0.000 & 1.000 \\
19 & \meth{median\_ensemble} & ensemble & \TE & 0.0435 & 1.032 & 0.007 & 1.000 & 23 & 0.0265 & 1.023 & 0.006 & 1.000 \\
20 & \meth{brezinski\_theta1} & brezinski & \LE & 0.0482 & 1.018 & 0.037 & 1.000 & 18 & 0.0248 & 0.987 & 0.007 & 1.000 \\
21 & \meth{levin\_t1} & levin & \LE & 0.0490 & 1.040 & 0.032 & 0.982 & 20 & 0.0265 & 1.027 & 0.007 & 1.000 \\
22 & \meth{wynn\_eps\_1} & wynn\_eps & \LE & 0.0490 & 1.040 & 0.032 & 0.982 & 22 & 0.0265 & 1.027 & 0.007 & 1.000 \\
23 & \meth{wynn\_eps\_2} & wynn\_eps & \LE & 0.0490 & 1.078 & 0.048 & 0.970 & 36 & 0.0363 & 1.064 & 0.009 & 1.000 \\
24 & \meth{log\_fit} & parametric & \TE & 0.0493 & 0.969 & 0.048 & 1.000 & 45 & 0.0563 & 0.940 & 0.007 & 1.000 \\
25 & \meth{stability\_weighted} & ensemble & \TE & 0.0493 & 1.047 & 0.019 & 1.000 & 34 & 0.0355 & 1.024 & 0.017 & 1.000 \\
26 & \meth{best\_shanks\_wynn} & ensemble & \LE & 0.0494 & 1.043 & 0.014 & 1.000 & 19 & 0.0265 & 1.027 & 0.007 & 1.000 \\
27 & \meth{shanks\_1} & shanks & \LE & 0.0494 & 1.043 & 0.014 & 1.000 & 21 & 0.0265 & 1.027 & 0.007 & 1.000 \\
28 & \meth{levin\_t2} & levin & \LE & 0.0497 & 1.051 & 0.040 & 0.981 & 31 & 0.0349 & 1.032 & 0.017 & 1.000 \\
29 & \meth{wynn\_eps\_3} & wynn\_eps & \LE & 0.0507 & 1.080 & 0.057 & 0.961 & 24 & 0.0284 & 1.058 & 0.022 & 0.998 \\
30 & \meth{shanks\_2} & shanks & \LE & 0.0541 & 1.091 & 0.022 & 0.999 & 38 & 0.0394 & 1.079 & 0.011 & 1.000 \\
31 & \meth{shanks\_3} & shanks & \LE & 0.0553 & 1.093 & 0.017 & 1.000 & 29 & 0.0320 & 1.062 & 0.007 & 0.998 \\
32 & \meth{shanks\_4} & shanks & \LE & 0.0570 & 1.098 & 0.036 & 0.982 & 28 & 0.0306 & 1.094 & 0.009 & 1.000 \\
33 & \meth{brezinski\_theta2} & brezinski & \LE & 0.0579 & 1.140 & 0.088 & 0.979 & 43 & 0.0439 & 1.144 & 0.030 & 0.998 \\
34 & \meth{wynn\_rho\_1} & wynn\_rho & \LE & 0.0598 & 1.060 & 0.032 & 0.981 & 37 & 0.0378 & 1.039 & 0.009 & 1.000 \\
35 & \meth{wynn\_rho\_2} & wynn\_rho & \LE & 0.0618 & 1.078 & 0.052 & 0.980 & 26 & 0.0299 & 1.065 & 0.085 & 0.943 \\
36 & \meth{levin\_v2} & levin & \LE & 0.0657 & 1.094 & 0.041 & 0.981 & 42 & 0.0434 & 1.092 & 0.009 & 1.000 \\
37 & \meth{levin\_u2} & levin & \LE & 0.0658 & 1.068 & 0.057 & 0.982 & 30 & 0.0325 & 1.039 & 0.017 & 1.000 \\
38 & \meth{wynn\_rho\_3} & wynn\_rho & \LE & 0.0660 & 1.092 & 0.071 & 0.975 & 35 & 0.0357 & 1.066 & 0.033 & 0.996 \\
39 & \meth{levin\_u1} & levin & \LE & 0.0665 & 1.079 & 0.069 & 0.982 & 39 & 0.0395 & 1.030 & 0.007 & 1.000 \\
40 & \meth{anderson\_3} & anderson & \LE & 0.0669 & 1.044 & 0.010 & 1.000 & 40 & 0.0404 & 1.021 & 0.006 & 1.000 \\
41 & \meth{levin\_v1} & levin & \LE & 0.0673 & 1.096 & 0.059 & 0.982 & 27 & 0.0300 & 1.071 & 0.006 & 1.000 \\
42 & \meth{anderson\_2} & anderson & \LE & 0.0679 & 1.038 & 0.009 & 1.000 & 41 & 0.0407 & 1.022 & 0.004 & 1.000 \\
43 & \meth{anderson\_1} & anderson & \LE & 0.0691 & 1.043 & 0.014 & 1.000 & 32 & 0.0351 & 1.016 & 0.006 & 1.000 \\
44 & \meth{richardson\_a20} & richardson & \TE & 0.0754 & 1.048 & 0.006 & 1.000 & 44 & 0.0543 & 1.019 & 0.006 & 1.000 \\
45 & $\dagger$\meth{pade\_32} & pade & \TE & 0.1046 & 2.932 & 0.338 & 0.936 & 33 & 0.0351 & 0.695 & 0.265 & 0.969 \\
46 & $\dagger$\meth{linear} & baseline & \TE & 0.1925 & 3.105 & 0.298 & 0.944 & 47 & 0.1335 & 3.318 & 0.261 & 0.995 \\
47 & $\dagger$\meth{pade\_21} & pade & \TE & 0.1934 & 2.711 & 0.352 & 0.922 & 46 & 0.0603 & 1.198 & 0.233 & 0.993 \\
\midrule
\multicolumn{13}{l}{\textit{Unranked: below the validity floor $\rhoV < 0.9$ on the core regimes (5 methods; shown, never ranked)}} \\
-- & $\dagger$\meth{pade\_31} & pade & \TE & 0.2006 & 5.674 & 0.526 & 0.889 & 48 & 0.1919 & 3.782 & 0.422 & 0.993 \\
-- & $\dagger$\meth{neville\_2} & neville & \TE & 1.2618 & 61.826 & 0.719 & 0.706 & -- & 2.2715 & 31.955 & 0.617 & 0.761 \\
-- & $\dagger$\meth{geom\_avg\_diff} & baseline & \LE & 0.0301 & 1.043 & 0.456 & 0.578 & -- & 0.0271 & 0.995 & 0.343 & 0.676 \\
-- & $\dagger$\meth{neville\_3} & neville & \TE & 27.4006 & 1435.589 & 0.759 & 0.563 & -- & 62.0249 & 940.748 & 0.661 & 0.654 \\
-- & $\dagger$\meth{neville\_4} & neville & \TE & 12.1425 & 2206.764 & 0.741 & 0.403 & -- & 115.7123 & 1768.344 & 0.667 & 0.446 \\
\bottomrule
\end{tabular}

\end{adjustbox}
\end{table}
\end{landscape}
\begin{landscape}
\begin{table}[p]
\centering
\caption{\texttt{f03\_skill\_summary.tex}: Skill summary: fraction of (method x curve family x noise) cells whose median skill is below 1 (the method beat the hindsight best-of-four deployable trivial (strict) on the median seed), per family, per stratum, core vs out-of-design}
\label{tab:f03_skill_summary}
\begin{adjustbox}{max width=\linewidth, max totalheight=0.86\textheight}
% AUTO-GENERATED -- do not hand-edit.  generated by scripts/make_paper_tables.py at code d82fd26-dirty, results as of commit 2232d3d
% Skill summary: fraction of (method x regime x noise) cells whose median skill is below 1 (the method beat the hindsight best-of-four deployable trivial (strict) on the median seed), per family, per stratum, core vs held-out
% source : results/phase1/phase1_aggregated.csv (per (method, regime, noise, g) cell: median skill over 30 seeds)
% filter : capped == 0, is_oracle == 0; a cell counts when its median-seed skill is < 1 (NaN median skill counts as not < 1)
% note   : K = methods in the family; T = method type (\TE trajectory extrapolator / \LE limit estimator)
\begin{tabular}{lrlrrrrrr}
\toprule
Family & $K$ & T & \multicolumn{2}{c}{$g = 0.5$} & \multicolumn{2}{c}{$g = 0.1$} & \multicolumn{2}{c}{$g = 0.02$} \\
\cmidrule(lr){4-5}\cmidrule(lr){6-7}\cmidrule(lr){8-9}
 & & & core & held-out & core & held-out & core & held-out \\
\midrule
richardson & 6 & \TE & 67\% & 72\% & 49\% & 59\% & 49\% & 61\% \\
parametric & 4 & \TE & 72\% & 82\% & 54\% & 55\% & 47\% & 45\% \\
pade & 12 & \TE & 49\% & 68\% & 37\% & 54\% & 33\% & 45\% \\
neville & 3 & \TE & 24\% & 33\% & 19\% & 29\% & 19\% & 29\% \\
baseline & 3 & mixed & 40\% & 48\% & 28\% & 40\% & 27\% & 33\% \\
ensemble & 3 & mixed & 26\% & 35\% & 20\% & 36\% & 21\% & 31\% \\
shanks & 4 & \LE & 20\% & 32\% & 18\% & 32\% & 18\% & 25\% \\
wynn\_eps & 3 & \LE & 23\% & 31\% & 19\% & 29\% & 18\% & 24\% \\
wynn\_rho & 3 & \LE & 19\% & 24\% & 14\% & 22\% & 13\% & 22\% \\
levin & 6 & \LE & 15\% & 26\% & 17\% & 23\% & 17\% & 21\% \\
brezinski & 2 & \LE & 28\% & 42\% & 19\% & 30\% & 18\% & 23\% \\
anderson & 3 & \LE & 23\% & 31\% & 4\% & 13\% & 4\% & 13\% \\
\midrule
trivial \textit{(4 deployable)} & 4 & \LE & 0\% & 0\% & 0\% & 0\% & 0\% & 0\% \\
\midrule
\textbf{all accelerators} & 52 & & 38\% & 49\% & 28\% & 40\% & 27\% & 35\% \\
\bottomrule
\end{tabular}

\end{adjustbox}
\end{table}
\end{landscape}
\begin{landscape}
\begin{table}[p]
\centering
\caption{\texttt{f04\_classical\_noop.tex}: Classical no-op table by stratum: median improvement factor vs the +/-10 \% band, win rate vs the last value, and strict skill}
\label{tab:f04_classical_noop}
\begin{adjustbox}{max width=\linewidth, max totalheight=0.86\textheight}
% AUTO-GENERATED -- do not hand-edit.  generated by scripts/make_paper_tables.py from the run started 2026-10-01T22:16:39 on code d82fd26 (full; 29,256 s; 7 workers)
% Classical no-op table by stratum: median improvement factor vs the +/-10 % band, win rate vs the last value, and strict skill
% source : results/phase1/phase1_global.csv (core regimes, all three strata, capped cells excluded)
% filter : core regimes, capped excluded; the 21 classical variants = families shanks, wynn_eps, wynn_rho, levin, brezinski, anderson (the Weniger pair was retired: numerically identical to levin_t1/t2); MI = med_improve (median over seeds and cells of current-value error / method error), bold = inside the +/-10 % band; win vs last = per-cell win rate vs last_value, bold = inside [0.4, 0.6]; med. skill = hindsight best-of-four (strict), bold = >= 0.9
\begin{tabular}{llrrrrrrrrr}
\toprule
Family & Method & \multicolumn{3}{c}{$g = 0.5$} & \multicolumn{3}{c}{$g = 0.1$} & \multicolumn{3}{c}{$g = 0.02$} \\
\cmidrule(lr){3-5}\cmidrule(lr){6-8}\cmidrule(lr){9-11}
 & & MI & win vs last & med.\ skill & MI & win vs last & med.\ skill & MI & win vs last & med.\ skill \\
\midrule
Shanks & \meth{shanks\_1} & \textbf{1.001} & 0.604 & \textbf{1.043} & \textbf{1.000} & 0.613 & \textbf{1.538} & \textbf{1.000} & 0.607 & \textbf{1.720} \\
Shanks & \meth{shanks\_2} & \textbf{0.965} & 0.331 & \textbf{1.091} & \textbf{0.980} & 0.353 & \textbf{1.672} & \textbf{0.980} & 0.347 & \textbf{1.928} \\
Shanks & \meth{shanks\_3} & \textbf{0.975} & 0.395 & \textbf{1.093} & \textbf{0.982} & \textbf{0.401} & \textbf{1.516} & \textbf{0.984} & 0.396 & \textbf{1.585} \\
Shanks & \meth{shanks\_4} & \textbf{0.971} & 0.385 & \textbf{1.098} & \textbf{0.973} & 0.370 & \textbf{1.581} & \textbf{0.974} & 0.362 & \textbf{1.686} \\
\addlinespace[2pt]
Wynn $\varepsilon$ & \meth{wynn\_eps\_1} & \textbf{1.001} & 0.604 & \textbf{1.040} & \textbf{1.000} & 0.613 & \textbf{1.523} & \textbf{1.001} & 0.607 & \textbf{1.654} \\
Wynn $\varepsilon$ & \meth{wynn\_eps\_2} & \textbf{0.982} & \textbf{0.488} & \textbf{1.078} & \textbf{0.989} & \textbf{0.471} & \textbf{1.521} & \textbf{0.990} & \textbf{0.466} & \textbf{1.595} \\
Wynn $\varepsilon$ & \meth{wynn\_eps\_3} & \textbf{0.987} & \textbf{0.498} & \textbf{1.080} & \textbf{0.991} & \textbf{0.488} & \textbf{1.518} & \textbf{0.990} & \textbf{0.482} & \textbf{1.606} \\
\addlinespace[2pt]
Wynn $\rho$ & \meth{wynn\_rho\_1} & \textbf{0.995} & \textbf{0.421} & \textbf{1.060} & \textbf{0.997} & \textbf{0.438} & \textbf{1.528} & \textbf{0.997} & \textbf{0.432} & \textbf{1.657} \\
Wynn $\rho$ & \meth{wynn\_rho\_2} & \textbf{1.000} & \textbf{0.465} & \textbf{1.078} & \textbf{1.000} & \textbf{0.454} & \textbf{1.651} & \textbf{1.000} & \textbf{0.464} & \textbf{1.871} \\
Wynn $\rho$ & \meth{wynn\_rho\_3} & \textbf{0.987} & \textbf{0.436} & \textbf{1.092} & \textbf{0.995} & \textbf{0.494} & \textbf{1.546} & \textbf{0.995} & \textbf{0.491} & \textbf{1.653} \\
\addlinespace[2pt]
Levin & \meth{levin\_t1} & \textbf{1.001} & 0.604 & \textbf{1.040} & \textbf{1.000} & 0.613 & \textbf{1.523} & \textbf{1.001} & 0.607 & \textbf{1.654} \\
Levin & \meth{levin\_t2} & \textbf{0.988} & \textbf{0.474} & \textbf{1.051} & \textbf{0.993} & \textbf{0.503} & \textbf{1.520} & \textbf{0.992} & \textbf{0.498} & \textbf{1.602} \\
Levin & \meth{levin\_u1} & \textbf{1.000} & \textbf{0.414} & \textbf{1.079} & \textbf{1.000} & \textbf{0.503} & \textbf{1.603} & \textbf{1.000} & \textbf{0.512} & \textbf{1.924} \\
Levin & \meth{levin\_u2} & \textbf{0.984} & 0.399 & \textbf{1.068} & \textbf{0.992} & \textbf{0.462} & \textbf{1.592} & \textbf{0.991} & \textbf{0.453} & \textbf{1.868} \\
Levin & \meth{levin\_v1} & \textbf{0.965} & 0.277 & \textbf{1.096} & \textbf{0.980} & 0.390 & \textbf{1.571} & \textbf{0.981} & 0.382 & \textbf{1.766} \\
Levin & \meth{levin\_v2} & \textbf{0.973} & \textbf{0.418} & \textbf{1.094} & \textbf{0.981} & \textbf{0.420} & \textbf{1.680} & \textbf{0.981} & \textbf{0.416} & \textbf{1.898} \\
\addlinespace[2pt]
Brezinski $\theta$ & \meth{brezinski\_theta1} & \textbf{1.048} & \textbf{0.546} & \textbf{1.018} & \textbf{1.037} & \textbf{0.592} & \textbf{1.518} & \textbf{1.033} & \textbf{0.584} & \textbf{1.900} \\
Brezinski $\theta$ & \meth{brezinski\_theta2} & \textbf{0.907} & 0.328 & \textbf{1.140} & \textbf{0.941} & 0.362 & \textbf{1.806} & \textbf{0.944} & 0.353 & \textbf{2.067} \\
\addlinespace[2pt]
Anderson & \meth{anderson\_1} & \textbf{1.001} & 0.604 & \textbf{1.043} & \textbf{1.000} & 0.613 & \textbf{1.573} & \textbf{1.000} & 0.607 & \textbf{1.921} \\
Anderson & \meth{anderson\_2} & \textbf{1.000} & \textbf{0.568} & \textbf{1.038} & \textbf{1.000} & \textbf{0.571} & \textbf{1.530} & \textbf{1.000} & \textbf{0.564} & \textbf{1.904} \\
Anderson & \meth{anderson\_3} & \textbf{0.999} & \textbf{0.540} & \textbf{1.044} & \textbf{0.999} & \textbf{0.540} & \textbf{1.559} & \textbf{0.999} & \textbf{0.534} & \textbf{1.895} \\
\midrule
\multicolumn{2}{l}{in the $\pm 10\%$ band (MI in $[0.9, 1.1]$)} & 21/21 &  &  & 21/21 &  &  & 21/21 &  &  \\
\multicolumn{2}{l}{win rate vs last in $[0.4, 0.6]$ (coin flip)} &  & 11/21 &  &  & 13/21 &  &  & 12/21 &  \\
\multicolumn{2}{l}{median skill $\geq 0.9$} &  &  & 21/21 &  &  & 21/21 &  &  & 21/21 \\
\bottomrule
\end{tabular}

\end{adjustbox}
\end{table}
\end{landscape}
\begin{landscape}
\begin{table}[p]
\centering
\caption{\texttt{f04b\_classical\_noop\_holdout.tex}: Classical no-op table by stratum: median improvement factor vs the +/-10 \% band, win rate vs the last value, and strict skill -- the out-of-design families (the layout and statistics of f04 on phase1\_global\_holdout.csv)}
\label{tab:f04b_classical_noop_holdout}
\begin{adjustbox}{max width=\linewidth, max totalheight=0.86\textheight}
% AUTO-GENERATED -- do not hand-edit.  generated by scripts/make_paper_tables.py from the run started 2026-10-01T22:16:39 on code d82fd26 (full; 29,256 s; 7 workers)
% Classical no-op table by stratum: median improvement factor vs the +/-10 % band, win rate vs the last value, and strict skill -- the held-out families (the layout and statistics of f04 on phase1_global_holdout.csv)
% source : results/phase1/phase1_global_holdout.csv (held-out regimes)
% filter : held-out families, capped excluded; the 21 classical variants = families shanks, wynn_eps, wynn_rho, levin, brezinski, anderson (the Weniger pair was retired: numerically identical to levin_t1/t2); MI = med_improve (median over seeds and cells of current-value error / method error), bold = inside the +/-10 % band; win vs last = per-cell win rate vs last_value, bold = inside [0.4, 0.6]; med. skill = hindsight best-of-four (strict), bold = >= 0.9
\begin{tabular}{llrrrrrrrrr}
\toprule
Family & Method & \multicolumn{3}{c}{$g = 0.5$} & \multicolumn{3}{c}{$g = 0.1$} & \multicolumn{3}{c}{$g = 0.02$} \\
\cmidrule(lr){3-5}\cmidrule(lr){6-8}\cmidrule(lr){9-11}
 & & MI & win vs last & med.\ skill & MI & win vs last & med.\ skill & MI & win vs last & med.\ skill \\
\midrule
Shanks & \meth{shanks\_1} & \textbf{1.002} & 0.654 & \textbf{1.027} & \textbf{1.006} & 0.711 & \textbf{1.458} & \textbf{1.006} & 0.711 & \textbf{1.764} \\
Shanks & \meth{shanks\_2} & \textbf{0.957} & 0.387 & \textbf{1.079} & \textbf{0.974} & 0.367 & \textbf{1.861} & \textbf{0.976} & 0.367 & \textbf{1.984} \\
Shanks & \meth{shanks\_3} & \textbf{0.960} & \textbf{0.505} & \textbf{1.062} & \textbf{0.978} & \textbf{0.491} & \textbf{1.554} & \textbf{0.979} & \textbf{0.491} & \textbf{1.847} \\
Shanks & \meth{shanks\_4} & \textbf{0.941} & \textbf{0.457} & \textbf{1.094} & \textbf{0.965} & \textbf{0.440} & \textbf{1.588} & \textbf{0.968} & \textbf{0.440} & \textbf{1.847} \\
\addlinespace[2pt]
Wynn $\varepsilon$ & \meth{wynn\_eps\_1} & \textbf{1.002} & 0.654 & \textbf{1.027} & \textbf{1.006} & 0.711 & \textbf{1.458} & \textbf{1.006} & 0.711 & \textbf{1.764} \\
Wynn $\varepsilon$ & \meth{wynn\_eps\_2} & \textbf{0.976} & \textbf{0.515} & \textbf{1.064} & \textbf{0.985} & \textbf{0.498} & \textbf{1.679} & \textbf{0.986} & \textbf{0.498} & \textbf{2.007} \\
Wynn $\varepsilon$ & \meth{wynn\_eps\_3} & \textbf{0.978} & \textbf{0.561} & \textbf{1.058} & \textbf{0.988} & \textbf{0.562} & \textbf{1.569} & \textbf{0.989} & \textbf{0.562} & \textbf{1.701} \\
\addlinespace[2pt]
Wynn $\rho$ & \meth{wynn\_rho\_1} & \textbf{0.991} & \textbf{0.406} & \textbf{1.039} & \textbf{0.998} & \textbf{0.516} & \textbf{1.719} & \textbf{0.998} & \textbf{0.516} & \textbf{1.845} \\
Wynn $\rho$ & \meth{wynn\_rho\_2} & \textbf{0.983} & \textbf{0.441} & \textbf{1.065} & \textbf{1.004} & \textbf{0.491} & \textbf{1.552} & \textbf{1.004} & \textbf{0.493} & \textbf{1.929} \\
Wynn $\rho$ & \meth{wynn\_rho\_3} & \textbf{0.976} & \textbf{0.470} & \textbf{1.066} & \textbf{0.988} & \textbf{0.467} & \textbf{1.937} & \textbf{0.989} & \textbf{0.469} & \textbf{1.853} \\
\addlinespace[2pt]
Levin & \meth{levin\_t1} & \textbf{1.002} & 0.654 & \textbf{1.027} & \textbf{1.006} & 0.711 & \textbf{1.458} & \textbf{1.006} & 0.711 & \textbf{1.764} \\
Levin & \meth{levin\_t2} & \textbf{0.986} & \textbf{0.550} & \textbf{1.032} & \textbf{0.992} & \textbf{0.547} & \textbf{1.434} & \textbf{0.992} & \textbf{0.547} & \textbf{1.867} \\
Levin & \meth{levin\_u1} & \textbf{1.001} & \textbf{0.480} & \textbf{1.030} & \textbf{1.001} & \textbf{0.504} & \textbf{1.769} & \textbf{1.001} & \textbf{0.571} & \textbf{1.885} \\
Levin & \meth{levin\_u2} & \textbf{0.983} & \textbf{0.493} & \textbf{1.039} & \textbf{0.991} & \textbf{0.547} & \textbf{1.663} & \textbf{0.992} & \textbf{0.547} & \textbf{1.864} \\
Levin & \meth{levin\_v1} & \textbf{0.960} & 0.367 & \textbf{1.071} & \textbf{0.979} & \textbf{0.464} & \textbf{1.696} & \textbf{0.981} & \textbf{0.464} & \textbf{1.951} \\
Levin & \meth{levin\_v2} & \textbf{0.952} & 0.368 & \textbf{1.092} & \textbf{0.977} & \textbf{0.460} & \textbf{1.666} & \textbf{0.979} & \textbf{0.460} & \textbf{1.886} \\
\addlinespace[2pt]
Brezinski $\theta$ & \meth{brezinski\_theta1} & \textbf{1.094} & 0.652 & \textbf{0.987} & 1.102 & 0.709 & \textbf{1.412} & \textbf{1.093} & 0.711 & \textbf{1.566} \\
Brezinski $\theta$ & \meth{brezinski\_theta2} & 0.894 & \textbf{0.428} & \textbf{1.144} & \textbf{0.937} & \textbf{0.422} & \textbf{1.835} & \textbf{0.942} & \textbf{0.422} & \textbf{2.410} \\
\addlinespace[2pt]
Anderson & \meth{anderson\_1} & \textbf{1.002} & 0.709 & \textbf{1.016} & \textbf{1.001} & 0.711 & \textbf{1.785} & \textbf{1.001} & 0.711 & \textbf{1.976} \\
Anderson & \meth{anderson\_2} & \textbf{1.001} & \textbf{0.587} & \textbf{1.022} & \textbf{1.000} & 0.631 & \textbf{1.754} & \textbf{1.000} & 0.631 & \textbf{2.110} \\
Anderson & \meth{anderson\_3} & \textbf{0.996} & \textbf{0.552} & \textbf{1.021} & \textbf{0.998} & \textbf{0.598} & \textbf{1.681} & \textbf{0.998} & \textbf{0.598} & \textbf{2.026} \\
\midrule
\multicolumn{2}{l}{in the $\pm 10\%$ band (MI in $[0.9, 1.1]$)} & 20/21 &  &  & 20/21 &  &  & 21/21 &  &  \\
\multicolumn{2}{l}{win rate vs last in $[0.4, 0.6]$ (coin flip)} &  & 13/21 &  &  & 14/21 &  &  & 14/21 &  \\
\multicolumn{2}{l}{median skill $\geq 0.9$} &  &  & 21/21 &  &  & 21/21 &  &  & 21/21 \\
\bottomrule
\end{tabular}

\end{adjustbox}
\end{table}
\end{landscape}
\begin{landscape}
\begin{table}[p]
\centering
\caption{\texttt{f05\_sweep1\_modes.tex}: Assumed-asymptote mode sweep: what happened when the Phase-2 cascade fired, under each L\_hat mode, core and out-of-design, per stratum}
\label{tab:f05_sweep1_modes}
\begin{adjustbox}{max width=\linewidth, max totalheight=0.86\textheight}
% AUTO-GENERATED -- do not hand-edit.  generated by scripts/make_paper_tables.py at code d82fd26-dirty, results as of commit 2232d3d
% Assumed-asymptote mode sweep: what happened when the Phase-2 cascade fired, under each L_hat mode, core and held-out, per stratum
% source : results/phase5b/phase5b_sweep1_global.csv (Phase 5b sweep 1a, pooled over regimes, capped excluded; src.panels.rule_panel)
% filter : all rows; one window (regime, noise, seed) is one cell of the cascade evaluation; fire = fraction of windows where the cascade routed to rational_fit; lower rec. / lower cell = fraction of fired windows where rational_fit's error was below richardson_1's (both valid; an invalid record never counts as lower; on one-record cells the two coincide); Delta med = median over fired windows of (rational_fit error - richardson_1 error) / richardson_1 error; V_alt = validity rate of rational_fit on the fired windows (errors conditional on validity)
% note   : the cascade features clamp L0 = max(0, min(L_hat, 0.5 * min(window))), so the four non-zero modes coincide whenever L_hat >= 0.5 * min(window); see FACTS.md
\begin{tabular}{lllrrrrrrrrrr}
\toprule
Regime set & $\hat{L}$ mode & $\sigma$ & \multicolumn{5}{c}{$g = 0.5$} & \multicolumn{5}{c}{$g = 0.1$} \\
\cmidrule(lr){4-8}\cmidrule(lr){9-13}
 & & & fire & lower rec. & lower cell & $\Delta$med & V$_{\mathrm{alt}}$ & fire & lower rec. & lower cell & $\Delta$med & V$_{\mathrm{alt}}$ \\
\midrule
core & zero & 0 & 0.233 & 0.417 & 0.417 & $+0.183$ & 1.000 & 0.197 & 0.524 & 0.524 & $-0.089$ & 1.000 \\
 &  & 0.005 & 0.233 & 0.452 & 0.452 & $+0.166$ & 1.000 & 0.197 & 0.571 & 0.571 & $-0.300$ & 1.000 \\
 &  & 0.02 & 0.400 & 0.514 & 0.514 & $-0.013$ & 1.000 & 0.328 & 0.543 & 0.543 & $-0.092$ & 1.000 \\
\addlinespace[1pt]
 & half & 0 & 0.222 & 0.438 & 0.438 & $+0.149$ & 1.000 & 0.188 & 0.550 & 0.550 & $-0.114$ & 1.000 \\
 &  & 0.005 & 0.222 & 0.450 & 0.450 & $+0.166$ & 1.000 & 0.188 & 0.583 & 0.583 & $-0.332$ & 1.000 \\
 &  & 0.02 & 0.400 & 0.514 & 0.514 & $-0.013$ & 1.000 & 0.328 & 0.543 & 0.543 & $-0.092$ & 1.000 \\
\addlinespace[1pt]
 & oracle \textit{(oracle)} & 0 & 0.222 & 0.438 & 0.438 & $+0.150$ & 1.000 & 0.188 & 0.550 & 0.550 & $-0.113$ & 1.000 \\
 &  & 0.005 & 0.222 & 0.450 & 0.450 & $+0.166$ & 1.000 & 0.188 & 0.583 & 0.583 & $-0.332$ & 1.000 \\
 &  & 0.02 & 0.400 & 0.514 & 0.514 & $-0.012$ & 1.000 & 0.328 & 0.543 & 0.543 & $-0.092$ & 1.000 \\
\addlinespace[1pt]
 & double & 0 & 0.222 & 0.438 & 0.438 & $+0.150$ & 1.000 & 0.188 & 0.550 & 0.550 & $-0.113$ & 1.000 \\
 &  & 0.005 & 0.222 & 0.463 & 0.463 & $+0.145$ & 1.000 & 0.188 & 0.583 & 0.583 & $-0.332$ & 1.000 \\
 &  & 0.02 & 0.400 & 0.514 & 0.514 & $-0.011$ & 1.000 & 0.328 & 0.543 & 0.543 & $-0.092$ & 1.000 \\
\addlinespace[1pt]
 & winmin & 0 & 0.222 & 0.438 & 0.438 & $+0.150$ & 1.000 & 0.188 & 0.550 & 0.550 & $-0.114$ & 1.000 \\
 &  & 0.005 & 0.222 & 0.450 & 0.450 & $+0.166$ & 1.000 & 0.188 & 0.583 & 0.583 & $-0.332$ & 1.000 \\
 &  & 0.02 & 0.400 & 0.514 & 0.514 & $-0.012$ & 1.000 & 0.328 & 0.543 & 0.543 & $-0.092$ & 1.000 \\
\midrule
held-out & zero & 0 & 0.017 & 0.000 & 0.000 & $+4.135$ & 1.000 & 0.017 & 0.000 & 0.000 & $+2.618$ & 1.000 \\
 &  & 0.005 & 0.025 & 0.667 & 0.667 & $-0.024$ & 1.000 & 0.025 & 0.667 & 0.667 & $-0.171$ & 1.000 \\
 &  & 0.02 & 0.267 & 0.375 & 0.375 & $+0.166$ & 1.000 & 0.269 & 0.438 & 0.438 & $+0.085$ & 1.000 \\
\addlinespace[1pt]
 & half & 0 & 0.008 & 0.000 & 0.000 & $+4.135$ & 1.000 & 0.008 & 0.000 & 0.000 & $+2.618$ & 1.000 \\
 &  & 0.005 & 0.017 & 0.500 & 0.500 & $-0.081$ & 1.000 & 0.017 & 0.500 & 0.500 & $-0.046$ & 1.000 \\
 &  & 0.02 & 0.267 & 0.375 & 0.375 & $+0.166$ & 1.000 & 0.269 & 0.438 & 0.438 & $+0.085$ & 1.000 \\
\addlinespace[1pt]
 & oracle \textit{(oracle)} & 0 & 0.008 & 0.000 & 0.000 & $+4.135$ & 1.000 & 0.008 & 0.000 & 0.000 & $+2.618$ & 1.000 \\
 &  & 0.005 & 0.017 & 0.500 & 0.500 & $-0.081$ & 1.000 & 0.017 & 0.500 & 0.500 & $-0.046$ & 1.000 \\
 &  & 0.02 & 0.267 & 0.375 & 0.375 & $+0.166$ & 1.000 & 0.269 & 0.438 & 0.438 & $+0.085$ & 1.000 \\
\addlinespace[1pt]
 & double & 0 & 0.008 & 0.000 & 0.000 & $+4.135$ & 1.000 & 0.008 & 0.000 & 0.000 & $+2.618$ & 1.000 \\
 &  & 0.005 & 0.017 & 0.500 & 0.500 & $-0.081$ & 1.000 & 0.017 & 0.500 & 0.500 & $-0.046$ & 1.000 \\
 &  & 0.02 & 0.267 & 0.375 & 0.375 & $+0.166$ & 1.000 & 0.269 & 0.438 & 0.438 & $+0.085$ & 1.000 \\
\addlinespace[1pt]
 & winmin & 0 & 0.008 & 0.000 & 0.000 & $+4.135$ & 1.000 & 0.008 & 0.000 & 0.000 & $+2.618$ & 1.000 \\
 &  & 0.005 & 0.017 & 0.500 & 0.500 & $-0.081$ & 1.000 & 0.017 & 0.500 & 0.500 & $-0.046$ & 1.000 \\
 &  & 0.02 & 0.267 & 0.375 & 0.375 & $+0.166$ & 1.000 & 0.269 & 0.438 & 0.438 & $+0.085$ & 1.000 \\
\bottomrule
\end{tabular}

\end{adjustbox}
\end{table}
\end{landscape}
\begin{landscape}
\begin{table}[p]
\centering
\caption{\texttt{f05b\_sweep1\_consumers\_core.tex}: Assumed-asymptote sweep 1b, core curve families: the L\_hat-consuming accelerators under every mode, and the value of knowing the floor}
\label{tab:f05b_sweep1_consumers_core}
\begin{adjustbox}{max width=\linewidth, max totalheight=0.86\textheight}
% AUTO-GENERATED -- do not hand-edit.  generated by scripts/make_paper_tables.py from the run started 2026-10-01T22:16:39 on code d82fd26 (full; 29,256 s; 7 workers)
% Assumed-asymptote sweep 1b, core curve families: the L_hat-consuming accelerators under every mode, and the value of knowing the floor
% source : results/phase5b/phase5b_sweep1_consumers.csv (Phase 5b sweep 1b; capped excluded; pooled over noise)
% filter : regime_set == core; med. err = median error of the method under that L_hat mode; win vs last = per-record win rate against last_value, whose prediction does not depend on the mode (the mode-invariant reference); the vs-L_hat columns stay in the CSV but are not the headline comparison because their reference changes with the mode
\begin{tabular}{llrrrrrrrrrr}
\toprule
Method & row & \multicolumn{5}{c}{$g = 0.5$} & \multicolumn{5}{c}{$g = 0.1$} \\
\cmidrule(lr){3-7}\cmidrule(lr){8-12}
 & & zero & half & oracle & double & winmin & zero & half & oracle & double & winmin \\
\midrule
\meth{richardson\_1} & med.\ err & 0.0173 & 0.0173 & 0.0173 & 0.0173 & 0.0173 & 0.0309 & 0.0309 & 0.0309 & 0.0309 & 0.0309 \\
 & win vs last & 0.774 & 0.774 & 0.774 & 0.774 & 0.774 & 0.799 & 0.799 & 0.799 & 0.800 & 0.799 \\
\addlinespace[1pt]
\meth{richardson\_2} & med.\ err & 0.0158 & 0.0156 & 0.0160 & 0.0159 & 0.0162 & 0.0263 & 0.0263 & 0.0263 & 0.0263 & 0.0265 \\
 & win vs last & 0.778 & 0.778 & 0.778 & 0.779 & 0.776 & 0.778 & 0.778 & 0.778 & 0.779 & 0.775 \\
\addlinespace[1pt]
\meth{richardson\_3} & med.\ err & 0.0164 & 0.0161 & 0.0161 & 0.0162 & 0.0161 & 0.0224 & 0.0230 & 0.0225 & 0.0225 & 0.0229 \\
 & win vs last & 0.766 & 0.769 & 0.771 & 0.772 & 0.778 & 0.754 & 0.758 & 0.758 & 0.765 & 0.766 \\
\addlinespace[1pt]
\meth{single\_exp\_fit} & med.\ err & 0.0214 & 0.0214 & 0.0214 & 0.0216 & 0.0214 & 0.0378 & 0.0378 & 0.0378 & 0.0384 & 0.0378 \\
 & win vs last & 0.863 & 0.863 & 0.863 & 0.863 & 0.863 & 0.920 & 0.919 & 0.920 & 0.912 & 0.920 \\
\addlinespace[1pt]
\meth{double\_exp\_fit} & med.\ err & 0.0176 & 0.0173 & 0.0161 & 0.0176 & 0.0199 & 0.0265 & 0.0263 & 0.0246 & 0.0263 & 0.0319 \\
 & win vs last & 0.850 & 0.850 & 0.846 & 0.852 & 0.859 & 0.880 & 0.888 & 0.871 & 0.890 & 0.884 \\
\addlinespace[1pt]
\meth{rational\_fit} & med.\ err & 0.0167 & 0.0167 & 0.0167 & 0.0167 & 0.0167 & 0.0290 & 0.0290 & 0.0290 & 0.0290 & 0.0290 \\
 & win vs last & 0.835 & 0.835 & 0.835 & 0.835 & 0.835 & 0.851 & 0.851 & 0.851 & 0.851 & 0.851 \\
\addlinespace[1pt]
\meth{log\_fit} & med.\ err & 0.0385 & 0.0385 & 0.0385 & 0.0385 & 0.0385 & 0.0684 & 0.0684 & 0.0684 & 0.0684 & 0.0684 \\
 & win vs last & 0.618 & 0.618 & 0.618 & 0.618 & 0.618 & 0.675 & 0.675 & 0.675 & 0.675 & 0.675 \\
\addlinespace[1pt]
\meth{log\_linear} & med.\ err & 0.0145 & 0.0138 & 0.0163 & 0.0253 & 0.0420 & 0.0240 & 0.0215 & 0.0162 & 0.0500 & 0.0752 \\
 & win vs last & 0.800 & 0.807 & 0.772 & 0.498 & 0.756 & 0.747 & 0.763 & 0.760 & 0.479 & 0.827 \\
\addlinespace[1pt]
\meth{stability\_weighted} & med.\ err & 0.0454 & 0.0454 & 0.0454 & 0.0454 & 0.0454 & 0.0454 & 0.0454 & 0.0454 & 0.0454 & 0.0454 \\
 & win vs last & 0.608 & 0.608 & 0.608 & 0.608 & 0.608 & 0.658 & 0.658 & 0.658 & 0.658 & 0.658 \\
\addlinespace[1pt]
\meth{median\_ensemble} & med.\ err & 0.0387 & 0.0387 & 0.0387 & 0.0387 & 0.0387 & 0.0457 & 0.0457 & 0.0457 & 0.0457 & 0.0457 \\
 & win vs last & 0.678 & 0.678 & 0.678 & 0.678 & 0.678 & 0.678 & 0.678 & 0.678 & 0.678 & 0.678 \\
\addlinespace[1pt]
predict $\hat L$ (\meth{constant\_assumed}) -- the value of knowing the floor & med.\ err & 0.1206 & 0.0972 & 0.0579 & 0.0616 & 0.0323 & 0.0593 & 0.0322 & 0.0077 & 0.0348 & 0.0507 \\
 & win vs last & 0.084 & 0.133 & 0.714 & 0.582 & 0.785 & 0.544 & 0.682 & 0.995 & 0.580 & 0.844 \\
\addlinespace[1pt]
\bottomrule
\end{tabular}

\end{adjustbox}
\end{table}
\end{landscape}
\begin{landscape}
\begin{table}[p]
\centering
\caption{\texttt{f05b\_sweep1\_consumers\_holdout.tex}: Assumed-asymptote sweep 1b, out-of-design curve families: the L\_hat-consuming accelerators under every mode, and the value of knowing the floor}
\label{tab:f05b_sweep1_consumers_holdout}
\begin{adjustbox}{max width=\linewidth, max totalheight=0.86\textheight}
% AUTO-GENERATED -- do not hand-edit.  generated by scripts/make_paper_tables.py from the run started 2026-10-01T22:16:39 on code d82fd26 (full; 29,256 s; 7 workers)
% Assumed-asymptote sweep 1b, holdout regimes: the L_hat-consuming accelerators under every mode, and the value of knowing the floor
% source : results/phase5b/phase5b_sweep1_consumers.csv (Phase 5b sweep 1b; capped excluded; pooled over noise)
% filter : regime_set == holdout; med. err = median error of the method under that L_hat mode; win vs last = per-record win rate against last_value, whose prediction does not depend on the mode (the mode-invariant reference); the vs-L_hat columns stay in the CSV but are not the headline comparison because their reference changes with the mode
\begin{tabular}{llrrrrrrrrrr}
\toprule
Method & row & \multicolumn{5}{c}{$g = 0.5$} & \multicolumn{5}{c}{$g = 0.1$} \\
\cmidrule(lr){3-7}\cmidrule(lr){8-12}
 & & zero & half & oracle & double & winmin & zero & half & oracle & double & winmin \\
\midrule
\meth{richardson\_1} & med.\ err & 0.0102 & 0.0102 & 0.0102 & 0.0102 & 0.0102 & 0.0212 & 0.0212 & 0.0212 & 0.0212 & 0.0212 \\
 & win vs last & 0.786 & 0.786 & 0.786 & 0.786 & 0.786 & 0.801 & 0.801 & 0.801 & 0.801 & 0.801 \\
\addlinespace[1pt]
\meth{richardson\_2} & med.\ err & 0.0105 & 0.0104 & 0.0107 & 0.0095 & 0.0104 & 0.0216 & 0.0215 & 0.0221 & 0.0217 & 0.0218 \\
 & win vs last & 0.775 & 0.778 & 0.786 & 0.786 & 0.769 & 0.795 & 0.798 & 0.809 & 0.807 & 0.790 \\
\addlinespace[1pt]
\meth{richardson\_3} & med.\ err & 0.0101 & 0.0100 & 0.0100 & 0.0096 & 0.0097 & 0.0195 & 0.0221 & 0.0207 & 0.0191 & 0.0198 \\
 & win vs last & 0.750 & 0.767 & 0.758 & 0.758 & 0.764 & 0.784 & 0.795 & 0.790 & 0.790 & 0.795 \\
\addlinespace[1pt]
\meth{single\_exp\_fit} & med.\ err & 0.0259 & 0.0259 & 0.0259 & 0.0259 & 0.0259 & 0.0350 & 0.0350 & 0.0350 & 0.0350 & 0.0350 \\
 & win vs last & 0.881 & 0.881 & 0.881 & 0.878 & 0.881 & 0.910 & 0.910 & 0.910 & 0.908 & 0.910 \\
\addlinespace[1pt]
\meth{double\_exp\_fit} & med.\ err & 0.0182 & 0.0178 & 0.0158 & 0.0176 & 0.0230 & 0.0333 & 0.0330 & 0.0319 & 0.0333 & 0.0342 \\
 & win vs last & 0.875 & 0.867 & 0.878 & 0.872 & 0.875 & 0.896 & 0.894 & 0.885 & 0.905 & 0.899 \\
\addlinespace[1pt]
\meth{rational\_fit} & med.\ err & 0.0086 & 0.0086 & 0.0086 & 0.0086 & 0.0086 & 0.0154 & 0.0154 & 0.0154 & 0.0154 & 0.0154 \\
 & win vs last & 0.792 & 0.792 & 0.792 & 0.792 & 0.792 & 0.824 & 0.824 & 0.824 & 0.824 & 0.824 \\
\addlinespace[1pt]
\meth{log\_fit} & med.\ err & 0.0383 & 0.0383 & 0.0383 & 0.0383 & 0.0383 & 0.0516 & 0.0516 & 0.0516 & 0.0516 & 0.0516 \\
 & win vs last & 0.603 & 0.603 & 0.603 & 0.603 & 0.603 & 0.706 & 0.706 & 0.706 & 0.706 & 0.706 \\
\addlinespace[1pt]
\meth{log\_linear} & med.\ err & 0.0144 & 0.0151 & 0.0215 & 0.0277 & 0.0455 & 0.0314 & 0.0303 & 0.0370 & 0.0578 & 0.0828 \\
 & win vs last & 0.750 & 0.769 & 0.758 & 0.394 & 0.775 & 0.809 & 0.860 & 0.927 & 0.507 & 0.838 \\
\addlinespace[1pt]
\meth{stability\_weighted} & med.\ err & 0.0392 & 0.0392 & 0.0392 & 0.0392 & 0.0392 & 0.0308 & 0.0308 & 0.0308 & 0.0308 & 0.0308 \\
 & win vs last & 0.672 & 0.672 & 0.672 & 0.672 & 0.672 & 0.767 & 0.767 & 0.767 & 0.767 & 0.767 \\
\addlinespace[1pt]
\meth{median\_ensemble} & med.\ err & 0.0254 & 0.0254 & 0.0254 & 0.0254 & 0.0254 & 0.0391 & 0.0391 & 0.0391 & 0.0391 & 0.0391 \\
 & win vs last & 0.722 & 0.722 & 0.722 & 0.722 & 0.722 & 0.781 & 0.781 & 0.781 & 0.781 & 0.781 \\
\addlinespace[1pt]
predict $\hat L$ (\meth{constant\_assumed}) -- the value of knowing the floor & med.\ err & 0.0985 & 0.0844 & 0.0455 & 0.0521 & 0.0233 & 0.0615 & 0.0348 & 0.0092 & 0.0432 & 0.0599 \\
 & win vs last & 0.039 & 0.094 & 0.736 & 0.592 & 0.878 & 0.527 & 0.698 & 0.994 & 0.602 & 0.947 \\
\addlinespace[1pt]
\bottomrule
\end{tabular}

\end{adjustbox}
\end{table}
\end{landscape}
\begin{landscape}
\begin{table}[p]
\centering
\caption{\texttt{f06\_generalisation.tex}: Out-of-design vs core generalisation: per-method rank shift per stratum, with median errors at the headline stratum; rational\_fit and log\_linear in bold}
\label{tab:f06_generalisation}
\begin{adjustbox}{max width=\linewidth, max totalheight=0.86\textheight}
% AUTO-GENERATED -- do not hand-edit.  generated by scripts/make_paper_tables.py from the run started 2026-10-01T22:16:39 on code d82fd26 (full; 29,256 s; 7 workers)
% Out-of-design vs core generalisation: per-method rank shift per stratum, with median errors at the headline stratum; rational_fit and log_linear in bold
% source : results/phase1/phase1_global.csv (core regimes, all three strata, capped cells excluded)
% source : results/phase1/phase1_global_holdout.csv (out-of-design regimes)
% filter : ranks among the 52 accelerators (rank-eligible only; -- = below the validity floor in that regime set); Delta = r_h - r_c (positive = worse on the out-of-design regimes); rows sorted by core rank at the headline stratum
\begin{tabular}{lrrrrrrrrrrr}
\toprule
Method & \multicolumn{3}{c}{$g = 0.5$} & \multicolumn{5}{c}{$g = 0.1$} & \multicolumn{3}{c}{$g = 0.02$} \\
\cmidrule(lr){2-4}\cmidrule(lr){5-9}\cmidrule(lr){10-12}
 & $r_c$ & $r_h$ & $\Delta$ & $r_c$ & $r_h$ & $\Delta$ & err$_c$ & err$_h$ & $r_c$ & $r_h$ & $\Delta$ \\
\midrule
\textbf{\meth{log\_linear}} & 1 & 8 & $+7$ & 1 & 10 & $+9$ & 0.0202 & 0.0305 & 4 & 12 & $+8$ \\
\meth{richardson\_3} & 2 & 5 & $+3$ & 2 & 8 & $+6$ & 0.0222 & 0.0260 & 2 & 3 & $+1$ \\
\meth{richardson\_2} & 7 & 7 & $+0$ & 3 & 11 & $+8$ & 0.0257 & 0.0316 & 1 & 6 & $+5$ \\
\meth{double\_exp\_fit} & 9 & 14 & $+5$ & 4 & 13 & $+9$ & 0.0287 & 0.0322 & 13 & 13 & $+0$ \\
\meth{pade\_34} & 16 & 13 & $-3$ & 5 & 3 & $-2$ & 0.0291 & 0.0211 & 7 & 7 & $+0$ \\
\textbf{\meth{rational\_fit}} & 8 & 1 & $-7$ & 6 & 1 & $-5$ & 0.0296 & 0.0074 & 3 & 1 & $-2$ \\
\meth{richardson\_1} & 14 & 6 & $-8$ & 7 & 12 & $+5$ & 0.0301 & 0.0318 & 5 & 5 & $+0$ \\
\meth{pade\_23} & 12 & 10 & $-2$ & 8 & 6 & $-2$ & 0.0302 & 0.0243 & 8 & 10 & $+2$ \\
\meth{pade\_12} & 13 & 11 & $-2$ & 9 & 5 & $-4$ & 0.0314 & 0.0225 & 6 & 8 & $+2$ \\
\meth{pade\_45} & 15 & 9 & $-6$ & 10 & 7 & $-3$ & 0.0318 & 0.0253 & 9 & 9 & $+0$ \\
\meth{pade\_33} & 5 & 3 & $-2$ & 11 & 2 & $-9$ & 0.0329 & 0.0153 & 11 & 2 & $-9$ \\
\meth{pade\_22} & 4 & 2 & $-2$ & 12 & 4 & $-8$ & 0.0333 & 0.0217 & 12 & 4 & $-8$ \\
\meth{single\_exp\_fit} & 11 & 16 & $+5$ & 13 & 15 & $+2$ & 0.0361 & 0.0342 & 14 & 14 & $+0$ \\
\meth{richardson\_a05} & 10 & 17 & $+7$ & 14 & 9 & $-5$ & 0.0384 & 0.0267 & 15 & 11 & $-4$ \\
\meth{pade\_13} & 17 & 15 & $-2$ & 15 & 33 & $+18$ & 0.0390 & 0.0563 & 10 & 36 & $+26$ \\
\meth{pade\_11} & 3 & 12 & $+9$ & 16 & 16 & $+0$ & 0.0397 & 0.0356 & 22 & 15 & $-7$ \\
\meth{pade\_44} & 6 & 4 & $-2$ & 17 & 14 & $-3$ & 0.0410 & 0.0338 & 19 & -- & -- \\
\meth{wynn\_eps\_3} & 29 & 24 & $-5$ & 18 & 29 & $+11$ & 0.0438 & 0.0480 & 17 & 24 & $+7$ \\
\meth{stability\_weighted} & 25 & 34 & $+9$ & 19 & 17 & $-2$ & 0.0456 & 0.0423 & 18 & 17 & $-1$ \\
\meth{median\_ensemble} & 19 & 23 & $+4$ & 20 & 18 & $-2$ & 0.0473 & 0.0436 & 21 & 18 & $-3$ \\
\meth{wynn\_eps\_2} & 23 & 36 & $+13$ & 21 & 23 & $+2$ & 0.0473 & 0.0476 & 24 & 22 & $-2$ \\
\meth{levin\_t2} & 28 & 31 & $+3$ & 22 & 21 & $-1$ & 0.0475 & 0.0471 & 25 & 20 & $-5$ \\
\meth{shanks\_4} & 32 & 28 & $-4$ & 23 & 31 & $+8$ & 0.0475 & 0.0490 & 23 & 26 & $+3$ \\
\meth{levin\_u2} & 37 & 30 & $-7$ & 24 & 36 & $+12$ & 0.0475 & 0.0742 & 26 & 33 & $+7$ \\
\meth{levin\_v1} & 41 & 27 & $-14$ & 25 & 37 & $+12$ & 0.0476 & 0.0759 & 30 & 35 & $+5$ \\
\meth{shanks\_3} & 31 & 29 & $-2$ & 26 & 22 & $-4$ & 0.0476 & 0.0473 & 29 & 21 & $-8$ \\
\meth{brezinski\_theta2} & 33 & 43 & $+10$ & 27 & 32 & $+5$ & 0.0478 & 0.0526 & 20 & 27 & $+7$ \\
\meth{shanks\_2} & 30 & 38 & $+8$ & 28 & 28 & $+0$ & 0.0479 & 0.0478 & 28 & 23 & $-5$ \\
\meth{wynn\_rho\_1} & 34 & 37 & $+3$ & 29 & 20 & $-9$ & 0.0489 & 0.0469 & 33 & 19 & $-14$ \\
\meth{wynn\_rho\_3} & 38 & 35 & $-3$ & 30 & 43 & $+13$ & 0.0502 & 0.1044 & 37 & 42 & $+5$ \\
\meth{levin\_v2} & 36 & 42 & $+6$ & 31 & 30 & $-1$ & 0.0503 & 0.0483 & 36 & 25 & $-11$ \\
\meth{levin\_t1} & 21 & 20 & $-1$ & 32 & 26 & $-6$ & 0.0505 & 0.0476 & 31 & 31 & $+0$ \\
\meth{wynn\_eps\_1} & 22 & 22 & $+0$ & 33 & 27 & $-6$ & 0.0505 & 0.0476 & 32 & 32 & $+0$ \\
\meth{wynn\_rho\_2} & 35 & 26 & $-9$ & 34 & 35 & $+1$ & 0.0510 & 0.0628 & 38 & 34 & $-4$ \\
\meth{richardson\_a10} & 18 & 25 & $+7$ & 35 & 34 & $-1$ & 0.0561 & 0.0614 & 16 & 16 & $+0$ \\
\meth{brezinski\_theta1} & 20 & 18 & $-2$ & 36 & 19 & $-17$ & 0.0584 & 0.0465 & 27 & 28 & $+1$ \\
\meth{best\_shanks\_wynn} & 26 & 19 & $-7$ & 37 & 24 & $-13$ & 0.0589 & 0.0476 & 34 & 29 & $-5$ \\
\meth{shanks\_1} & 27 & 21 & $-6$ & 38 & 25 & $-13$ & 0.0589 & 0.0476 & 35 & 30 & $-5$ \\
\meth{log\_fit} & 24 & 45 & $+21$ & 39 & 42 & $+3$ & 0.0759 & 0.1008 & 39 & 41 & $+2$ \\
\meth{levin\_u1} & 39 & 39 & $+0$ & 40 & 41 & $+1$ & 0.0863 & 0.0889 & 40 & 37 & $-3$ \\
\meth{anderson\_2} & 42 & 41 & $-1$ & 41 & 40 & $-1$ & 0.1036 & 0.0877 & 41 & 40 & $-1$ \\
\meth{anderson\_3} & 40 & 40 & $+0$ & 42 & 39 & $-3$ & 0.1038 & 0.0870 & 42 & 39 & $-3$ \\
\meth{anderson\_1} & 43 & 32 & $-11$ & 43 & 38 & $-5$ & 0.1039 & 0.0853 & 43 & 38 & $-5$ \\
\meth{richardson\_a20} & 44 & 44 & $+0$ & 44 & 44 & $+0$ & 0.1113 & 0.1378 & 44 & 43 & $-1$ \\
$\dagger$\meth{pade\_31} & -- & 48 & -- & -- & 45 & -- & 1.9535 & 2.8890 & -- & -- & -- \\
$\dagger$\meth{geom\_avg\_diff} & -- & -- & -- & -- & -- & -- & 0.0505 & 0.0503 & -- & -- & -- \\
$\dagger$\meth{linear} & 46 & 47 & $+1$ & -- & -- & -- & 0.3057 & 0.5665 & -- & -- & -- \\
$\dagger$\meth{neville\_2} & -- & -- & -- & -- & -- & -- & 4.5269 & 4.9545 & -- & -- & -- \\
$\dagger$\meth{neville\_3} & -- & -- & -- & -- & -- & -- & 111.0798 & 137.1949 & -- & -- & -- \\
$\dagger$\meth{neville\_4} & -- & -- & -- & -- & -- & -- & 96.0674 & 0.1069 & -- & -- & -- \\
$\dagger$\meth{pade\_21} & 47 & 46 & $-1$ & -- & -- & -- & 0.3364 & 0.1909 & -- & -- & -- \\
$\dagger$\meth{pade\_32} & 45 & 33 & $-12$ & -- & -- & -- & 0.2574 & 0.0680 & -- & -- & -- \\
\midrule
\multicolumn{12}{l}{Spearman $\rho$ of core vs out-of-design rank over methods ranked in both: $g = 0.5$: 0.877 ($n = 47$), $g = 0.1$: 0.833 ($n = 44$), $g = 0.02$: 0.873 ($n = 43$)} \\
\bottomrule
\end{tabular}

\end{adjustbox}
\end{table}
\end{landscape}
\begin{landscape}
\begin{table}[p]
\centering
\caption{\texttt{f07\_real\_data\_v2.tex}: Real-data re-evaluation over the (depth x target) grid: routed method, cascade error and skill per cell}
\label{tab:f07_real_data_v2}
\begin{adjustbox}{max width=\linewidth, max totalheight=0.86\textheight}
% AUTO-GENERATED -- do not hand-edit.  generated by scripts/make_paper_tables.py at code d82fd26-dirty, results as of commit 2232d3d
% Real-data re-evaluation over the (depth x target) grid: routed method, cascade error and skill per cell
% source : results/real_data/real_data_summary_v2.csv (recorded XGBoost validation curves re-evaluated under the v2 design; L_hat mode zero; cascade = Phase-2 rules on the window features)
% filter : all 90 cells; routed: R1 = richardson_1, RF = rational_fit; err = |cascade prediction - recorded value at the target round|; skill = hindsight best-of-four (strict): err / best-of-four trivial error on the same cell, bold when >= 1; a trailing + marks a post-minimum target (target round beyond the recorded curve's argmin)
\begin{tabular}{lrlrrlrrlrr}
\toprule
Dataset & $\nobs$ & \multicolumn{3}{c}{target round 300} & \multicolumn{3}{c}{target round 400} & \multicolumn{3}{c}{target round 500} \\
\cmidrule(lr){3-5}\cmidrule(lr){6-8}\cmidrule(lr){9-11}
 & & routed & err & skill & routed & err & skill & routed & err & skill \\
\midrule
\texttt{adult} & 30 & R1 & 0.0225 & 0.36 & R1 & 0.0330 & 0.53$^{+}$ & R1 & 0.0415 & 0.68$^{+}$ \\
 & 60 & R1 & 0.0398 & \textbf{1.53} & R1 & 0.0507 & \textbf{1.97}$^{+}$ & R1 & 0.0595 & \textbf{2.40}$^{+}$ \\
 & 90 & R1 & 0.0019 & 0.13 & R1 & 0.0039 & 0.28$^{+}$ & R1 & 0.0057 & 0.44$^{+}$ \\
 & 120 & RF & 0.0055 & 0.72 & RF & 0.0082 & \textbf{1.12}$^{+}$ & RF & 0.0106 & \textbf{1.63}$^{+}$ \\
 & 150 & RF & 0.0033 & 0.75 & RF & 0.0055 & \textbf{1.35}$^{+}$ & RF & 0.0074 & \textbf{2.33}$^{+}$ \\
\addlinespace[2pt]
\texttt{bank\_marketing} & 30 & R1 & 0.0113 & 0.27$^{+}$ & R1 & 0.0183 & 0.45$^{+}$ & R1 & 0.0246 & 0.63$^{+}$ \\
 & 60 & R1 & 0.0209 & \textbf{1.13}$^{+}$ & R1 & 0.0283 & \textbf{1.59}$^{+}$ & R1 & 0.0348 & \textbf{2.16}$^{+}$ \\
 & 90 & R1 & 0.0020 & 0.20$^{+}$ & R1 & 0.0042 & 0.46$^{+}$ & R1 & 0.0067 & 0.88$^{+}$ \\
 & 120 & RF & 0.0081 & \textbf{1.73}$^{+}$ & RF & 0.0116 & \textbf{2.99}$^{+}$ & RF & 0.0149 & \textbf{6.51}$^{+}$ \\
 & 150 & RF & 0.0040 & \textbf{1.82}$^{+}$ & RF & 0.0063 & \textbf{4.56}$^{+}$ & RF & 0.0088 & \textbf{51.69}$^{+}$ \\
\addlinespace[2pt]
\texttt{covertype} & 30 & R1 & 0.0911 & 0.34 & R1 & 0.1023 & 0.34 & R1 & 0.1114 & 0.34 \\
 & 60 & R1 & 0.0306 & 0.20 & R1 & 0.0407 & 0.22 & R1 & 0.0492 & 0.23 \\
 & 90 & R1 & 0.0064 & 0.06 & R1 & 0.0253 & 0.18 & R1 & 0.0412 & 0.25 \\
 & 120 & R1 & 0.0174 & 0.23 & R1 & 0.0388 & 0.35 & R1 & 0.0567 & 0.42 \\
 & 150 & R1 & 0.0133 & 0.23 & R1 & 0.0328 & 0.36 & R1 & 0.0491 & 0.42 \\
\addlinespace[2pt]
\texttt{higgs} & 30 & RF & 0.0096 & 0.17 & RF & 0.0106 & 0.19 & RF & 0.0116 & 0.20 \\
 & 60 & RF & 0.0091 & 0.32 & RF & 0.0101 & 0.33 & RF & 0.0111 & 0.36 \\
 & 90 & RF & 0.0036 & 0.22 & RF & 0.0040 & 0.22 & RF & 0.0046 & 0.24 \\
 & 120 & RF & 0.0020 & 0.20 & RF & 0.0021 & 0.17 & RF & 0.0025 & 0.19 \\
 & 150 & RF & 0.0008 & 0.11 & RF & 0.0006 & 0.07 & RF & 0.0008 & 0.09 \\
\addlinespace[2pt]
\texttt{jannis} & 30 & RF & 0.0490 & 0.37 & RF & 0.0543 & 0.40 & RF & 0.0583 & 0.43 \\
 & 60 & R1 & 0.0430 & 0.67 & R1 & 0.0592 & 0.88 & R1 & 0.0727 & \textbf{1.06} \\
 & 90 & R1 & 0.0100 & 0.27 & R1 & 0.0134 & 0.33 & R1 & 0.0165 & 0.40 \\
 & 120 & RF & 0.0067 & 0.30 & RF & 0.0089 & 0.35 & RF & 0.0110 & 0.41 \\
 & 150 & RF & 0.0045 & 0.32 & RF & 0.0060 & 0.35 & RF & 0.0077 & 0.42 \\
\addlinespace[2pt]
\texttt{miniboone} & 30 & R1 & 0.0232 & 0.21 & R1 & 0.0149 & 0.13 & R1 & 0.0083 & 0.07 \\
 & 60 & R1 & 0.0068 & 0.14 & R1 & 0.0147 & 0.28 & R1 & 0.0210 & 0.39 \\
 & 90 & R1 & 0.0071 & 0.26 & R1 & 0.0085 & 0.28 & R1 & 0.0096 & 0.31 \\
 & 120 & R1 & 0.0020 & 0.11 & R1 & 0.0023 & 0.12 & R1 & 0.0027 & 0.13 \\
 & 150 & R1 & 0.0001 & 0.01 & R1 & 0.0001 & 0.01 & R1 & 0.0001 & 0.01 \\
\midrule
\multicolumn{2}{l}{median skill / cells with skill $\geq 1$ (of 30)} &  & & 0.26 / 4 &  & & 0.34 / 6 &  & & 0.41 / 7 \\
\multicolumn{2}{l}{pre-minimum targets: cells / skill $\geq 1$ / median skill} &  & & 25 / 1 / 0.23 &  & & 20 / 0 / 0.28 &  & & 20 / 1 / 0.32 \\
\multicolumn{2}{l}{post-minimum targets: cells / skill $\geq 1$ / median skill} &  & & 5 / 3 / 1.13 &  & & 10 / 6 / 1.23 &  & & 10 / 6 / 1.89 \\
\bottomrule
\end{tabular}

\end{adjustbox}
\end{table}
\end{landscape}
\begin{landscape}
\begin{table}[p]
\centering
\caption{\texttt{f07b\_real\_data\_legacy18.tex}: Legacy 18-cell real-data table (pre-redesign design), preserved unchanged}
\label{tab:f07b_real_data_legacy18}
\begin{adjustbox}{max width=\linewidth, max totalheight=0.86\textheight}
% AUTO-GENERATED -- do not hand-edit.  generated by scripts/make_paper_tables.py from the run started 2026-10-01T22:16:39 on code d82fd26 (full; 29,256 s; 7 workers)
% Legacy 18-cell real-data table (pre-redesign design), preserved unchanged
% source : results/real_data/real_data_results.csv (the pre-redesign 18-cell run: fixed target = final recorded round, legacy assumed asymptote 0.01; kept for the record, regenerable with scripts/analyze_real_diagnostics_legacy.py)
% filter : all 18 rows; bold = cascade worse than the last observed value
\begin{tabular}{lrlrrrr}
\toprule
Dataset & $\nobs$ & routed & predicted & cascade err & last-value err & improv. \\
\midrule
\texttt{adult} & 30 & \meth{richardson\_1} & 0.2314 & 0.0415 & 0.0611 & $+32.0$\% \\
\texttt{adult} & 60 & \meth{richardson\_1} & 0.2134 & \textbf{0.0595} & 0.0248 & $-139.6$\% \\
\texttt{adult} & 90 & \meth{richardson\_1} & 0.2673 & 0.0057 & 0.0130 & $+56.1$\% \\
\addlinespace[2pt]
\texttt{bank\_marketing} & 30 & \meth{richardson\_1} & 0.1765 & 0.0246 & 0.0389 & $+36.7$\% \\
\texttt{bank\_marketing} & 60 & \meth{richardson\_1} & 0.1664 & \textbf{0.0348} & 0.0161 & $-115.6$\% \\
\texttt{bank\_marketing} & 90 & \meth{richardson\_1} & 0.1944 & 0.0067 & 0.0077 & $+12.1$\% \\
\addlinespace[2pt]
\texttt{covertype} & 30 & \meth{richardson\_1} & 0.4605 & 0.1114 & 0.3303 & $+66.3$\% \\
\texttt{covertype} & 60 & \meth{richardson\_1} & 0.3983 & 0.0492 & 0.2122 & $+76.8$\% \\
\texttt{covertype} & 90 & \meth{richardson\_1} & 0.3903 & 0.0412 & 0.1633 & $+74.8$\% \\
\addlinespace[2pt]
\texttt{higgs} & 30 & \meth{rational\_fit} & 0.5185 & 0.0116 & 0.0581 & $+80.0$\% \\
\texttt{higgs} & 60 & \meth{rational\_fit} & 0.5190 & 0.0111 & 0.0311 & $+64.4$\% \\
\texttt{higgs} & 90 & \meth{rational\_fit} & 0.5255 & 0.0046 & 0.0193 & $+76.0$\% \\
\addlinespace[2pt]
\texttt{jannis} & 30 & \meth{rational\_fit} & 0.6312 & 0.0583 & 0.1371 & $+57.5$\% \\
\texttt{jannis} & 60 & \meth{richardson\_1} & 0.6169 & \textbf{0.0727} & 0.0684 & $-6.2$\% \\
\texttt{jannis} & 90 & \meth{richardson\_1} & 0.6731 & 0.0165 & 0.0416 & $+60.4$\% \\
\addlinespace[2pt]
\texttt{miniboone} & 30 & \meth{richardson\_1} & 0.1454 & 0.0083 & 0.1159 & $+92.8$\% \\
\texttt{miniboone} & 60 & \meth{richardson\_1} & 0.1161 & 0.0210 & 0.0534 & $+60.7$\% \\
\texttt{miniboone} & 90 & \meth{richardson\_1} & 0.1275 & 0.0096 & 0.0312 & $+69.1$\% \\
\midrule
\multicolumn{7}{l}{reduction in mean error: $\nobs = 30$: $+65.5$\%, $\nobs = 60$: $+38.9$\%, $\nobs = 90$: $+69.4$\%; all cells $+58.7$\%} \\
\bottomrule
\end{tabular}

\end{adjustbox}
\end{table}
\end{landscape}
\begin{landscape}
\begin{table}[p]
\centering
\caption{\texttt{f08\_real\_perturb\_diagnostic.tex}: Real-data perturbation diagnostic: does the routed method's perturb\_iqr separate failing cells from succeeding ones?}
\label{tab:f08_real_perturb_diagnostic}
\begin{adjustbox}{max width=\linewidth, max totalheight=0.86\textheight}
% AUTO-GENERATED -- do not hand-edit.  generated by scripts/make_paper_tables.py from the run started 2026-10-01T22:16:39 on code d82fd26 (full; 29,256 s; 7 workers)
% Real-data perturbation diagnostic: does the routed method's perturb_iqr separate failing cells from succeeding ones?
% source : results/real_data/real_data_summary_v2.csv (perturb_iqr = IQR of the routed method's prediction over 5 evaluations on 2 %-perturbed windows, crc32 seed per (dataset, depth))
% filter : all 90 cells with a finite perturb_iqr, then pre-minimum (target round <= argmin round of the recorded curve) and post-minimum targets, then by routed method; AUC = P(IQR_fail > IQR_succ) from the Mann-Whitney U statistic (ties count one half); p two-sided
% note   : On the 90-cell grid the perturb_iqr of the routed method weakly separates failing cells (skill >= 1 or negative improvement; n = 17) from succeeding ones (n = 73) -- but in the INVERSE direction (failing cells have the LOWER IQR): AUC = 0.317 with a higher IQR read as 'failure', two-sided Mann-Whitney p = 0.019; ordering: failing < succeeding (median IQR 0.0013 vs 0.0054). With skill >= 1 alone: AUC = 0.317, p = 0.019, failing < succeeding (n_fail = 17). With negative improvement alone: AUC = 0.317, p = 0.019, failing < succeeding (n_fail = 17). Within pre-minimum targets only: AUC = 0.198, p = 0.175, failing < succeeding (n_fail = 2 of 65). Within post-minimum targets only: AUC = 0.433, p = 0.598, failing < succeeding (n_fail = 15 of 25). Routed richardson_1 only: AUC = 0.219, p = 0.015, failing < succeeding (n_fail = 7 of 54). Routed rational_fit only: AUC = 0.300, p = 0.069, failing < succeeding (n_fail = 10 of 36). The two failure definitions select the same cells.
\begin{tabular}{llrrrrrrl}
\toprule
Cells & failure definition & $n_{\mathrm{fail}}$ & $n_{\mathrm{succ}}$ & med.\ IQR$_{\mathrm{fail}}$ & med.\ IQR$_{\mathrm{succ}}$ & AUC & $p$ & ordering \\
\midrule
all 90 & skill $\geq 1$ & 17 & 73 & 0.0013 & 0.0054 & 0.317 & 0.019 & failing $<$ succeeding \\
all 90 & improvement $< 0$ & 17 & 73 & 0.0013 & 0.0054 & 0.317 & 0.019 & failing $<$ succeeding \\
all 90 & either & 17 & 73 & 0.0013 & 0.0054 & 0.317 & 0.019 & failing $<$ succeeding \\
\addlinespace[2pt]
pre-minimum targets (65) & either & 2 & 63 & 0.0010 & 0.0056 & 0.198 & 0.175 & failing $<$ succeeding \\
post-minimum targets (25) & either & 15 & 10 & 0.0013 & 0.0028 & 0.433 & 0.598 & failing $<$ succeeding \\
\addlinespace[2pt]
routed \meth{richardson\_1} (54) & either & 7 & 47 & 0.0008 & 0.0035 & 0.219 & 0.015 & failing $<$ succeeding \\
routed \meth{rational\_fit} (36) & either & 10 & 26 & 0.0023 & 0.0083 & 0.300 & 0.069 & failing $<$ succeeding \\
\bottomrule
\end{tabular}

\end{adjustbox}
\end{table}
\end{landscape}
\begin{landscape}
\begin{table}[p]
\centering
\caption{\texttt{f09a\_selectors\_phase3.tex}: Phase 3 selectors: conditional on validity; validity rate alongside}
\label{tab:f09a_selectors_phase3}
\begin{adjustbox}{max width=\linewidth, max totalheight=0.86\textheight}
% AUTO-GENERATED -- do not hand-edit.  generated by scripts/make_paper_tables.py at code d82fd26-dirty, results as of commit 2232d3d
% Phase 3 selectors: conditional on validity; validity rate alongside
% source : results/phase3/phase3_selector_comparison.csv (core regimes, capped cells excluded; src.panels.error_panel over the chosen records)
% filter : all selectors, per stratum; a selector picks one method per cell and is scored on every record of that cell (the chosen method's record; an invalid choice stays invalid, no fallback); valid rate and n_valid/n_total over all records; med. err, q25--q75 and p90 over the valid records only; win vs last = fraction of all records where the chosen record is valid and below the last-value error; the cascades' thresholds are fixed constants, so the per-regime panel (phase3_regime_results.csv) is the per-family evidence and the held-out families of Phase 5a the out-of-sample test
\begin{tabular}{lrrrrrr}
\toprule
Selector & valid rate & $n_{\mathrm{valid}}/n_{\mathrm{total}}$ & med.\ err & q25--q75 & p90 & win vs last \\
\midrule
\multicolumn{7}{l}{\textit{$g = 0.5$}; 1170 cells, 0 capped cells excluded} \\
oracle \textit{(hindsight: lowest cell-median error)} & 0.999 & 23365/23400 & 0.0079 & 0.0019--0.0253 & 0.0497 & 0.945 \\
Phase-2 cascade & 1.000 & 23400/23400 & 0.0235 & 0.0060--0.0629 & 0.1270 & 0.796 \\
enhanced cascade & 0.994 & 23268/23400 & 0.0235 & 0.0064--0.0585 & 0.0972 & 0.823 \\
fixed \meth{rational\_fit} & 1.000 & 23400/23400 & 0.0209 & 0.0059--0.0515 & 0.0783 & 0.861 \\
fixed \meth{richardson\_1} & 1.000 & 23400/23400 & 0.0231 & 0.0060--0.0643 & 0.1345 & 0.785 \\
fixed \meth{single\_exp\_fit} & 1.000 & 23400/23400 & 0.0383 & 0.0081--0.0820 & 0.1167 & 0.756 \\
fixed \meth{last\_value} (the trivial floor) & 1.000 & 23400/23400 & 0.0684 & 0.0225--0.1046 & 0.1525 & 0.000 \\
\addlinespace[2pt]
\multicolumn{7}{l}{\textit{$g = 0.1$} (headline); 1020 cells, 150 capped cells excluded} \\
oracle \textit{(hindsight: lowest cell-median error)} & 0.999 & 20374/20400 & 0.0119 & 0.0029--0.0392 & 0.0792 & 0.939 \\
Phase-2 cascade & 1.000 & 20400/20400 & 0.0348 & 0.0081--0.0907 & 0.1615 & 0.867 \\
enhanced cascade & 0.994 & 20275/20400 & 0.0638 & 0.0179--0.1260 & 0.1896 & 0.782 \\
fixed \meth{rational\_fit} & 1.000 & 20400/20400 & 0.0308 & 0.0094--0.0698 & 0.1234 & 0.884 \\
fixed \meth{richardson\_1} & 1.000 & 20400/20400 & 0.0353 & 0.0077--0.0982 & 0.1884 & 0.843 \\
fixed \meth{single\_exp\_fit} & 1.000 & 20400/20400 & 0.0615 & 0.0109--0.1376 & 0.2086 & 0.802 \\
fixed \meth{last\_value} (the trivial floor) & 1.000 & 20400/20400 & 0.1105 & 0.0339--0.1831 & 0.2769 & 0.000 \\
\addlinespace[2pt]
\multicolumn{7}{l}{\textit{$g = 0.02$}; 960 cells, 210 capped cells excluded} \\
oracle \textit{(hindsight: lowest cell-median error)} & 0.999 & 19179/19200 & 0.0126 & 0.0034--0.0357 & 0.0846 & 0.929 \\
Phase-2 cascade & 1.000 & 19200/19200 & 0.0351 & 0.0086--0.0953 & 0.1605 & 0.855 \\
enhanced cascade & 0.994 & 19082/19200 & 0.0654 & 0.0146--0.1394 & 0.2171 & 0.756 \\
fixed \meth{rational\_fit} & 1.000 & 19200/19200 & 0.0314 & 0.0099--0.0691 & 0.1361 & 0.873 \\
fixed \meth{richardson\_1} & 1.000 & 19200/19200 & 0.0357 & 0.0084--0.0954 & 0.1648 & 0.854 \\
fixed \meth{single\_exp\_fit} & 1.000 & 19200/19200 & 0.0691 & 0.0100--0.1559 & 0.2328 & 0.790 \\
fixed \meth{last\_value} (the trivial floor) & 1.000 & 19200/19200 & 0.1215 & 0.0338--0.2022 & 0.3077 & 0.000 \\
\bottomrule
\end{tabular}

\end{adjustbox}
\end{table}
\end{landscape}
\begin{landscape}
\begin{table}[p]
\centering
\caption{\texttt{f09b\_ensemble\_phase5a\_core.tex}: Phase 5a selectors and ensembles, core curve families: validity rate and n\_valid/n\_total next to the median error and median skill per stratum (mean error at the headline stratum), conditional on validity; oracle\_52 vs constant\_oracle vs the fixed defaults}
\label{tab:f09b_ensemble_phase5a_core}
\begin{adjustbox}{max width=\linewidth, max totalheight=0.86\textheight}
% AUTO-GENERATED -- do not hand-edit.  generated by scripts/make_paper_tables.py at code d82fd26-dirty, results as of commit 2232d3d
% Phase 5a selectors and ensembles, core regimes: validity rate and n_valid/n_total next to the median error and median skill per stratum (mean error at the headline stratum), conditional on validity; oracle_52 vs constant_oracle vs the fixed defaults
% source : results/phase5a/phase5a_ensemble.csv (core regimes; the Phase-5a depth grid; capped cells excluded; src.panels.error_panel over each selector's records)
% filter : all selectors; a selector's record on a cell is the record of the method it chose or the ensemble value, an invalid choice stays invalid; rho_V = validity rate over all records, n_v/n = n_valid/n_total; mean err and med. err over the valid records only (conditional on validity); skill = hindsight best-of-four (strict): err / best-of-four trivial error per record, median over the valid records
\begin{tabular}{lrrrrrrrrrrrrr}
\toprule
Selector & \multicolumn{4}{c}{$g = 0.5$} & \multicolumn{5}{c}{$g = 0.1$} & \multicolumn{4}{c}{$g = 0.02$} \\
\cmidrule(lr){2-5}\cmidrule(lr){6-10}\cmidrule(lr){11-14}
 & $\rhoV$ & $n_v/n$ & med.\ err & med.\ skill & $\rhoV$ & $n_v/n$ & mean err & med.\ err & med.\ skill & $\rhoV$ & $n_v/n$ & med.\ err & med.\ skill \\
\midrule
oracle over the 52-method pool & 1.000 & 4320/4320 & 0.0011 & 0.042 & 1.000 & 3780/3780 & 0.0065 & 0.0023 & 0.076 & 1.000 & 3540/3540 & 0.0027 & 0.100 \\
\meth{constant\_oracle} \textit{(reference)} & 1.000 & 4320/4320 & 0.0679 & 1.009 & 1.000 & 3780/3780 & 0.0135 & 0.0124 & 0.207 & 1.000 & 3540/3540 & 0.0022 & 0.037 \\
oracle over the 9-method pool & 1.000 & 4320/4320 & 0.0053 & 0.183 & 1.000 & 3780/3780 & 0.0240 & 0.0080 & 0.282 & 1.000 & 3540/3540 & 0.0087 & 0.372 \\
\addlinespace[2pt]
fixed \meth{rational\_fit} & 1.000 & 4320/4320 & 0.0256 & 0.581 & 1.000 & 3780/3780 & 0.0531 & 0.0333 & 1.000 & 1.000 & 3540/3540 & 0.0356 & 1.165 \\
fixed \meth{richardson\_1} & 1.000 & 4320/4320 & 0.0258 & 0.627 & 1.000 & 3780/3780 & 0.0830 & 0.0379 & 1.003 & 1.000 & 3540/3540 & 0.0394 & 1.097 \\
Phase-2 cascade & 1.000 & 4320/4320 & 0.0271 & 0.671 & 1.000 & 3780/3780 & 0.0676 & 0.0382 & 1.000 & 1.000 & 3540/3540 & 0.0395 & 1.066 \\
\addlinespace[2pt]
equal ensemble (9) & 1.000 & 4320/4320 & 0.0406 & 0.925 & 1.000 & 3780/3780 & 0.0967 & 0.0645 & 1.391 & 1.000 & 3540/3540 & 0.0692 & 1.654 \\
diagnostic-weighted ensemble (9) & 1.000 & 4320/4320 & 0.0397 & 0.855 & 1.000 & 3780/3780 & 0.0944 & 0.0622 & 1.493 & 1.000 & 3540/3540 & 0.0660 & 1.717 \\
equal ensemble (52) & 1.000 & 4320/4320 & 0.0420 & 1.053 & 1.000 & 3780/3780 & 0.1068 & 0.0659 & 1.513 & 1.000 & 3540/3540 & 0.0778 & 1.762 \\
diagnostic-weighted ensemble (52) & 1.000 & 4320/4320 & 0.0309 & 0.812 & 1.000 & 3780/3780 & 0.1030 & 0.0569 & 1.496 & 1.000 & 3540/3540 & 0.0634 & 1.898 \\
capped diagnostic-weighted (52) & 1.000 & 4320/4320 & 0.0309 & 0.812 & 1.000 & 3780/3780 & 0.1030 & 0.0569 & 1.496 & 1.000 & 3540/3540 & 0.0634 & 1.898 \\
equal ensemble (safe) & 1.000 & 4320/4320 & 0.0436 & 1.059 & 1.000 & 3780/3780 & 0.1065 & 0.0675 & 1.517 & 1.000 & 3540/3540 & 0.0817 & 1.808 \\
diagnostic-weighted (safe) & 1.000 & 4320/4320 & 0.0344 & 0.859 & 1.000 & 3780/3780 & 0.0995 & 0.0568 & 1.384 & 1.000 & 3540/3540 & 0.0617 & 1.718 \\
capped diagnostic-weighted (safe) & 1.000 & 4320/4320 & 0.0344 & 0.859 & 1.000 & 3780/3780 & 0.0995 & 0.0568 & 1.384 & 1.000 & 3540/3540 & 0.0617 & 1.718 \\
threshold ensemble, IQR $\leq 0.10$ & 1.000 & 4320/4320 & 0.0444 & 1.060 & 1.000 & 3780/3780 & 0.1069 & 0.0681 & 1.526 & 1.000 & 3540/3540 & 0.0825 & 1.812 \\
threshold ensemble, IQR $\leq 0.50$ & 1.000 & 4320/4320 & 0.0437 & 1.059 & 1.000 & 3780/3780 & 0.1067 & 0.0675 & 1.517 & 1.000 & 3540/3540 & 0.0817 & 1.817 \\
threshold ensemble, no threshold & 1.000 & 4320/4320 & 0.0436 & 1.059 & 1.000 & 3780/3780 & 0.1065 & 0.0675 & 1.517 & 1.000 & 3540/3540 & 0.0817 & 1.808 \\
\addlinespace[2pt]
\meth{constant\_assumed} & 1.000 & 4320/4320 & 0.1345 & 2.543 & 1.000 & 3780/3780 & 0.1213 & 0.0617 & 1.000 & 1.000 & 3540/3540 & 0.0496 & 1.000 \\
\meth{last\_value} & 1.000 & 4320/4320 & 0.0715 & 1.030 & 1.000 & 3780/3780 & 0.1327 & 0.1185 & 2.284 & 1.000 & 3540/3540 & 0.1310 & 2.618 \\
\meth{window\_min} & 1.000 & 4320/4320 & 0.0579 & 1.000 & 1.000 & 3780/3780 & 0.1221 & 0.1012 & 1.810 & 1.000 & 3540/3540 & 0.1221 & 2.497 \\
\meth{window\_mean} & 1.000 & 4320/4320 & 0.1144 & 2.147 & 1.000 & 3780/3780 & 0.1995 & 0.1670 & 3.834 & 1.000 & 3540/3540 & 0.1946 & 4.181 \\
\bottomrule
\end{tabular}

\end{adjustbox}
\end{table}
\end{landscape}
\begin{landscape}
\begin{table}[p]
\centering
\caption{\texttt{f09c\_ensemble\_phase5a\_holdout.tex}: Phase 5a selectors and ensembles, out-of-design curve families: validity rate and n\_valid/n\_total next to the median error and median skill per stratum (mean error at the headline stratum), conditional on validity; oracle\_52 vs constant\_oracle vs the fixed defaults}
\label{tab:f09c_ensemble_phase5a_holdout}
\begin{adjustbox}{max width=\linewidth, max totalheight=0.86\textheight}
% AUTO-GENERATED -- do not hand-edit.  generated by scripts/make_paper_tables.py at code d82fd26-dirty, results as of commit 2232d3d
% Phase 5a selectors and ensembles, held-out regimes: validity rate and n_valid/n_total next to the median error and median skill per stratum (mean error at the headline stratum), conditional on validity; oracle_52 vs constant_oracle vs the fixed defaults
% source : results/phase5a/phase5a_ensemble_holdout.csv (held-out regimes; the Phase-5a depth grid; capped cells excluded; src.panels.error_panel over each selector's records)
% filter : all selectors; a selector's record on a cell is the record of the method it chose or the ensemble value, an invalid choice stays invalid; rho_V = validity rate over all records, n_v/n = n_valid/n_total; mean err and med. err over the valid records only (conditional on validity); skill = hindsight best-of-four (strict): err / best-of-four trivial error per record, median over the valid records
\begin{tabular}{lrrrrrrrrrrrrr}
\toprule
Selector & \multicolumn{4}{c}{$g = 0.5$} & \multicolumn{5}{c}{$g = 0.1$} & \multicolumn{4}{c}{$g = 0.02$} \\
\cmidrule(lr){2-5}\cmidrule(lr){6-10}\cmidrule(lr){11-14}
 & $\rhoV$ & $n_v/n$ & med.\ err & med.\ skill & $\rhoV$ & $n_v/n$ & mean err & med.\ err & med.\ skill & $\rhoV$ & $n_v/n$ & med.\ err & med.\ skill \\
\midrule
oracle over the 52-method pool & 1.000 & 1440/1440 & 0.0009 & 0.031 & 1.000 & 1434/1434 & 0.0091 & 0.0024 & 0.058 & 1.000 & 1263/1263 & 0.0027 & 0.073 \\
\meth{constant\_oracle} \textit{(reference)} & 1.000 & 1440/1440 & 0.0617 & 1.016 & 1.000 & 1434/1434 & 0.0139 & 0.0124 & 0.203 & 1.000 & 1263/1263 & 0.0026 & 0.045 \\
oracle over the 9-method pool & 1.000 & 1440/1440 & 0.0045 & 0.136 & 1.000 & 1434/1434 & 0.0349 & 0.0078 & 0.250 & 1.000 & 1263/1263 & 0.0106 & 0.311 \\
\addlinespace[2pt]
fixed \meth{rational\_fit} & 1.000 & 1440/1440 & 0.0153 & 0.385 & 1.000 & 1434/1434 & 0.0635 & 0.0369 & 0.611 & 1.000 & 1263/1263 & 0.0429 & 0.669 \\
fixed \meth{richardson\_1} & 1.000 & 1440/1440 & 0.0217 & 0.639 & 1.000 & 1434/1434 & 0.0781 & 0.0351 & 0.783 & 1.000 & 1263/1263 & 0.0462 & 0.937 \\
Phase-2 cascade & 1.000 & 1440/1440 & 0.0226 & 0.644 & 1.000 & 1434/1434 & 0.0739 & 0.0359 & 0.775 & 1.000 & 1263/1263 & 0.0450 & 0.931 \\
\addlinespace[2pt]
equal ensemble (9) & 1.000 & 1440/1440 & 0.0305 & 0.814 & 1.000 & 1434/1434 & 0.0946 & 0.0624 & 1.128 & 1.000 & 1263/1263 & 0.0752 & 1.323 \\
diagnostic-weighted ensemble (9) & 1.000 & 1440/1440 & 0.0363 & 0.834 & 1.000 & 1434/1434 & 0.0942 & 0.0621 & 1.075 & 1.000 & 1263/1263 & 0.0709 & 1.252 \\
equal ensemble (52) & 1.000 & 1440/1440 & 0.0265 & 1.026 & 1.000 & 1434/1434 & 0.1031 & 0.0592 & 1.309 & 1.000 & 1263/1263 & 0.0788 & 1.533 \\
diagnostic-weighted ensemble (52) & 1.000 & 1440/1440 & 0.0262 & 0.713 & 1.000 & 1434/1434 & 0.1040 & 0.0479 & 0.986 & 1.000 & 1263/1263 & 0.0614 & 1.206 \\
capped diagnostic-weighted (52) & 1.000 & 1440/1440 & 0.0262 & 0.713 & 1.000 & 1434/1434 & 0.1040 & 0.0479 & 0.986 & 1.000 & 1263/1263 & 0.0614 & 1.206 \\
equal ensemble (safe) & 1.000 & 1440/1440 & 0.0281 & 1.033 & 1.000 & 1434/1434 & 0.1047 & 0.0614 & 1.316 & 1.000 & 1263/1263 & 0.0791 & 1.525 \\
diagnostic-weighted (safe) & 1.000 & 1440/1440 & 0.0274 & 0.794 & 1.000 & 1434/1434 & 0.1025 & 0.0488 & 0.985 & 1.000 & 1263/1263 & 0.0637 & 1.194 \\
capped diagnostic-weighted (safe) & 1.000 & 1440/1440 & 0.0274 & 0.794 & 1.000 & 1434/1434 & 0.1025 & 0.0488 & 0.985 & 1.000 & 1263/1263 & 0.0637 & 1.194 \\
threshold ensemble, IQR $\leq 0.10$ & 1.000 & 1440/1440 & 0.0286 & 1.041 & 1.000 & 1434/1434 & 0.1057 & 0.0636 & 1.329 & 1.000 & 1263/1263 & 0.0807 & 1.526 \\
threshold ensemble, IQR $\leq 0.50$ & 1.000 & 1440/1440 & 0.0281 & 1.035 & 1.000 & 1434/1434 & 0.1055 & 0.0623 & 1.325 & 1.000 & 1263/1263 & 0.0791 & 1.525 \\
threshold ensemble, no threshold & 1.000 & 1440/1440 & 0.0281 & 1.033 & 1.000 & 1434/1434 & 0.1047 & 0.0614 & 1.316 & 1.000 & 1263/1263 & 0.0791 & 1.525 \\
\addlinespace[2pt]
\meth{constant\_assumed} & 1.000 & 1440/1440 & 0.1422 & 2.603 & 1.000 & 1434/1434 & 0.1182 & 0.0675 & 1.000 & 1.000 & 1263/1263 & 0.0564 & 1.000 \\
\meth{last\_value} & 1.000 & 1440/1440 & 0.0589 & 1.000 & 1.000 & 1434/1434 & 0.1262 & 0.1078 & 1.749 & 1.000 & 1263/1263 & 0.1277 & 2.163 \\
\meth{window\_min} & 1.000 & 1440/1440 & 0.0498 & 1.000 & 1.000 & 1434/1434 & 0.1182 & 0.0902 & 1.222 & 1.000 & 1263/1263 & 0.1190 & 1.546 \\
\meth{window\_mean} & 1.000 & 1440/1440 & 0.1300 & 2.307 & 1.000 & 1434/1434 & 0.2030 & 0.1857 & 3.600 & 1.000 & 1263/1263 & 0.2108 & 4.399 \\
\bottomrule
\end{tabular}

\end{adjustbox}
\end{table}
\end{landscape}
\begin{landscape}
\begin{table}[p]
\centering
\caption{\texttt{f09d\_ablation\_phase5a.tex}: Phase 5a ablation: paired mean error differences between selectors per stratum, with the count of records where both are valid}
\label{tab:f09d_ablation_phase5a}
\begin{adjustbox}{max width=\linewidth, max totalheight=0.86\textheight}
% AUTO-GENERATED -- do not hand-edit.  generated by scripts/make_paper_tables.py from the run started 2026-10-01T22:16:39 on code d82fd26 (full; 29,256 s; 7 workers)
% Phase 5a ablation: paired mean error differences between selectors per stratum, with the count of records where both are valid
% source : results/phase5a/phase5a_ablation.csv (core regimes, capped excluded)
% filter : all comparisons, printed as '<first> vs <second>' with the selector labels of f09b (the (first, second) pairs are the comparisons list of phases/phase5a.py; the comparison names of the CSV are threshold_vs_equal_52, threshold_vs_rational, capped_diag_vs_equal, diag_vs_equal_52, weighting_gain_9, filter_benefit, pool_expansion_gain, threshold_vs_constant_assumed, rational_vs_window_min); mean_improvement = mean over the records where both selectors are valid of (error of the second selector - error of the first); n_both/n = those records over all records at the headline stratum
\begin{tabular}{lrrrr}
\toprule
Comparison (positive = first is better) & $g = 0.5$ & $g = 0.1$ & $g = 0.02$ & $n_{\mathrm{both}}/n$ \\
\midrule
threshold ensemble, IQR $\leq 0.10$ vs equal ensemble (52) & $-0.0009$ & $-0.0001$ & $-0.0003$ & 3780/3780 \\
threshold ensemble, IQR $\leq 0.10$ vs fixed \meth{rational\_fit} & $-0.0239$ & $-0.0538$ & $-0.0636$ & 3780/3780 \\
capped diagnostic-weighted (52) vs equal ensemble (52) & $-0.0021$ & $+0.0038$ & $+0.0089$ & 3780/3780 \\
diagnostic-weighted ensemble (52) vs equal ensemble (52) & $-0.0021$ & $+0.0038$ & $+0.0089$ & 3780/3780 \\
diagnostic-weighted ensemble (9) vs equal ensemble (9) & $-0.0013$ & $+0.0023$ & $+0.0087$ & 3780/3780 \\
threshold ensemble, IQR $\leq 0.10$ vs threshold ensemble, no threshold & $-0.0001$ & $-0.0004$ & $-0.0005$ & 3780/3780 \\
equal ensemble (52) vs equal ensemble (9) & $-0.0057$ & $-0.0101$ & $-0.0120$ & 3780/3780 \\
threshold ensemble, IQR $\leq 0.10$ vs \meth{constant\_assumed} & $+0.1148$ & $+0.0144$ & $-0.0093$ & 3780/3780 \\
fixed \meth{rational\_fit} vs \meth{window\_min} & $+0.0311$ & $+0.0690$ & $+0.0828$ & 3780/3780 \\
\bottomrule
\end{tabular}

\end{adjustbox}
\end{table}
\end{landscape}
\begin{landscape}
\begin{table}[p]
\centering
\caption{\texttt{f10a\_diagnostics\_correlations.tex}: Phase 4 diagnostics: correlation of perturb\_IQR and shift\_IQR with error, core pooled and out-of-design}
\label{tab:f10a_diagnostics_correlations}
\begin{adjustbox}{max width=\linewidth, max totalheight=0.86\textheight}
% AUTO-GENERATED -- do not hand-edit.  generated by scripts/make_paper_tables.py at code d82fd26-dirty, results as of commit 2232d3d
% Phase 4 diagnostics: correlation of perturb_IQR and shift_IQR with error, core pooled and held-out
% source : results/phase4/phase4_diagnostic_correlations.csv (core regimes pooled over obs 30/60/90/120, noise, seeds, strata; capped excluded)
% source : results/phase4/phase4_diagnostic_correlations_holdout.csv (held-out regimes)
% filter : all rows; Spearman correlation between the diagnostic and the absolute error over valid records; -- = diagnostic undefined (a NaN Spearman r in the stored table: the diagnostic is constant or undefined for that method)
\begin{tabular}{lrrrrrr}
\toprule
Method & \multicolumn{3}{c}{core (pooled)} & \multicolumn{3}{c}{held-out} \\
\cmidrule(lr){2-4}\cmidrule(lr){5-7}
 & $r$(\diag{perturb\_IQR}) & $r$(\diag{shift\_IQR}) & $n$ & $r$(\diag{perturb\_IQR}) & $r$(\diag{shift\_IQR}) & $n$ \\
\midrule
\textit{all methods pooled} & 0.383 & 0.131 & 103648 & 0.316 & 0.122 & 37034 \\
\addlinespace[2pt]
\meth{last\_value} & 0.518 & -- & 11640 & 0.560 & -- & 4137 \\
\meth{richardson\_1} & 0.322 & 0.583 & 11640 & 0.110 & 0.536 & 4137 \\
\meth{richardson\_a10} & 0.351 & 0.779 & 11640 & 0.252 & 0.716 & 4137 \\
\meth{single\_exp\_fit} & 0.325 & 0.406 & 11640 & 0.374 & 0.444 & 4137 \\
\meth{rational\_fit} & 0.290 & 0.487 & 11640 & 0.288 & 0.356 & 4137 \\
\meth{pade\_22} & 0.413 & 0.417 & 11260 & 0.282 & 0.347 & 4012 \\
\meth{log\_linear} & 0.485 & 0.629 & 11250 & 0.407 & 0.625 & 4075 \\
\meth{levin\_t2} & 0.454 & -- & 11301 & 0.409 & -- & 4131 \\
\meth{anderson\_1} & 0.452 & -- & 11637 & 0.470 & -- & 4131 \\
\bottomrule
\end{tabular}

\end{adjustbox}
\end{table}
\end{landscape}
\begin{landscape}
\begin{table}[p]
\centering
\caption{\texttt{f10b\_diagnostics\_depth.tex}: Phase 4: perturb\_IQR reliability by observation depth at the headline stratum}
\label{tab:f10b_diagnostics_depth}
\begin{adjustbox}{max width=\linewidth, max totalheight=0.86\textheight}
% AUTO-GENERATED -- do not hand-edit.  generated by scripts/make_paper_tables.py from the run started 2026-10-01T22:16:39 on code d82fd26 (full; 29,256 s; 7 workers)
% Phase 4: perturb_IQR reliability by observation depth at the headline stratum
% source : results/phase4/phase4_obs_reliability.csv (core + out-of-design, per depth)
% filter : diagnostic == perturb_iqr, target_g == 0.1; Spearman r between perturb_IQR and error at each observation depth
\begin{tabular}{lrrrr}
\toprule
Method & $\nobs = 30$ & $\nobs = 60$ & $\nobs = 90$ & $\nobs = 120$ \\
\midrule
\meth{last\_value} & 0.563 & 0.552 & 0.429 & 0.417 \\
\meth{richardson\_1} & 0.473 & 0.328 & 0.310 & 0.169 \\
\meth{richardson\_a10} & 0.497 & 0.510 & 0.481 & 0.433 \\
\meth{single\_exp\_fit} & 0.277 & 0.375 & 0.340 & 0.410 \\
\meth{rational\_fit} & 0.355 & 0.260 & 0.128 & 0.222 \\
\meth{pade\_22} & 0.370 & 0.423 & 0.513 & 0.301 \\
\meth{log\_linear} & 0.559 & 0.453 & 0.575 & 0.524 \\
\meth{levin\_t2} & 0.468 & 0.477 & 0.416 & 0.350 \\
\meth{anderson\_1} & 0.425 & 0.487 & 0.410 & 0.341 \\
\bottomrule
\end{tabular}

\end{adjustbox}
\end{table}
\end{landscape}
\begin{landscape}
\begin{table}[p]
\centering
\caption{\texttt{f10c\_diagnostics\_ensemble.tex}: Phase 4 selectors with the trivial references at the headline stratum: validity rate and n\_valid/n\_total next to the error and skill, conditional on validity; validity rate alongside}
\label{tab:f10c_diagnostics_ensemble}
\begin{adjustbox}{max width=\linewidth, max totalheight=0.86\textheight}
% AUTO-GENERATED -- do not hand-edit.  generated by scripts/make_paper_tables.py from the run started 2026-10-01T22:16:39 on code d82fd26 (full; 29,256 s; 7 workers)
% Phase 4 selectors with the trivial references at the headline stratum: validity rate and n_valid/n_total next to the error and skill, conditional on validity; validity rate alongside
% source : results/phase4/phase4_ensemble.csv (core regimes, headline stratum, capped excluded; src.panels.error_panel over each selector's records)
% filter : all rows; a selector's record on a window is the record of the method it chose or the combined value, an invalid choice stays invalid; rho_V = validity rate over all windows, n_v/n = n_valid/n_total; mean err and med. err over the valid records only (conditional on validity); selectors printed with the labels of f09b: oracle = oracle over the 9-method pool; phase2_proxy = the lower error of richardson_1 and rational_fit (a hindsight proxy of the Phase-2 cascade); diag_ensemble = perturb_IQR-weighted ensemble of the 9-method pool (eight accelerators plus the last observed value), weights 1/(perturb_IQR + eps) over the paired perturbations
\begin{tabular}{lrrrrrr}
\toprule
Selector & $\rhoV$ & $n_v/n$ & mean err & med.\ err & med.\ skill & win vs last \\
\midrule
oracle over the 9-method pool & 1.000 & 3780/3780 & 0.0240 & 0.0080 & 0.282 & 0.983 \\
fixed \meth{rational\_fit} & 1.000 & 3780/3780 & 0.0531 & 0.0333 & 1.000 & 0.884 \\
fixed \meth{richardson\_1} & 1.000 & 3780/3780 & 0.0830 & 0.0379 & 1.003 & 0.865 \\
lower error of \meth{richardson\_1} and \meth{rational\_fit} (hindsight) & 1.000 & 3780/3780 & 0.0436 & 0.0246 & 0.715 & 0.922 \\
diagnostic-weighted ensemble (9) & 1.000 & 3780/3780 & 0.0944 & 0.0622 & 1.493 & 0.899 \\
\meth{constant\_assumed} & 1.000 & 3780/3780 & 0.1213 & 0.0617 & 1.000 & 0.586 \\
\meth{last\_value} & 1.000 & 3780/3780 & 0.1327 & 0.1185 & 2.284 & 0.000 \\
\meth{window\_mean} & 1.000 & 3780/3780 & 0.1995 & 0.1670 & 3.834 & 0.115 \\
\meth{window\_min} & 1.000 & 3780/3780 & 0.1221 & 0.1012 & 1.810 & 0.456 \\
\bottomrule
\end{tabular}

\end{adjustbox}
\end{table}
\end{landscape}
\begin{landscape}
\begin{table}[p]
\centering
\caption{\texttt{f10d\_diagnostics\_filter.tex}: Phase 4: the perturb\_IQR screen applied to the Phase-2 cascade, conditional on validity; validity rate alongside}
\label{tab:f10d_diagnostics_filter}
\begin{adjustbox}{max width=\linewidth, max totalheight=0.86\textheight}
% AUTO-GENERATED -- do not hand-edit.  generated by scripts/make_paper_tables.py from the run started 2026-10-01T22:16:39 on code d82fd26 (full; 29,256 s; 7 workers)
% Phase 4: the perturb_IQR screen applied to the Phase-2 cascade, conditional on validity; validity rate alongside
% source : results/phase4/phase4_cascade_filter.csv (core regimes, headline stratum; src.panels.error_panel over the chosen records)
% filter : all rows; a window whose routed prediction has perturb_IQR above the threshold is routed to the last observed value (last_value), which is part of the screened method; the chosen record is then taken as it is (an invalid choice stays invalid, no evaluation fallback); rho_V = validity rate, n_v/n = n_valid/n_total; med. err and mean err over the valid records (conditional on validity); win vs last = fraction of all windows where the chosen record is valid and below the last-value error
\begin{tabular}{lrrrrr}
\toprule
Screen & $\rhoV$ & $n_v/n$ & med.\ err & mean err & win vs last \\
\midrule
no\_filter & 1.000 & 3780/3780 & 0.0382 & 0.0676 & 0.869 \\
perturb\_iqr>1.0 & 1.000 & 3780/3780 & 0.0382 & 0.0676 & 0.869 \\
perturb\_iqr>0.5 & 1.000 & 3780/3780 & 0.0382 & 0.0676 & 0.869 \\
perturb\_iqr>0.2 & 1.000 & 3780/3780 & 0.0381 & 0.0675 & 0.868 \\
perturb\_iqr>0.1 & 1.000 & 3780/3780 & 0.0382 & 0.0677 & 0.864 \\
perturb\_iqr>0.05 & 1.000 & 3780/3780 & 0.0389 & 0.0687 & 0.848 \\
perturb\_iqr>0.02 & 1.000 & 3780/3780 & 0.0403 & 0.0707 & 0.821 \\
\bottomrule
\end{tabular}

\end{adjustbox}
\end{table}
\end{landscape}
\begin{landscape}
\begin{table}[p]
\centering
\caption{\texttt{f11\_capped\_block.tex}: Capped block: every capped (curve family x stratum) cell of the main run with the achieved gap fraction, and the capped-cell counts of the depth grids}
\label{tab:f11_capped_block}
\begin{adjustbox}{max width=\linewidth, max totalheight=0.86\textheight}
% AUTO-GENERATED -- do not hand-edit.  generated by scripts/make_paper_tables.py from the run started 2026-10-01T22:16:39 on code d82fd26 (full; 29,256 s; 7 workers)
% Capped block: every capped (regime x stratum) cell of the main run with the achieved gap fraction, and the capped-cell counts of the depth grids
% source : results/phase1/phase1_capped.csv (the capped block of the main run: every method x capped cell)
% source : results/phase1/phase1_horizons.csv (n_f search per regime and stratum at n_obs = 90; seed 0 for seed-dependent shapes)
% source : results/phase2/phase2_capped.csv; results/phase4/phase4_capped.csv; results/phase5a/phase5a_capped.csv; results/phase3/phase3_capped_cells.csv
% filter : capped == 1; n_f = HORIZON_N_CAP = 50000 unless the shape is seed-dependent (median over seeds shown); achieved g = gap(n_f) / gap(n_obs) actually reached; these cells are excluded from every pooled table and figure
\begin{tabular}{lllrrrrl}
\toprule
Regime & set & $g$ & $\nf$ & achieved $g$ & gap$(\nobs)$ & gap$(\nf)$ & note \\
\midrule
\texttt{broken\_power\_law} & core & 0.02 & 50000 & 0.0802 & 0.1317 & 0.0106 &  \\
\texttt{log\_oscillatory} & core & 0.1 & 50000 & 0.2993 & 0.2050 & 0.0614 &  \\
\texttt{log\_oscillatory} & core & 0.02 & 50000 & 0.2993 & 0.2050 & 0.0614 &  \\
\texttt{log\_slow} & core & 0.1 & 50000 & 0.4186 & 0.1766 & 0.0739 &  \\
\texttt{log\_slow} & core & 0.02 & 50000 & 0.4186 & 0.1766 & 0.0739 &  \\
\texttt{random\_knots} & held-out & 0.1 & 3850 & 0.1000 & 0.0247 & 0.0025 & seed-dependent shape: capped on some seeds (median $\nf$ shown) \\
\texttt{random\_knots} & held-out & 0.02 & 45170 & 0.0200 & 0.0247 & 0.0008 & seed-dependent shape: capped on some seeds (median $\nf$ shown) \\
\midrule
\multicolumn{8}{l}{Capped cells in the depth grids (regime $\times$ depth $\times$ $g$): Phase 2: 72 over 13 depths; Phase 4: 29 over 4 depths; Phase 5a: 29 over 4 depths; Phase 3: 360 (regime, depth, noise, $g$) cells} \\
\multicolumn{8}{l}{Regimes ever capped in the depth grids: \texttt{broken\_power\_law}, \texttt{damped\_osc\_pow}, \texttt{inv\_sqrt\_log}, \texttt{log\_oscillatory}, \texttt{log\_slow}, \texttt{random\_knots}} \\
\bottomrule
\end{tabular}

\end{adjustbox}
\end{table}
\end{landscape}
\begin{landscape}
\begin{table}[p]
\centering
\caption{\texttt{f12\_validity\_by\_depth.tex}: Validity by observation depth: every method that falls below the rank floor at some depth}
\label{tab:f12_validity_by_depth}
\begin{adjustbox}{max width=\linewidth, max totalheight=0.86\textheight}
% AUTO-GENERATED -- do not hand-edit.  generated by scripts/make_paper_tables.py at code d82fd26-dirty, results as of commit 2232d3d
% Validity by observation depth: every method that falls below the rank floor at some depth
% source : results/phase5a/phase5a_validity_by_depth.csv (method x obs_idx x noise valid rate over every regime, seed and stratum)
% filter : methods whose noise-pooled valid rate is below 0.9 at any depth (bold); last column = the minimum over (depth, noise) cells; dagger = dangerous artifact
\begin{tabular}{lrrrrr}
\toprule
Method & $\nobs = 30$ & $\nobs = 60$ & $\nobs = 90$ & $\nobs = 120$ & min (depth, $\sigma$) \\
\midrule
$\dagger$\meth{neville\_4} & \textbf{0.470} & \textbf{0.334} & \textbf{0.321} & \textbf{0.327} & 0.033 \\
$\dagger$\meth{linear} & \textbf{0.345} & \textbf{0.358} & \textbf{0.551} & \textbf{0.621} & 0.341 \\
$\dagger$\meth{neville\_3} & \textbf{0.648} & \textbf{0.486} & \textbf{0.400} & \textbf{0.406} & 0.111 \\
\meth{richardson\_3} & \textbf{0.480} & \textbf{0.573} & 0.978 & 0.999 & 0.348 \\
$\dagger$\meth{pade\_21} & \textbf{0.628} & \textbf{0.544} & \textbf{0.607} & \textbf{0.664} & 0.235 \\
$\dagger$\meth{pade\_31} & \textbf{0.598} & \textbf{0.575} & \textbf{0.743} & \textbf{0.772} & 0.392 \\
$\dagger$\meth{geom\_avg\_diff} & \textbf{0.587} & \textbf{0.597} & \textbf{0.604} & \textbf{0.619} & 0.471 \\
$\dagger$\meth{neville\_2} & \textbf{0.707} & \textbf{0.682} & \textbf{0.632} & \textbf{0.674} & 0.480 \\
$\dagger$\meth{pade\_32} & \textbf{0.666} & \textbf{0.728} & \textbf{0.632} & \textbf{0.666} & 0.494 \\
\midrule
\multicolumn{6}{l}{9 of 52 accelerators fall below $\rhoV = 0.9$ at some depth; the dangerous artifact (obs 90 only) flags 8 of them} \\
\bottomrule
\end{tabular}

\end{adjustbox}
\end{table}
\end{landscape}
\begin{landscape}
\begin{table}[p]
\centering
\caption{\texttt{f13\_by\_Ltrue\_g0.02.tex}: Phase-1 results by L\textasciicircum{}* tercile at g = 0.02: the deployable trivials, the oracle reference, the top-10 accelerators (by pooled median error at g = 0.1, core) and the 21 classical variants pooled; core | out-of-design}
\label{tab:f13_by_Ltrue_g0.02}
\begin{adjustbox}{max width=\linewidth, max totalheight=0.86\textheight}
% AUTO-GENERATED -- do not hand-edit.  generated by scripts/analyze_by_ltrue.py from the run started 2026-10-01T22:16:39 on code d82fd26 (full; 29,256 s; 7 workers)
% Phase-1 results by L^* tercile at g = 0.02: the deployable trivials, the oracle reference, the top-10 accelerators (by pooled median error at g = 0.1, core) and the 21 classical variants pooled; core | held-out
% source : results/phase1/phase1_by_Ltrue.csv (the per-tercile aggregate this fragment prints; the aggregate of the git-ignored results/phase1/phase1_records.csv, capped == 0)
% filter : target_g == 0.02; log-scale terciles of L^* (L_true) in [0.005, 0.5]: T1 [0.005, 0.0232), T2 [0.0232, 0.1077), T3 [0.1077, 0.5]; med. err = median over valid records; win vs X = mean of the per-record win indicator (invalid never wins); skill vs last = median over valid records of err / err(last_value)
\begin{tabular}{llrrrrrrrrrrrr}
\toprule
$g$ / row & & \multicolumn{4}{c}{T1: $L^{*} \in [0.005, 0.0232)$} & \multicolumn{4}{c}{T2: $[0.0232, 0.1077)$} & \multicolumn{4}{c}{T3: $[0.1077, 0.5]$} \\
\cmidrule(lr){3-6}\cmidrule(lr){7-10}\cmidrule(lr){11-14}
 & set & med.\ err & win/last & win/$\hat L$ & skill/last & med.\ err & win/last & win/$\hat L$ & skill/last & med.\ err & win/last & win/$\hat L$ & skill/last \\
\midrule
\multicolumn{14}{l}{\textit{$g = 0.02$}} \\
predict $\hat L$ (\meth{constant\_assumed}) & core & 0.0125 & 0.859 & -- & 0.156 & 0.0506 & 0.649 & -- & 0.542 & 0.2252 & 0.138 & -- & 2.657 \\
 & held-out & 0.0137 & 0.978 & -- & 0.180 & 0.0535 & 0.739 & -- & 0.599 & 0.2050 & 0.119 & -- & 2.673 \\
\meth{last\_value} & core & 0.0567 & -- & 0.141 & -- & 0.0978 & -- & 0.351 & -- & 0.0970 & -- & 0.862 & -- \\
 & held-out & 0.0934 & -- & 0.022 & -- & 0.0919 & -- & 0.261 & -- & 0.0935 & -- & 0.881 & -- \\
\meth{window\_mean} & core & 0.1298 & 0.102 & 0.125 & 1.395 & 0.1749 & 0.095 & 0.221 & 1.413 & 0.1749 & 0.107 & 0.704 & 1.413 \\
 & held-out & 0.2074 & 0.000 & 0.000 & 1.418 & 0.2077 & 0.000 & 0.130 & 1.463 & 0.2076 & 0.000 & 0.678 & 1.400 \\
\meth{window\_min} & core & 0.0549 & 0.338 & 0.083 & 1.000 & 0.0975 & 0.290 & 0.355 & 1.000 & 0.0970 & 0.287 & 0.864 & 1.000 \\
 & held-out & 0.0930 & 0.333 & 0.066 & 1.000 & 0.0919 & 0.348 & 0.268 & 1.000 & 0.0935 & 0.345 & 0.893 & 1.000 \\
\meth{constant\_oracle} \textit{(ref.)} & core & 0.0011 & 1.000 & 1.000 & 0.020 & 0.0011 & 1.000 & 1.000 & 0.020 & 0.0011 & 1.000 & 1.000 & 0.020 \\
 & held-out & 0.0019 & 1.000 & 1.000 & 0.020 & 0.0019 & 1.000 & 1.000 & 0.020 & 0.0019 & 1.000 & 1.000 & 0.020 \\
\addlinespace[2pt]
\meth{log\_linear} & core & 0.0087 & 0.905 & 0.650 & 0.233 & 0.0255 & 0.799 & 0.881 & 0.352 & 0.0599 & 0.588 & 1.000 & 0.743 \\
 & held-out & 0.0236 & 0.995 & 0.492 & 0.345 & 0.0322 & 0.855 & 0.870 & 0.349 & 0.0522 & 0.593 & 1.000 & 0.506 \\
\meth{richardson\_3} & core & 0.0122 & 0.917 & 0.528 & 0.279 & 0.0199 & 0.827 & 0.838 & 0.269 & 0.0333 & 0.693 & 0.952 & 0.594 \\
 & held-out & 0.0187 & 0.885 & 0.481 & 0.266 & 0.0244 & 0.913 & 0.746 & 0.230 & 0.0223 & 0.853 & 0.927 & 0.218 \\
\meth{richardson\_2} & core & 0.0152 & 0.949 & 0.523 & 0.302 & 0.0193 & 0.855 & 0.831 & 0.241 & 0.0424 & 0.741 & 0.963 & 0.540 \\
 & held-out & 0.0250 & 0.852 & 0.486 & 0.265 & 0.0241 & 0.804 & 0.739 & 0.261 & 0.0196 & 0.904 & 0.966 & 0.227 \\
\meth{double\_exp\_fit} & core & 0.0135 & 0.896 & 0.537 & 0.261 & 0.0278 & 0.926 & 0.712 & 0.447 & 0.0688 & 0.833 & 0.969 & 0.673 \\
 & held-out & 0.0169 & 0.984 & 0.421 & 0.421 & 0.0306 & 0.971 & 0.681 & 0.612 & 0.0797 & 0.870 & 1.000 & 0.733 \\
\meth{pade\_34} & core & 0.0135 & 0.880 & 0.593 & 0.235 & 0.0272 & 0.755 & 0.842 & 0.344 & 0.1112 & 0.417 & 0.868 & 1.363 \\
 & held-out & 0.0108 & 0.945 & 0.727 & 0.180 & 0.0397 & 0.659 & 0.768 & 0.683 & 0.1451 & 0.316 & 0.853 & 2.014 \\
\meth{rational\_fit} & core & 0.0203 & 0.910 & 0.403 & 0.271 & 0.0161 & 0.887 & 0.868 & 0.239 & 0.0329 & 0.735 & 0.993 & 0.290 \\
 & held-out & 0.0061 & 0.852 & 0.607 & 0.250 & 0.0158 & 0.877 & 0.703 & 0.291 & 0.0073 & 0.994 & 1.000 & 0.219 \\
\meth{richardson\_1} & core & 0.0132 & 0.917 & 0.525 & 0.320 & 0.0221 & 0.905 & 0.805 & 0.237 & 0.0420 & 0.735 & 0.993 & 0.594 \\
 & held-out & 0.0250 & 0.852 & 0.492 & 0.197 & 0.0180 & 0.717 & 0.717 & 0.236 & 0.0194 & 0.955 & 1.000 & 0.196 \\
\meth{pade\_23} & core & 0.0124 & 0.887 & 0.632 & 0.216 & 0.0264 & 0.745 & 0.866 & 0.337 & 0.1318 & 0.386 & 0.890 & 1.445 \\
 & held-out & 0.0119 & 0.973 & 0.705 & 0.184 & 0.0414 & 0.674 & 0.783 & 0.704 & 0.1434 & 0.299 & 0.859 & 1.965 \\
\meth{pade\_12} & core & 0.0140 & 0.894 & 0.593 & 0.208 & 0.0264 & 0.736 & 0.879 & 0.358 & 0.1292 & 0.368 & 0.925 & 1.440 \\
 & held-out & 0.0138 & 0.978 & 0.656 & 0.202 & 0.0470 & 0.754 & 0.717 & 0.695 & 0.1331 & 0.305 & 0.893 & 1.674 \\
\meth{pade\_45} & core & 0.0134 & 0.887 & 0.641 & 0.239 & 0.0280 & 0.755 & 0.853 & 0.384 & 0.1140 & 0.397 & 0.908 & 1.383 \\
 & held-out & 0.0108 & 0.945 & 0.716 & 0.184 & 0.0409 & 0.638 & 0.746 & 0.700 & 0.1221 & 0.328 & 0.847 & 1.918 \\
\addlinespace[2pt]
21 classical variants (pooled) & core & 0.0519 & 0.496 & 0.218 & 1.000 & 0.0559 & 0.471 & 0.439 & 1.003 & 0.0531 & 0.485 & 0.853 & 1.001 \\
 & held-out & 0.0531 & 0.552 & 0.123 & 0.996 & 0.0529 & 0.566 & 0.429 & 0.996 & 0.0657 & 0.559 & 0.883 & 0.998 \\
\bottomrule
\end{tabular}

\end{adjustbox}
\end{table}
\end{landscape}
\begin{landscape}
\begin{table}[p]
\centering
\caption{\texttt{f13\_by\_Ltrue\_g0.1.tex}: Phase-1 results by L\textasciicircum{}* tercile at g = 0.1: the deployable trivials, the oracle reference, the top-10 accelerators (by pooled median error at g = 0.1, core) and the 21 classical variants pooled; core | out-of-design}
\label{tab:f13_by_Ltrue_g0.1}
\begin{adjustbox}{max width=\linewidth, max totalheight=0.86\textheight}
% AUTO-GENERATED -- do not hand-edit.  generated by scripts/analyze_by_ltrue.py from the run started 2026-10-01T22:16:39 on code d82fd26 (full; 29,256 s; 7 workers)
% Phase-1 results by L^* tercile at g = 0.1: the deployable trivials, the oracle reference, the top-10 accelerators (by pooled median error at g = 0.1, core) and the 21 classical variants pooled; core | out-of-design
% source : results/phase1/phase1_by_Ltrue.csv (the per-tercile aggregate this fragment prints; the aggregate of the git-ignored results/phase1/phase1_records.csv, capped == 0)
% filter : target_g == 0.1; log-scale terciles of L^* (L_true) in [0.005, 0.5]: T1 [0.005, 0.0232), T2 [0.0232, 0.1077), T3 [0.1077, 0.5]; med. err = median over valid records; win vs X = mean of the per-record win indicator (invalid never wins); skill vs last = median over valid records of err / err(last_value)
\begin{tabular}{llrrrrrrrrrrrr}
\toprule
$g$ / row & & \multicolumn{4}{c}{T1: $L^{*} \in [0.005, 0.0232)$} & \multicolumn{4}{c}{T2: $[0.0232, 0.1077)$} & \multicolumn{4}{c}{T3: $[0.1077, 0.5]$} \\
\cmidrule(lr){3-6}\cmidrule(lr){7-10}\cmidrule(lr){11-14}
 & set & med.\ err & win/last & win/$\hat L$ & skill/last & med.\ err & win/last & win/$\hat L$ & skill/last & med.\ err & win/last & win/$\hat L$ & skill/last \\
\midrule
\multicolumn{14}{l}{\textit{$g = 0.1$} (headline)} \\
predict $\hat L$ (\meth{constant\_assumed}) & core & 0.0214 & 0.867 & -- & 0.252 & 0.0621 & 0.654 & -- & 0.661 & 0.2406 & 0.070 & -- & 2.866 \\
 & out-of-design & 0.0213 & 0.928 & -- & 0.275 & 0.0615 & 0.687 & -- & 0.742 & 0.2175 & 0.072 & -- & 3.126 \\
\meth{last\_value} & core & 0.0904 & -- & 0.133 & -- & 0.1157 & -- & 0.346 & -- & 0.1133 & -- & 0.930 & -- \\
 & out-of-design & 0.0850 & -- & 0.072 & -- & 0.0845 & -- & 0.313 & -- & 0.0862 & -- & 0.928 & -- \\
\meth{window\_mean} & core & 0.1256 & 0.166 & 0.116 & 1.427 & 0.1670 & 0.140 & 0.241 & 1.441 & 0.1663 & 0.168 & 0.771 & 1.431 \\
 & out-of-design & 0.1943 & 0.000 & 0.000 & 1.443 & 0.1945 & 0.000 & 0.143 & 1.463 & 0.1823 & 0.000 & 0.738 & 1.428 \\
\meth{window\_min} & core & 0.0494 & 0.385 & 0.155 & 1.000 & 0.0891 & 0.325 & 0.401 & 1.000 & 0.0822 & 0.335 & 0.935 & 1.000 \\
 & out-of-design & 0.0850 & 0.338 & 0.113 & 1.000 & 0.0845 & 0.367 & 0.320 & 1.000 & 0.0862 & 0.364 & 0.944 & 1.000 \\
\meth{constant\_oracle} \textit{(ref.)} & core & 0.0056 & 1.000 & 1.000 & 0.111 & 0.0077 & 1.000 & 1.000 & 0.111 & 0.0099 & 1.000 & 1.000 & 0.111 \\
 & out-of-design & 0.0092 & 1.000 & 1.000 & 0.110 & 0.0092 & 1.000 & 1.000 & 0.110 & 0.0092 & 1.000 & 1.000 & 0.111 \\
\addlinespace[2pt]
\meth{log\_linear} & core & 0.0101 & 0.843 & 0.617 & 0.246 & 0.0224 & 0.805 & 0.868 & 0.289 & 0.0430 & 0.626 & 0.933 & 0.550 \\
 & out-of-design & 0.0291 & 1.000 & 0.538 & 0.326 & 0.0294 & 0.857 & 0.878 & 0.453 & 0.0364 & 0.631 & 1.000 & 0.380 \\
\meth{richardson\_3} & core & 0.0176 & 0.830 & 0.516 & 0.352 & 0.0225 & 0.815 & 0.813 & 0.302 & 0.0234 & 0.681 & 0.894 & 0.386 \\
 & out-of-design & 0.0255 & 0.841 & 0.549 & 0.222 & 0.0186 & 0.721 & 0.748 & 0.218 & 0.0190 & 0.862 & 0.933 & 0.213 \\
\meth{richardson\_2} & core & 0.0211 & 0.843 & 0.512 & 0.367 & 0.0220 & 0.848 & 0.809 & 0.303 & 0.0267 & 0.730 & 0.904 & 0.478 \\
 & out-of-design & 0.0312 & 0.862 & 0.508 & 0.263 & 0.0213 & 0.714 & 0.741 & 0.251 & 0.0188 & 0.913 & 0.969 & 0.212 \\
\meth{double\_exp\_fit} & core & 0.0201 & 0.910 & 0.527 & 0.349 & 0.0248 & 0.949 & 0.780 & 0.416 & 0.0596 & 0.875 & 0.988 & 0.622 \\
 & out-of-design & 0.0212 & 0.985 & 0.462 & 0.387 & 0.0271 & 0.959 & 0.837 & 0.590 & 0.0539 & 0.944 & 1.000 & 0.707 \\
\meth{pade\_34} & core & 0.0167 & 0.875 & 0.630 & 0.276 & 0.0259 & 0.786 & 0.815 & 0.355 & 0.0915 & 0.405 & 0.820 & 1.292 \\
 & out-of-design & 0.0111 & 0.964 & 0.831 & 0.165 & 0.0289 & 0.769 & 0.816 & 0.557 & 0.1169 & 0.395 & 0.887 & 1.437 \\
\meth{rational\_fit} & core & 0.0286 & 0.895 & 0.477 & 0.316 & 0.0188 & 0.903 & 0.870 & 0.246 & 0.0329 & 0.834 & 1.000 & 0.316 \\
 & out-of-design & 0.0100 & 0.862 & 0.641 & 0.214 & 0.0119 & 0.741 & 0.735 & 0.263 & 0.0079 & 0.995 & 1.000 & 0.208 \\
\meth{richardson\_1} & core & 0.0272 & 0.824 & 0.488 & 0.385 & 0.0273 & 0.862 & 0.794 & 0.297 & 0.0331 & 0.748 & 0.933 & 0.404 \\
 & out-of-design & 0.0312 & 0.862 & 0.518 & 0.251 & 0.0190 & 0.741 & 0.748 & 0.223 & 0.0156 & 0.872 & 1.000 & 0.178 \\
\meth{pade\_23} & core & 0.0157 & 0.897 & 0.619 & 0.256 & 0.0252 & 0.788 & 0.833 & 0.316 & 0.1195 & 0.427 & 0.826 & 1.307 \\
 & out-of-design & 0.0113 & 0.964 & 0.810 & 0.165 & 0.0349 & 0.789 & 0.830 & 0.432 & 0.1224 & 0.374 & 0.872 & 1.409 \\
\meth{pade\_12} & core & 0.0168 & 0.884 & 0.609 & 0.255 & 0.0251 & 0.778 & 0.835 & 0.363 & 0.1105 & 0.401 & 0.857 & 1.366 \\
 & out-of-design & 0.0133 & 0.974 & 0.749 & 0.180 & 0.0389 & 0.782 & 0.844 & 0.542 & 0.1058 & 0.369 & 0.897 & 1.336 \\
\meth{pade\_45} & core & 0.0155 & 0.884 & 0.624 & 0.253 & 0.0273 & 0.798 & 0.823 & 0.391 & 0.0835 & 0.427 & 0.845 & 1.302 \\
 & out-of-design & 0.0116 & 0.949 & 0.800 & 0.170 & 0.0369 & 0.707 & 0.776 & 0.570 & 0.1030 & 0.431 & 0.882 & 1.250 \\
\addlinespace[2pt]
21 classical variants (pooled) & core & 0.0499 & 0.499 & 0.236 & 1.000 & 0.0554 & 0.476 & 0.452 & 1.002 & 0.0508 & 0.491 & 0.900 & 1.000 \\
 & out-of-design & 0.0481 & 0.554 & 0.206 & 0.997 & 0.0483 & 0.573 & 0.472 & 0.994 & 0.0552 & 0.566 & 0.929 & 0.997 \\
\bottomrule
\end{tabular}

\end{adjustbox}
\end{table}
\end{landscape}
\begin{landscape}
\begin{table}[p]
\centering
\caption{\texttt{f13\_by\_Ltrue\_g0.5.tex}: Phase-1 results by L\textasciicircum{}* tercile at g = 0.5: the deployable trivials, the oracle reference, the top-10 accelerators (by pooled median error at g = 0.1, core) and the 21 classical variants pooled; core | out-of-design}
\label{tab:f13_by_Ltrue_g0.5}
\begin{adjustbox}{max width=\linewidth, max totalheight=0.86\textheight}
% AUTO-GENERATED -- do not hand-edit.  generated by scripts/analyze_by_ltrue.py from the run started 2026-10-01T22:16:39 on code d82fd26 (full; 29,256 s; 7 workers)
% Phase-1 results by L^* tercile at g = 0.5: the deployable trivials, the oracle reference, the top-10 accelerators (by pooled median error at g = 0.1, core) and the 21 classical variants pooled; core | out-of-design
% source : results/phase1/phase1_by_Ltrue.csv (the per-tercile aggregate this fragment prints; the aggregate of the git-ignored results/phase1/phase1_records.csv, capped == 0)
% filter : target_g == 0.5; log-scale terciles of L^* (L_true) in [0.005, 0.5]: T1 [0.005, 0.0232), T2 [0.0232, 0.1077), T3 [0.1077, 0.5]; med. err = median over valid records; win vs X = mean of the per-record win indicator (invalid never wins); skill vs last = median over valid records of err / err(last_value)
\begin{tabular}{llrrrrrrrrrrrr}
\toprule
$g$ / row & & \multicolumn{4}{c}{T1: $L^{*} \in [0.005, 0.0232)$} & \multicolumn{4}{c}{T2: $[0.0232, 0.1077)$} & \multicolumn{4}{c}{T3: $[0.1077, 0.5]$} \\
\cmidrule(lr){3-6}\cmidrule(lr){7-10}\cmidrule(lr){11-14}
 & set & med.\ err & win/last & win/$\hat L$ & skill/last & med.\ err & win/last & win/$\hat L$ & skill/last & med.\ err & win/last & win/$\hat L$ & skill/last \\
\midrule
\multicolumn{14}{l}{\textit{$g = 0.5$}} \\
predict $\hat L$ (\meth{constant\_assumed}) & core & 0.0713 & 0.054 & -- & 1.210 & 0.1133 & 0.064 & -- & 1.836 & 0.2963 & 0.000 & -- & 5.464 \\
 & out-of-design & 0.0524 & 0.000 & -- & 1.282 & 0.0897 & 0.000 & -- & 2.085 & 0.2453 & 0.000 & -- & 6.459 \\
\meth{last\_value} & core & 0.0658 & -- & 0.946 & -- & 0.0738 & -- & 0.936 & -- & 0.0665 & -- & 1.000 & -- \\
 & out-of-design & 0.0482 & -- & 1.000 & -- & 0.0481 & -- & 1.000 & -- & 0.0498 & -- & 1.000 & -- \\
\meth{window\_mean} & core & 0.1059 & 0.200 & 0.252 & 1.702 & 0.1075 & 0.203 & 0.599 & 1.709 & 0.1074 & 0.187 & 0.967 & 1.709 \\
 & out-of-design & 0.1225 & 0.000 & 0.212 & 1.807 & 0.1323 & 0.000 & 0.551 & 1.883 & 0.1143 & 0.000 & 0.923 & 1.782 \\
\meth{window\_min} & core & 0.0458 & 0.418 & 0.906 & 1.000 & 0.0499 & 0.387 & 0.934 & 1.000 & 0.0499 & 0.369 & 1.000 & 1.000 \\
 & out-of-design & 0.0475 & 0.333 & 1.000 & 1.000 & 0.0481 & 0.367 & 1.000 & 1.000 & 0.0498 & 0.364 & 1.000 & 1.000 \\
\meth{constant\_oracle} \textit{(ref.)} & core & 0.0658 & 0.721 & 1.000 & 0.994 & 0.0500 & 0.740 & 1.000 & 0.992 & 0.0579 & 0.722 & 1.000 & 0.993 \\
 & out-of-design & 0.0455 & 0.667 & 1.000 & 0.996 & 0.0455 & 0.748 & 1.000 & 0.993 & 0.0455 & 0.692 & 1.000 & 0.995 \\
\addlinespace[2pt]
\meth{log\_linear} & core & 0.0127 & 0.829 & 0.842 & 0.335 & 0.0144 & 0.873 & 0.918 & 0.296 & 0.0153 & 0.770 & 0.950 & 0.347 \\
 & out-of-design & 0.0197 & 0.859 & 0.924 & 0.276 & 0.0142 & 0.721 & 1.000 & 0.425 & 0.0109 & 0.764 & 1.000 & 0.719 \\
\meth{richardson\_3} & core & 0.0200 & 0.751 & 0.836 & 0.379 & 0.0158 & 0.854 & 0.913 & 0.344 & 0.0093 & 0.769 & 0.920 & 0.263 \\
 & out-of-design & 0.0160 & 0.843 & 0.859 & 0.223 & 0.0085 & 0.714 & 0.844 & 0.157 & 0.0053 & 0.774 & 0.938 & 0.139 \\
\meth{richardson\_2} & core & 0.0159 & 0.757 & 0.829 & 0.391 & 0.0172 & 0.838 & 0.914 & 0.322 & 0.0105 & 0.781 & 0.930 & 0.243 \\
 & out-of-design & 0.0160 & 0.859 & 0.864 & 0.278 & 0.0094 & 0.728 & 0.816 & 0.218 & 0.0048 & 0.805 & 0.969 & 0.128 \\
\meth{double\_exp\_fit} & core & 0.0120 & 0.871 & 0.948 & 0.308 & 0.0152 & 0.902 & 0.995 & 0.308 & 0.0186 & 0.861 & 1.000 & 0.420 \\
 & out-of-design & 0.0107 & 0.975 & 1.000 & 0.407 & 0.0154 & 0.952 & 1.000 & 0.562 & 0.0157 & 0.974 & 1.000 & 0.456 \\
\meth{pade\_34} & core & 0.0130 & 0.701 & 0.751 & 0.526 & 0.0242 & 0.654 & 0.813 & 0.581 & 0.0540 & 0.457 & 0.819 & 1.163 \\
 & out-of-design & 0.0047 & 0.919 & 0.990 & 0.108 & 0.0181 & 0.769 & 0.986 & 0.384 & 0.0426 & 0.538 & 0.959 & 0.865 \\
\meth{rational\_fit} & core & 0.0199 & 0.827 & 0.873 & 0.310 & 0.0167 & 0.877 & 0.979 & 0.355 & 0.0089 & 0.841 & 1.000 & 0.264 \\
 & out-of-design & 0.0058 & 0.854 & 0.864 & 0.178 & 0.0071 & 0.728 & 0.816 & 0.218 & 0.0067 & 0.882 & 1.000 & 0.205 \\
\meth{richardson\_1} & core & 0.0241 & 0.726 & 0.803 & 0.391 & 0.0204 & 0.825 & 0.900 & 0.322 & 0.0124 & 0.794 & 0.967 & 0.238 \\
 & out-of-design & 0.0212 & 0.859 & 0.864 & 0.278 & 0.0130 & 0.735 & 0.755 & 0.216 & 0.0052 & 0.836 & 1.000 & 0.152 \\
\meth{pade\_23} & core & 0.0131 & 0.697 & 0.744 & 0.513 & 0.0225 & 0.660 & 0.820 & 0.528 & 0.0473 & 0.456 & 0.831 & 1.090 \\
 & out-of-design & 0.0046 & 0.924 & 0.990 & 0.109 & 0.0185 & 0.803 & 0.980 & 0.430 & 0.0486 & 0.503 & 0.949 & 0.971 \\
\meth{pade\_12} & core & 0.0136 & 0.719 & 0.769 & 0.493 & 0.0188 & 0.660 & 0.818 & 0.537 & 0.0508 & 0.457 & 0.867 & 1.090 \\
 & out-of-design & 0.0059 & 0.934 & 1.000 & 0.137 & 0.0149 & 0.782 & 0.980 & 0.354 & 0.0410 & 0.513 & 0.944 & 0.932 \\
\meth{pade\_45} & core & 0.0137 & 0.705 & 0.782 & 0.552 & 0.0251 & 0.640 & 0.840 & 0.752 & 0.0477 & 0.476 & 0.894 & 1.112 \\
 & out-of-design & 0.0055 & 0.899 & 0.980 & 0.128 & 0.0168 & 0.748 & 0.980 & 0.350 & 0.0381 & 0.559 & 0.923 & 0.697 \\
\addlinespace[2pt]
21 classical variants (pooled) & core & 0.0517 & 0.471 & 0.854 & 1.005 & 0.0638 & 0.455 & 0.890 & 1.007 & 0.0539 & 0.474 & 0.970 & 1.004 \\
 & out-of-design & 0.0356 & 0.511 & 0.856 & 1.000 & 0.0371 & 0.523 & 0.952 & 0.999 & 0.0427 & 0.523 & 0.984 & 0.999 \\
\bottomrule
\end{tabular}

\end{adjustbox}
\end{table}
\end{landscape}
\begin{landscape}
\begin{table}[p]
\centering
\caption{\texttt{f14\_real\_fixed\_methods.tex}: Real data: fixed rational\_fit and richardson\_1 next to the cascade, per dataset and pooled, pre- and post-minimum targets: n, win vs last, fail, valid, median error and skill}
\label{tab:f14_real_fixed_methods}
\begin{adjustbox}{max width=\linewidth, max totalheight=0.86\textheight}
% AUTO-GENERATED -- do not hand-edit.  generated by scripts/make_paper_tables.py from the run started 2026-10-01T22:16:39 on code d82fd26 (full; 29,256 s; 7 workers)
% Real data (Table tab:realmethods): fixed rational_fit and richardson_1 next to the cascade, per dataset and pooled, pre- and post-minimum targets: n, win vs last, fail, valid, median error and skill
% source : results/real_data/real_data_results_v2.csv (recorded curves re-evaluated on the (depth x target) grid; L_hat mode zero)
% filter : method in {rational_fit, richardson_1} and the cascade row (is_cascade == 1); split = pre-minimum (target_round <= argmin_round) / post-minimum / total; win vs last = rows with win_vs_last == 1 (error below the last observed value's); fail = skill >= 1 or error not finite (no better than the hindsight best trivial; the definition of real_data_strata_v2.csv); valid = rows with a finite error; med. err and med. skill over the valid rows only (conditional on validity)
\begin{tabular}{lllrrrrrr}
\toprule
Dataset & split & method & $n$ & win vs last & fail & valid & med.\ err & med.\ skill \\
\midrule
\textbf{pooled} & pre-minimum & \meth{rational\_fit} & 65 & 65/65 & 0/65 & 65/65 & 0.0106 & 0.270 \\
 &  & \meth{richardson\_1} & 65 & 63/65 & 2/65 & 65/65 & 0.0119 & 0.268 \\
 &  & cascade & 65 & 63/65 & 2/65 & 65/65 & 0.0101 & 0.268 \\
\addlinespace[1pt]
 & post-minimum & \meth{rational\_fit} & 25 & 13/25 & 12/25 & 25/25 & 0.0079 & 0.783 \\
 &  & \meth{richardson\_1} & 25 & 10/25 & 15/25 & 25/25 & 0.0113 & 1.312 \\
 &  & cascade & 25 & 10/25 & 15/25 & 25/25 & 0.0106 & 1.345 \\
\addlinespace[1pt]
 & total & \meth{rational\_fit} & 90 & 78/90 & 12/90 & 90/90 & 0.0090 & 0.355 \\
 &  & \meth{richardson\_1} & 90 & 73/90 & 17/90 & 90/90 & 0.0116 & 0.345 \\
 &  & cascade & 90 & 73/90 & 17/90 & 90/90 & 0.0103 & 0.336 \\
\midrule
\texttt{adult} & pre-minimum & \meth{rational\_fit} & 5 & 5/5 & 0/5 & 5/5 & 0.0055 & 0.721 \\
 &  & \meth{richardson\_1} & 5 & 4/5 & 1/5 & 5/5 & 0.0064 & 0.532 \\
 &  & cascade & 5 & 4/5 & 1/5 & 5/5 & 0.0055 & 0.721 \\
\addlinespace[1pt]
 & post-minimum & \meth{rational\_fit} & 10 & 4/10 & 6/10 & 10/10 & 0.0094 & 1.077 \\
 &  & \meth{richardson\_1} & 10 & 4/10 & 6/10 & 10/10 & 0.0110 & 1.159 \\
 &  & cascade & 10 & 4/10 & 6/10 & 10/10 & 0.0094 & 1.232 \\
\addlinespace[1pt]
 & total & \meth{rational\_fit} & 15 & 9/15 & 6/15 & 15/15 & 0.0082 & 0.755 \\
 &  & \meth{richardson\_1} & 15 & 8/15 & 7/15 & 15/15 & 0.0096 & 0.841 \\
 &  & cascade & 15 & 8/15 & 7/15 & 15/15 & 0.0082 & 0.755 \\
\addlinespace[3pt]
\texttt{bank\_marketing} & pre-minimum & \meth{rational\_fit} & 0 & 0/0 & 0/0 & 0/0 & -- & -- \\
 &  & \meth{richardson\_1} & 0 & 0/0 & 0/0 & 0/0 & -- & -- \\
 &  & cascade & 0 & 0/0 & 0/0 & 0/0 & -- & -- \\
\addlinespace[1pt]
 & post-minimum & \meth{rational\_fit} & 15 & 9/15 & 6/15 & 15/15 & 0.0072 & 0.758 \\
 &  & \meth{richardson\_1} & 15 & 6/15 & 9/15 & 15/15 & 0.0113 & 1.389 \\
 &  & cascade & 15 & 6/15 & 9/15 & 15/15 & 0.0113 & 1.595 \\
\addlinespace[1pt]
 & total & \meth{rational\_fit} & 15 & 9/15 & 6/15 & 15/15 & 0.0072 & 0.758 \\
 &  & \meth{richardson\_1} & 15 & 6/15 & 9/15 & 15/15 & 0.0113 & 1.389 \\
 &  & cascade & 15 & 6/15 & 9/15 & 15/15 & 0.0113 & 1.595 \\
\addlinespace[3pt]
\texttt{covertype} & pre-minimum & \meth{rational\_fit} & 15 & 15/15 & 0/15 & 15/15 & 0.0436 & 0.270 \\
 &  & \meth{richardson\_1} & 15 & 15/15 & 0/15 & 15/15 & 0.0407 & 0.252 \\
 &  & cascade & 15 & 15/15 & 0/15 & 15/15 & 0.0407 & 0.252 \\
\addlinespace[1pt]
 & post-minimum & \meth{rational\_fit} & 0 & 0/0 & 0/0 & 0/0 & -- & -- \\
 &  & \meth{richardson\_1} & 0 & 0/0 & 0/0 & 0/0 & -- & -- \\
 &  & cascade & 0 & 0/0 & 0/0 & 0/0 & -- & -- \\
\addlinespace[1pt]
 & total & \meth{rational\_fit} & 15 & 15/15 & 0/15 & 15/15 & 0.0436 & 0.270 \\
 &  & \meth{richardson\_1} & 15 & 15/15 & 0/15 & 15/15 & 0.0407 & 0.252 \\
 &  & cascade & 15 & 15/15 & 0/15 & 15/15 & 0.0407 & 0.252 \\
\addlinespace[3pt]
\texttt{higgs} & pre-minimum & \meth{rational\_fit} & 15 & 15/15 & 0/15 & 15/15 & 0.0040 & 0.199 \\
 &  & \meth{richardson\_1} & 15 & 15/15 & 0/15 & 15/15 & 0.0069 & 0.247 \\
 &  & cascade & 15 & 15/15 & 0/15 & 15/15 & 0.0040 & 0.199 \\
\addlinespace[1pt]
 & post-minimum & \meth{rational\_fit} & 0 & 0/0 & 0/0 & 0/0 & -- & -- \\
 &  & \meth{richardson\_1} & 0 & 0/0 & 0/0 & 0/0 & -- & -- \\
 &  & cascade & 0 & 0/0 & 0/0 & 0/0 & -- & -- \\
\addlinespace[1pt]
 & total & \meth{rational\_fit} & 15 & 15/15 & 0/15 & 15/15 & 0.0040 & 0.199 \\
 &  & \meth{richardson\_1} & 15 & 15/15 & 0/15 & 15/15 & 0.0069 & 0.247 \\
 &  & cascade & 15 & 15/15 & 0/15 & 15/15 & 0.0040 & 0.199 \\
\addlinespace[3pt]
\texttt{jannis} & pre-minimum & \meth{rational\_fit} & 15 & 15/15 & 0/15 & 15/15 & 0.0089 & 0.369 \\
 &  & \meth{richardson\_1} & 15 & 14/15 & 1/15 & 15/15 & 0.0134 & 0.382 \\
 &  & cascade & 15 & 14/15 & 1/15 & 15/15 & 0.0134 & 0.396 \\
\addlinespace[1pt]
 & post-minimum & \meth{rational\_fit} & 0 & 0/0 & 0/0 & 0/0 & -- & -- \\
 &  & \meth{richardson\_1} & 0 & 0/0 & 0/0 & 0/0 & -- & -- \\
 &  & cascade & 0 & 0/0 & 0/0 & 0/0 & -- & -- \\
\addlinespace[1pt]
 & total & \meth{rational\_fit} & 15 & 15/15 & 0/15 & 15/15 & 0.0089 & 0.369 \\
 &  & \meth{richardson\_1} & 15 & 14/15 & 1/15 & 15/15 & 0.0134 & 0.382 \\
 &  & cascade & 15 & 14/15 & 1/15 & 15/15 & 0.0134 & 0.396 \\
\addlinespace[3pt]
\texttt{miniboone} & pre-minimum & \meth{rational\_fit} & 15 & 15/15 & 0/15 & 15/15 & 0.0064 & 0.233 \\
 &  & \meth{richardson\_1} & 15 & 15/15 & 0/15 & 15/15 & 0.0071 & 0.130 \\
 &  & cascade & 15 & 15/15 & 0/15 & 15/15 & 0.0071 & 0.130 \\
\addlinespace[1pt]
 & post-minimum & \meth{rational\_fit} & 0 & 0/0 & 0/0 & 0/0 & -- & -- \\
 &  & \meth{richardson\_1} & 0 & 0/0 & 0/0 & 0/0 & -- & -- \\
 &  & cascade & 0 & 0/0 & 0/0 & 0/0 & -- & -- \\
\addlinespace[1pt]
 & total & \meth{rational\_fit} & 15 & 15/15 & 0/15 & 15/15 & 0.0064 & 0.233 \\
 &  & \meth{richardson\_1} & 15 & 15/15 & 0/15 & 15/15 & 0.0071 & 0.130 \\
 &  & cascade & 15 & 15/15 & 0/15 & 15/15 & 0.0071 & 0.130 \\
\addlinespace[3pt]
\bottomrule
\end{tabular}

\end{adjustbox}
\end{table}
\end{landscape}
\begin{landscape}
\begin{table}[p]
\centering
\caption{\texttt{f15\_bootstrap.tex}: Family bootstrap over the core curve families: rational\_fit's win rate vs the last value and its cell-median error against each comparator, point and 2.5 / 97.5 percentiles}
\label{tab:f15_bootstrap}
\begin{adjustbox}{max width=\linewidth, max totalheight=0.86\textheight}
% AUTO-GENERATED -- do not hand-edit.  generated by scripts/make_paper_tables.py from the run started 2026-10-01T22:16:39 on code d82fd26 (full; 29,256 s; 7 workers)
% Family bootstrap over core regimes (Table tab:boot): rational_fit's win rate vs the last value and its cell-median error against each comparator, point and 2.5 / 97.5 percentiles
% source : results/phase1/phase1_aggregated.csv (core regimes, uncapped cells at the headline stratum)
% filter : is_holdout == 0, capped == 0, target_g == 0.1; resampling unit = core regime (16 regimes, 48 cells), 5000 draws with replacement from RandomState(20260927); point = unresampled cells; interval = 2.5th / 97.5th percentile of the draws; lower-error fraction and delta median over the cells where both methods have a finite cell-median error (excl. = cells dropped because either method has none; cells + excl. = the cell total); positive delta = rational_fit worse
\begin{tabular}{llrrrrr}
\toprule
Statistic & comparator & point & 2.5\% & 97.5\% & cells & excl. \\
\midrule
\meth{rational\_fit} win rate vs last & -- & 0.877 & 0.751 & 0.980 & 48 & -- \\
\addlinespace[2pt]
lower-error fraction & \meth{log\_linear} & 0.604 & 0.375 & 0.833 & 48 & 0 \\
$\Delta$ median error & \meth{log\_linear} & $+0.00938$ & $-0.01043$ & $+0.02002$ & 48 & 0 \\
lower-error fraction & \meth{richardson\_1} & 0.542 & 0.312 & 0.771 & 48 & 0 \\
$\Delta$ median error & \meth{richardson\_1} & $-0.00050$ & $-0.02005$ & $+0.00730$ & 48 & 0 \\
lower-error fraction & \meth{richardson\_2} & 0.562 & 0.333 & 0.792 & 48 & 0 \\
$\Delta$ median error & \meth{richardson\_2} & $+0.00388$ & $-0.02116$ & $+0.01091$ & 48 & 0 \\
lower-error fraction & \meth{richardson\_3} & 0.562 & 0.333 & 0.792 & 48 & 0 \\
$\Delta$ median error & \meth{richardson\_3} & $+0.00737$ & $-0.01716$ & $+0.01556$ & 48 & 0 \\
lower-error fraction & \meth{double\_exp\_fit} & 0.417 & 0.208 & 0.646 & 48 & 0 \\
$\Delta$ median error & \meth{double\_exp\_fit} & $+0.00083$ & $-0.04059$ & $+0.02309$ & 48 & 0 \\
\bottomrule
\end{tabular}

\end{adjustbox}
\end{table}
\end{landscape}
\begin{landscape}
\begin{table}[p]
\centering
\caption{\texttt{f16\_generalisation\_summary.tex}: Generalisation summary: eligible accelerators per set of curve families, rank correlation core vs out-of-design, and the two headline fits}
\label{tab:f16_generalisation_summary}
\begin{adjustbox}{max width=\linewidth, max totalheight=0.86\textheight}
% AUTO-GENERATED -- do not hand-edit.  generated by scripts/make_paper_tables.py from the run started 2026-10-01T22:16:39 on code d82fd26 (full; 29,256 s; 7 workers)
% Generalisation summary: eligible accelerators per set of curve families, rank correlation core vs out-of-design, and the two headline fits
% source : results/phase1/phase1_global.csv (core regimes, all three strata, capped cells excluded)
% source : results/phase1/phase1_global_holdout.csv (out-of-design regimes)
% filter : per stratum; eligible = accelerators with valid_rate >= 0.9 (rank_eligible == 1) in that regime set; rho = Spearman correlation of the core and out-of-design med_error ranks over the accelerators eligible on both (the generalisation FACTS rows); ranks are among the eligible accelerators of that regime set
\begin{tabular}{lrrrrll}
\toprule
$g$ & eligible core & eligible out-of-design & eligible both & Spearman $\rho$ & \meth{rational\_fit} core / out-of-design & \meth{log\_linear} core / out-of-design \\
\midrule
0.5 & 47/52 & 48/52 & 47 & 0.877 & 8 / 1 & 1 / 8 \\
0.1 & 44/52 & 45/52 & 44 & 0.833 & 6 / 1 & 1 / 10 \\
0.02 & 44/52 & 43/52 & 43 & 0.873 & 3 / 1 & 4 / 12 \\
\bottomrule
\end{tabular}

\end{adjustbox}
\end{table}
\end{landscape}
\begin{landscape}
\begin{table}[p]
\centering
\caption{\texttt{f18\_ladders\_classical.tex}: Order ladders of the classical families (Phase 0b): every order next to the roster orders, core and out-of-design}
\label{tab:f18_ladders_classical}
\begin{adjustbox}{max width=\linewidth, max totalheight=0.86\textheight}
% AUTO-GENERATED -- do not hand-edit.  generated by scripts/make_paper_tables.py from the run started 2026-10-01T22:16:39 on code d82fd26 (full; 29,256 s; 7 workers)
% Order ladders of the classical families (Phase 0b): every order next to the roster orders, core and held-out
% source : results/phase0b/order_ladders_panels.csv (Phase 0b: every order of every family with an order parameter, evaluated on Phase 1's grid at the headline stratum; noise pooled; uncapped cells; the roster orders reproduce phase1_records.csv exactly, results/phase0b/order_ladders_agreement.txt)
% filter : regime_set core (48 cells) | held-out (18 cells), noise == pooled; bullet = roster method; rho_V = valid rate over all records; med. err, q25--q75 over the valid records (conditional on validity); win vs last = fraction of all records where the variant is valid and below the last-value error
\begin{tabular}{llrrrrrrr}
\toprule
Family & order & \multicolumn{4}{c}{core} & \multicolumn{3}{c}{held-out} \\
\cmidrule(lr){3-6}\cmidrule(lr){7-9}
 & & $\rhoV$ & med.\ err & q25--q75 & win vs last & $\rhoV$ & med.\ err & win vs last \\
\midrule
shanks & k=1 \textbullet & 1.000 & 0.0674 & 0.0221--0.1413 & 0.613 & 1.000 & 0.0476 & 0.713 \\
 & k=2 \textbullet & 0.999 & 0.0478 & 0.0208--0.1457 & 0.353 & 1.000 & 0.0481 & 0.387 \\
 & k=3 \textbullet & 1.000 & 0.0475 & 0.0150--0.1405 & 0.401 & 0.998 & 0.0420 & 0.510 \\
 & k=4 \textbullet & 0.979 & 0.0475 & 0.0146--0.1402 & 0.370 & 1.000 & 0.0491 & 0.458 \\
 & k=5 & 0.979 & 0.0475 & 0.0146--0.1405 & 0.374 & 1.000 & 0.0486 & 0.464 \\
 & k=6 & 0.979 & 0.0475 & 0.0146--0.1422 & 0.348 & 1.000 & 0.0492 & 0.426 \\
 & k=7 & 0.979 & 0.0475 & 0.0146--0.1435 & 0.342 & 0.998 & 0.0492 & 0.408 \\
 & k=8 & 0.979 & 0.0475 & 0.0146--0.1450 & 0.327 & 1.000 & 0.0491 & 0.399 \\
 & k=9 & 0.979 & 0.0482 & 0.0145--0.1460 & 0.311 & 1.000 & 0.0523 & 0.382 \\
 & k=10 & 0.979 & 0.0484 & 0.0147--0.1469 & 0.304 & 1.000 & 0.0508 & 0.367 \\
\addlinespace[2pt]
wynn\_eps & k=1 \textbullet & 0.979 & 0.0516 & 0.0209--0.1361 & 0.613 & 1.000 & 0.0476 & 0.713 \\
 & k=2 \textbullet & 0.967 & 0.0473 & 0.0207--0.1361 & 0.471 & 1.000 & 0.0478 & 0.516 \\
 & k=3 \textbullet & 0.957 & 0.0410 & 0.0160--0.1364 & 0.488 & 0.998 & 0.0476 & 0.572 \\
 & k=4 & 0.956 & 0.0411 & 0.0160--0.1380 & 0.470 & 0.968 & 0.0489 & 0.531 \\
 & k=5 & 0.954 & 0.0432 & 0.0160--0.1377 & 0.481 & 0.996 & 0.0467 & 0.562 \\
 & k=6 & 0.953 & 0.0430 & 0.0147--0.1392 & 0.498 & 0.993 & 0.0455 & 0.562 \\
 & k=7 & 0.951 & 0.0447 & 0.0148--0.1382 & 0.503 & 0.998 & 0.0461 & 0.566 \\
 & k=8 & 0.947 & 0.0406 & 0.0152--0.1327 & 0.568 & 0.996 & 0.0446 & 0.620 \\
\addlinespace[2pt]
wynn\_rho & k=1 \textbullet & 0.979 & 0.0500 & 0.0188--0.1362 & 0.438 & 1.000 & 0.0467 & 0.521 \\
 & k=2 \textbullet & 0.978 & 0.0544 & 0.0185--0.1445 & 0.454 & 0.942 & 0.0509 & 0.523 \\
 & k=3 \textbullet & 0.973 & 0.0501 & 0.0174--0.1419 & 0.494 & 0.996 & 0.0526 & 0.486 \\
 & k=4 & 0.973 & 0.0686 & 0.0175--0.1564 & 0.440 & 0.993 & 0.0507 & 0.536 \\
 & k=5 & 0.961 & 0.0683 & 0.0189--0.1633 & 0.450 & 0.987 & 0.0451 & 0.562 \\
 & k=6 & 0.958 & 0.0610 & 0.0223--0.1733 & 0.392 & 0.980 & 0.0435 & 0.512 \\
 & k=7 & 0.933 & 0.0578 & 0.0183--0.1720 & 0.437 & 0.963 & 0.0410 & 0.544 \\
\addlinespace[2pt]
levin\_t & k=1 \textbullet & 0.979 & 0.0516 & 0.0209--0.1361 & 0.613 & 1.000 & 0.0476 & 0.713 \\
 & k=2 \textbullet & 0.979 & 0.0478 & 0.0215--0.1349 & 0.503 & 1.000 & 0.0478 & 0.553 \\
 & k=3 & 0.978 & 0.0487 & 0.0240--0.1386 & 0.423 & 0.998 & 0.0481 & 0.490 \\
 & k=4 & 0.978 & 0.0481 & 0.0125--0.1453 & 0.401 & 1.000 & 0.0479 & 0.451 \\
 & k=5 & 0.977 & 0.0482 & 0.0175--0.1449 & 0.397 & 1.000 & 0.0504 & 0.445 \\
 & k=6 & 0.979 & 0.0482 & 0.0155--0.1506 & 0.387 & 1.000 & 0.0492 & 0.367 \\
 & k=7 & 0.958 & 0.0475 & 0.0186--0.1425 & 0.388 & 1.000 & 0.0492 & 0.352 \\
 & k=8 & 0.973 & 0.0479 & 0.0162--0.1505 & 0.367 & 1.000 & 0.0495 & 0.415 \\
\addlinespace[2pt]
levin\_u & k=1 \textbullet & 0.979 & 0.0863 & 0.0208--0.1542 & 0.503 & 1.000 & 0.0839 & 0.538 \\
 & k=2 \textbullet & 0.979 & 0.0478 & 0.0170--0.1432 & 0.462 & 1.000 & 0.0626 & 0.551 \\
 & k=3 & 0.979 & 0.0487 & 0.0139--0.1422 & 0.442 & 1.000 & 0.0478 & 0.490 \\
 & k=4 & 0.958 & 0.0493 & 0.0132--0.1463 & 0.379 & 1.000 & 0.0478 & 0.447 \\
 & k=5 & 0.977 & 0.0482 & 0.0132--0.1494 & 0.378 & 1.000 & 0.0497 & 0.447 \\
 & k=6 & 0.978 & 0.0517 & 0.0142--0.1474 & 0.404 & 1.000 & 0.0491 & 0.367 \\
 & k=7 & 0.979 & 0.0511 & 0.0146--0.1473 & 0.387 & 1.000 & 0.0490 & 0.352 \\
 & k=8 & 0.958 & 0.0525 & 0.0246--0.1509 & 0.341 & 0.963 & 0.0585 & 0.344 \\
\addlinespace[2pt]
levin\_v & k=1 \textbullet & 0.979 & 0.0480 & 0.0183--0.1457 & 0.390 & 1.000 & 0.0566 & 0.482 \\
 & k=2 \textbullet & 0.979 & 0.0509 & 0.0132--0.1487 & 0.420 & 1.000 & 0.0482 & 0.484 \\
 & k=3 & 0.977 & 0.0513 & 0.0133--0.1499 & 0.353 & 1.000 & 0.0486 & 0.417 \\
 & k=4 & 0.979 & 0.0514 & 0.0131--0.1471 & 0.391 & 1.000 & 0.0487 & 0.391 \\
 & k=5 & 0.956 & 0.0505 & 0.0128--0.1419 & 0.368 & 1.000 & 0.0493 & 0.322 \\
 & k=6 & 0.978 & 0.0494 & 0.0216--0.1505 & 0.362 & 1.000 & 0.0497 & 0.328 \\
 & k=7 & 0.979 & 0.0955 & 0.0256--0.1604 & 0.333 & 0.998 & 0.0850 & 0.302 \\
 & k=8 & 0.976 & 0.1002 & 0.0264--0.1518 & 0.276 & 1.000 & 0.0870 & 0.236 \\
\addlinespace[2pt]
brezinski\_theta & k=1 \textbullet & 1.000 & 0.0664 & 0.0206--0.1390 & 0.592 & 1.000 & 0.0467 & 0.711 \\
 & k=2 \textbullet & 0.976 & 0.0499 & 0.0249--0.1477 & 0.362 & 0.998 & 0.0530 & 0.436 \\
 & k=3 & 0.976 & 0.0507 & 0.0209--0.1357 & 0.521 & 0.993 & 0.0349 & 0.698 \\
 & k=4 & 0.967 & 0.0727 & 0.0274--0.1755 & 0.386 & 0.991 & 0.0543 & 0.467 \\
 & k=5 & 0.953 & 0.0962 & 0.0271--0.2329 & 0.398 & 0.952 & 0.0907 & 0.438 \\
 & k=6 & 0.899 & 0.1451 & 0.0386--0.3291 & 0.338 & 0.933 & 0.1026 & 0.428 \\
\addlinespace[2pt]
anderson & m=1 \textbullet & 1.000 & 0.1063 & 0.0280--0.1507 & 0.613 & 1.000 & 0.0703 & 0.713 \\
 & m=2 \textbullet & 1.000 & 0.1037 & 0.0269--0.1453 & 0.571 & 1.000 & 0.0755 & 0.644 \\
 & m=3 \textbullet & 1.000 & 0.1040 & 0.0271--0.1450 & 0.540 & 1.000 & 0.0795 & 0.609 \\
 & m=4 & 1.000 & 0.1035 & 0.0270--0.1448 & 0.547 & 1.000 & 0.0785 & 0.615 \\
 & m=5 & 1.000 & 0.1032 & 0.0269--0.1448 & 0.547 & 1.000 & 0.0787 & 0.602 \\
 & m=6 & 1.000 & 0.1026 & 0.0270--0.1448 & 0.542 & 1.000 & 0.0784 & 0.600 \\
\addlinespace[2pt]
neville & d=1 & 0.915 & 0.0707 & 0.0316--0.1782 & 0.462 & 0.972 & 0.0614 & 0.566 \\
 & d=2 \textbullet & 0.661 & 1.3841 & 0.0209--9.1455 & 0.231 & 0.726 & 0.7464 & 0.335 \\
 & d=3 \textbullet & 0.480 & 0.1798 & 0.0103--88.2801 & 0.229 & 0.538 & 0.0280 & 0.333 \\
 & d=4 \textbullet & 0.335 & 0.0209 & 0.0047--0.3368 & 0.229 & 0.374 & 0.0070 & 0.278 \\
 & d=5 & 0.276 & 0.0124 & 0.0018--0.0412 & 0.229 & 0.339 & 0.0051 & 0.278 \\
 & d=6 & 0.260 & 0.0083 & 0.0015--0.0345 & 0.229 & 0.222 & 0.0017 & 0.222 \\
 & d=7 & 0.250 & 0.0028 & 0.0010--0.0292 & 0.229 & 0.276 & 0.0042 & 0.220 \\
 & d=8 & 0.230 & 0.0027 & 0.0004--0.0262 & 0.208 & 0.220 & 0.0008 & 0.220 \\
\addlinespace[2pt]
pade & [1,1] \textbullet & 0.999 & 0.0377 & 0.0112--0.0791 & 0.700 & 0.996 & 0.0331 & 0.700 \\
 & [1,2] \textbullet & 0.995 & 0.0284 & 0.0100--0.0772 & 0.684 & 0.993 & 0.0389 & 0.702 \\
 & [1,3] \textbullet & 0.972 & 0.0379 & 0.0109--0.1252 & 0.575 & 0.989 & 0.0692 & 0.538 \\
 & [1,4] & 0.986 & 0.0424 & 0.0190--0.1445 & 0.549 & 0.991 & 0.0425 & 0.603 \\
 & [1,5] & 0.992 & 0.0507 & 0.0242--0.1502 & 0.548 & 1.000 & 0.0507 & 0.564 \\
 & [2,1] \textbullet & 0.602 & 0.1334 & 0.0071--0.6466 & 0.154 & 0.629 & 0.1603 & 0.201 \\
 & [2,2] \textbullet & 0.924 & 0.0325 & 0.0029--0.0706 & 0.657 & 0.948 & 0.0196 & 0.838 \\
 & [2,3] \textbullet & 0.996 & 0.0262 & 0.0093--0.0754 & 0.701 & 0.994 & 0.0256 & 0.702 \\
 & [2,4] & 0.977 & 0.0311 & 0.0099--0.1233 & 0.589 & 0.976 & 0.0554 & 0.549 \\
 & [2,5] & 0.989 & 0.0427 & 0.0188--0.1435 & 0.570 & 0.998 & 0.0422 & 0.615 \\
 & [3,1] \textbullet & 0.867 & 1.2396 & 0.1535--42.6027 & 0.115 & 0.898 & 3.5051 & 0.000 \\
 & [3,2] \textbullet & 0.688 & 0.0827 & 0.0131--0.6355 & 0.258 & 0.626 & 0.0575 & 0.384 \\
 & [3,3] \textbullet & 0.995 & 0.0310 & 0.0033--0.0978 & 0.678 & 0.967 & 0.0118 & 0.836 \\
 & [3,4] \textbullet & 0.992 & 0.0247 & 0.0085--0.0776 & 0.685 & 0.993 & 0.0248 & 0.704 \\
 & [3,5] & 0.989 & 0.0311 & 0.0112--0.1037 & 0.612 & 0.981 & 0.0469 & 0.544 \\
 & [4,1] & 0.320 & 0.0848 & 0.0128--0.4777 & 0.052 & 0.298 & 0.3392 & 0.000 \\
 & [4,2] & 0.823 & 0.8550 & 0.1436--34.9851 & 0.083 & 0.894 & 3.5889 & 0.116 \\
 & [4,3] & 0.725 & 0.1169 & 0.0059--0.8531 & 0.253 & 0.659 & 0.0627 & 0.393 \\
 & [4,4] \textbullet & 0.976 & 0.0389 & 0.0037--0.0970 & 0.628 & 0.953 & 0.0179 & 0.778 \\
 & [4,5] \textbullet & 0.994 & 0.0264 & 0.0092--0.0759 & 0.700 & 0.998 & 0.0255 & 0.695 \\
 & [5,1] & 0.365 & 0.2448 & 0.0441--5.9324 & 0.068 & 0.506 & 34.3650 & 0.006 \\
 & [5,2] & 0.451 & 0.2768 & 0.0419--4.6062 & 0.049 & 0.361 & 0.3264 & 0.116 \\
 & [5,3] & 0.781 & 0.6483 & 0.0640--12.9048 & 0.117 & 0.911 & 3.6748 & 0.168 \\
 & [5,4] & 0.694 & 0.1078 & 0.0085--0.5634 & 0.235 & 0.698 & 0.0685 & 0.402 \\
 & [5,5] & 0.960 & 0.0351 & 0.0040--0.1155 & 0.628 & 0.978 & 0.0219 & 0.771 \\
 & [5,6] & 0.996 & 0.0268 & 0.0093--0.0828 & 0.681 & 0.993 & 0.0243 & 0.706 \\
 & [7,8] & 0.991 & 0.0271 & 0.0109--0.0924 & 0.653 & 0.996 & 0.0262 & 0.670 \\
 & [9,10] & 0.988 & 0.0357 & 0.0133--0.1150 & 0.625 & 0.981 & 0.0255 & 0.678 \\
\addlinespace[2pt]
\bottomrule
\end{tabular}

\end{adjustbox}
\end{table}
\end{landscape}
\begin{landscape}
\begin{table}[p]
\centering
\caption{\texttt{f19\_ladders\_fits.tex}: Order ladders of the fits (Phase 0b): Richardson terms, fixed exponents and the parametric models, core and out-of-design}
\label{tab:f19_ladders_fits}
\begin{adjustbox}{max width=\linewidth, max totalheight=0.86\textheight}
% AUTO-GENERATED -- do not hand-edit.  generated by scripts/make_paper_tables.py from the run started 2026-10-01T22:16:39 on code d82fd26 (full; 29,256 s; 7 workers)
% Order ladders of the fits (Phase 0b): Richardson terms, fixed exponents and the parametric models, core and held-out
% source : results/phase0b/order_ladders_panels.csv (Phase 0b: every order of every family with an order parameter, evaluated on Phase 1's grid at the headline stratum; noise pooled; uncapped cells; the roster orders reproduce phase1_records.csv exactly, results/phase0b/order_ladders_agreement.txt)
% filter : regime_set core (48 cells) | held-out (18 cells), noise == pooled; bullet = roster method; rho_V = valid rate over all records; med. err, q25--q75 over the valid records (conditional on validity); win vs last = fraction of all records where the variant is valid and below the last-value error
\begin{tabular}{llrrrrrrr}
\toprule
Family & order & \multicolumn{4}{c}{core} & \multicolumn{3}{c}{held-out} \\
\cmidrule(lr){3-6}\cmidrule(lr){7-9}
 & & $\rhoV$ & med.\ err & q25--q75 & win vs last & $\rhoV$ & med.\ err & win vs last \\
\midrule
Richardson, free exponents & 1 term \textbullet & 1.000 & 0.0280 & 0.0044--0.0763 & 0.811 & 1.000 & 0.0169 & 0.832 \\
 & 2 terms \textbullet & 0.992 & 0.0249 & 0.0056--0.0691 & 0.806 & 0.989 & 0.0195 & 0.840 \\
 & 3 terms \textbullet & 0.980 & 0.0205 & 0.0060--0.0612 & 0.774 & 0.972 & 0.0194 & 0.816 \\
 & 4 terms & 0.890 & 0.0172 & 0.0060--0.0520 & 0.706 & 0.937 & 0.0163 & 0.808 \\
\addlinespace[2pt]
Richardson, fixed exponent & $\alpha = 0.25$ & 1.000 & 0.0383 & 0.0195--0.0753 & 0.594 & 1.000 & 0.0263 & 0.603 \\
 & $\alpha = 0.5$ \textbullet & 1.000 & 0.0384 & 0.0120--0.0524 & 0.712 & 1.000 & 0.0261 & 0.981 \\
 & $\alpha = 0.75$ & 1.000 & 0.0400 & 0.0155--0.0811 & 0.781 & 1.000 & 0.0354 & 1.000 \\
 & $\alpha = 1$ \textbullet & 1.000 & 0.0557 & 0.0185--0.1056 & 0.846 & 1.000 & 0.0612 & 1.000 \\
 & $\alpha = 1.5$ & 1.000 & 0.0923 & 0.0087--0.1342 & 0.849 & 1.000 & 0.1037 & 0.819 \\
 & $\alpha = 2$ \textbullet & 1.000 & 0.1112 & 0.0254--0.1456 & 0.719 & 1.000 & 0.1372 & 0.711 \\
 & $\alpha = 3$ & 1.000 & 0.1139 & 0.0327--0.1655 & 0.174 & 1.000 & 0.1589 & 0.019 \\
\addlinespace[2pt]
parametric fit & single-exp \textbullet & 1.000 & 0.0348 & 0.0064--0.0867 & 0.962 & 1.000 & 0.0345 & 0.974 \\
 & double-exp \textbullet & 0.979 & 0.0265 & 0.0044--0.0729 & 0.911 & 1.000 & 0.0303 & 0.963 \\
 & triple-exp & 0.908 & 0.0248 & 0.0041--0.0729 & 0.857 & 0.955 & 0.0275 & 0.916 \\
 & rational one-term \textbullet & 1.000 & 0.0262 & 0.0063--0.0439 & 0.877 & 1.000 & 0.0097 & 0.877 \\
 & rational two-term & 1.000 & 0.0218 & 0.0050--0.0545 & 0.859 & 0.983 & 0.0106 & 0.832 \\
 & log \textbullet & 1.000 & 0.0688 & 0.0182--0.0995 & 0.692 & 1.000 & 0.0528 & 0.697 \\
\addlinespace[2pt]
\bottomrule
\end{tabular}

\end{adjustbox}
\end{table}
\end{landscape}
\begin{landscape}
\begin{table}[p]
\centering
\caption{\texttt{f20\_roster.tex}: Method roster by family, with type and member count, derived from the registry}
\label{tab:f20_roster}
\begin{adjustbox}{max width=\linewidth, max totalheight=0.86\textheight}
% AUTO-GENERATED -- do not hand-edit.  generated by scripts/make_paper_tables.py from the run started 2026-10-01T22:16:39 on code d82fd26 (full; 29,256 s; 7 workers)
% Method roster by family (Table tab:roster), with type and member count, derived from the registry
% source : src/accelerators.py (METHODS registry)
% source : src/evaluation.py (FAMILY, USES_FUTURE_X)
% filter : families in order of first appearance in the registry; type TE = evaluates at the target index (USES_FUTURE_X), LE = estimates the limit, mixed = a family with both; K = accelerators in the family; the trivial comparators are listed as comparators and not counted
\begin{tabular}{llrp{0.62\linewidth}}
\toprule
Family & type & $K$ & members \\
\midrule
baseline & mixed & 3 & \meth{linear}, \meth{log\_linear}, \meth{geom\_avg\_diff} \\
richardson & \TE & 6 & \meth{richardson\_1}, \meth{richardson\_2}, \meth{richardson\_3}, \meth{richardson\_a05}, \meth{richardson\_a10}, \meth{richardson\_a20} \\
parametric & \TE & 4 & \meth{single\_exp\_fit}, \meth{double\_exp\_fit}, \meth{rational\_fit}, \meth{log\_fit} \\
shanks & \LE & 4 & \meth{shanks\_1}, \meth{shanks\_2}, \meth{shanks\_3}, \meth{shanks\_4} \\
wynn\_eps & \LE & 3 & \meth{wynn\_eps\_1}, \meth{wynn\_eps\_2}, \meth{wynn\_eps\_3} \\
wynn\_rho & \LE & 3 & \meth{wynn\_rho\_1}, \meth{wynn\_rho\_2}, \meth{wynn\_rho\_3} \\
pade & \TE & 12 & \meth{pade\_11}, \meth{pade\_12}, \meth{pade\_13}, \meth{pade\_21}, \meth{pade\_22}, \meth{pade\_23}, \meth{pade\_31}, \meth{pade\_32}, \meth{pade\_33}, \meth{pade\_34}, \meth{pade\_44}, \meth{pade\_45} \\
levin & \LE & 6 & \meth{levin\_t1}, \meth{levin\_t2}, \meth{levin\_u1}, \meth{levin\_u2}, \meth{levin\_v1}, \meth{levin\_v2} \\
brezinski & \LE & 2 & \meth{brezinski\_theta1}, \meth{brezinski\_theta2} \\
neville & \TE & 3 & \meth{neville\_2}, \meth{neville\_3}, \meth{neville\_4} \\
anderson & \LE & 3 & \meth{anderson\_1}, \meth{anderson\_2}, \meth{anderson\_3} \\
ensemble & mixed & 3 & \meth{median\_ensemble}, \meth{stability\_weighted}, \meth{best\_shanks\_wynn} \\
\midrule
\textbf{accelerators} & & 52 & 12 families \\
\addlinespace[2pt]
trivial comparators (not counted) & & 5 & \meth{constant\_assumed}, \meth{constant\_oracle}, \meth{window\_mean}, \meth{window\_min}, \meth{last\_value} \\
\bottomrule
\end{tabular}

\end{adjustbox}
\end{table}
\end{landscape}
\begin{landscape}
\begin{table}[p]
\centering
\caption{\texttt{f21\_noop\_by\_noise.tex}: Classical no-op by noise class: the classical variants against the +/-10 \% band in each noise class, per set and stratum, with the last value and rational\_fit alongside}
\label{tab:f21_noop_by_noise}
\begin{adjustbox}{max width=\linewidth, max totalheight=0.86\textheight}
% AUTO-GENERATED -- do not hand-edit.  generated by scripts/make_paper_tables.py from the run started 2026-10-01T22:16:39 on code d82fd26 (full; 29,256 s; 7 workers)
% Classical no-op by noise class: the classical variants against the +/-10 % band in each noise class, per set and stratum, with the last value and rational_fit alongside
% source : results/phase1/phase1_aggregated.csv (per (method, regime, noise, g) cell; capped == 0)
% filter : capped == 0; noise class of a (family, noise) cell: noise-free = sigma 0 and no intrinsic noise; intrinsic noise only = sigma 0 and one of the intrinsic-noise families (noisy_plateau, heavy_noise_plat); otherwise sigma = the level; families = distinct curve families in the class, cells = (family, noise) cells; for a classical variant and a class: MI = median over the class's cells of med_improve, win = mean over cells of win_rate_vs_last, error = median over cells of med_error; in band = variants with MI in [0.9, 1.1] (k/N, N = 21); MI and win summarised over the variants; median error of the last value, of the classical variants (median over variants), of the best variant and of rational_fit; the 'all cells' row reproduces the band count of f04 (core) and f04b (out-of-design)
\begin{tabular}{lllrrrrrrrrrrrrr}
\toprule
Set & $g$ & noise class & families & cells & in band & \multicolumn{3}{c}{MI over the variants} & \multicolumn{2}{c}{win vs last} & \multicolumn{4}{c}{median error} & \meth{rational\_fit} \\
\cmidrule(lr){7-9}\cmidrule(lr){10-11}\cmidrule(lr){12-15}
 & & & & & $k/N$ & min & median & max & mean & max & last value & classical & best classical & \meth{rational\_fit} & win vs last \\
\midrule
core & 0.5 & noise-free & 16 & 16 & 12/21 & 0.859 & 1.054 & 1.926 & 0.642 & 0.812 & 0.0742 & 0.0419 & 0.0237 & 0.0185 & 0.858 \\
 &  & intrinsic noise only & 2 & 2 & 14/21 & 0.206 & 1.029 & 1.713 & 0.466 & 0.617 & 0.0023 & 0.0020 & 0.0013 & 0.0007 & 0.833 \\
 &  & $\sigma = 0.001$ & 18 & 18 & 21/21 & 0.921 & 0.985 & 1.017 & 0.380 & 0.520 & 0.0699 & 0.0711 & 0.0702 & 0.0163 & 0.843 \\
 &  & $\sigma = 0.005$ & 18 & 18 & 20/21 & 0.845 & 0.983 & 1.073 & 0.397 & 0.520 & 0.0682 & 0.0689 & 0.0659 & 0.0157 & 0.848 \\
 &  & all cells & 18 & 54 & 21/21 & 0.907 & 0.988 & 1.048 & 0.466 & 0.604 & 0.0690 & 0.0579 & 0.0482 & 0.0159 & 0.849 \\
\addlinespace[2pt]
 & 0.1 & noise-free & 14 & 14 & 5/21 & 1.007 & 2.777 & 8.827 & 0.752 & 0.857 & 0.1255 & 0.0261 & 0.0147 & 0.0334 & 0.900 \\
 &  & intrinsic noise only & 2 & 2 & 14/21 & 0.193 & 1.028 & 1.614 & 0.468 & 0.617 & 0.0023 & 0.0020 & 0.0013 & 0.0009 & 0.750 \\
 &  & $\sigma = 0.001$ & 16 & 16 & 21/21 & 0.954 & 0.995 & 1.010 & 0.370 & 0.519 & 0.1038 & 0.1055 & 0.1029 & 0.0303 & 0.879 \\
 &  & $\sigma = 0.005$ & 16 & 16 & 20/21 & 0.871 & 0.990 & 1.042 & 0.380 & 0.517 & 0.1022 & 0.1033 & 0.1001 & 0.0296 & 0.871 \\
 &  & all cells & 16 & 48 & 21/21 & 0.941 & 0.995 & 1.037 & 0.489 & 0.613 & 0.1030 & 0.0502 & 0.0438 & 0.0296 & 0.877 \\
\addlinespace[2pt]
 & 0.02 & noise-free & 13 & 13 & 4/21 & 1.010 & 2.744 & 17.304 & 0.741 & 0.846 & 0.1442 & 0.0319 & 0.0116 & 0.0303 & 0.869 \\
 &  & intrinsic noise only & 2 & 2 & 14/21 & 0.194 & 1.028 & 1.596 & 0.467 & 0.617 & 0.0023 & 0.0020 & 0.0013 & 0.0011 & 0.700 \\
 &  & $\sigma = 0.001$ & 15 & 15 & 21/21 & 0.957 & 0.996 & 1.010 & 0.371 & 0.516 & 0.0971 & 0.0994 & 0.0943 & 0.0228 & 0.840 \\
 &  & $\sigma = 0.005$ & 15 & 15 & 20/21 & 0.869 & 0.990 & 1.038 & 0.377 & 0.513 & 0.0954 & 0.0972 & 0.0939 & 0.0230 & 0.842 \\
 &  & all cells & 15 & 45 & 21/21 & 0.944 & 0.995 & 1.033 & 0.484 & 0.607 & 0.0970 & 0.0523 & 0.0456 & 0.0228 & 0.843 \\
\addlinespace[2pt]
\midrule
out-of-design & 0.5 & noise-free & 6 & 6 & 5/21 & 0.685 & 1.911 & 5.146 & 0.793 & 1.000 & 0.0394 & 0.0245 & 0.0052 & 0.0057 & 0.850 \\
 &  & $\sigma = 0.001$ & 6 & 6 & 20/21 & 0.896 & 0.980 & 1.025 & 0.385 & 0.567 & 0.0399 & 0.0429 & 0.0369 & 0.0057 & 0.850 \\
 &  & $\sigma = 0.005$ & 6 & 6 & 18/21 & 0.768 & 0.974 & 1.091 & 0.378 & 0.561 & 0.0390 & 0.0447 & 0.0359 & 0.0063 & 0.789 \\
 &  & all cells & 6 & 18 & 20/21 & 0.894 & 0.983 & 1.094 & 0.519 & 0.709 & 0.0391 & 0.0349 & 0.0248 & 0.0062 & 0.830 \\
\addlinespace[2pt]
 & 0.1 & noise-free & 5 & 5 & 4/21 & 0.969 & 2.696 & 7.376 & 0.895 & 1.000 & 0.0862 & 0.0217 & 0.0088 & 0.0053 & 0.860 \\
 &  & $\sigma = 0.001$ & 5 & 5 & 21/21 & 0.937 & 0.988 & 1.016 & 0.381 & 0.567 & 0.0862 & 0.0877 & 0.0724 & 0.0074 & 0.860 \\
 &  & $\sigma = 0.005$ & 5 & 5 & 20/21 & 0.842 & 0.990 & 1.067 & 0.376 & 0.567 & 0.0846 & 0.0909 & 0.0847 & 0.0102 & 0.840 \\
 &  & all cells & 5 & 15 & 20/21 & 0.937 & 0.992 & 1.102 & 0.550 & 0.711 & 0.0862 & 0.0483 & 0.0465 & 0.0074 & 0.853 \\
\addlinespace[2pt]
 & 0.02 & noise-free & 5 & 5 & 3/21 & 1.019 & 2.413 & 6.865 & 0.904 & 1.000 & 0.0935 & 0.0243 & 0.0045 & 0.0051 & 0.907 \\
 &  & $\sigma = 0.001$ & 5 & 5 & 21/21 & 0.942 & 0.989 & 1.015 & 0.381 & 0.567 & 0.0935 & 0.0950 & 0.0747 & 0.0056 & 0.907 \\
 &  & $\sigma = 0.005$ & 5 & 5 & 20/21 & 0.854 & 0.991 & 1.062 & 0.376 & 0.567 & 0.0919 & 0.0982 & 0.0920 & 0.0065 & 0.887 \\
 &  & all cells & 5 & 15 & 21/21 & 0.942 & 0.992 & 1.093 & 0.554 & 0.711 & 0.0935 & 0.0601 & 0.0510 & 0.0061 & 0.900 \\
\bottomrule
\end{tabular}

\end{adjustbox}
\end{table}
\end{landscape}
\begin{landscape}
\begin{table}[p]
\centering
\caption{\texttt{f21b\_noop\_by\_noise\_variants\_core.tex}: Classical no-op by noise class, core families, headline stratum: every classical variant's median improvement factor and win rate against the last value in each noise class}
\label{tab:f21b_noop_by_noise_variants_core}
\begin{adjustbox}{max width=\linewidth, max totalheight=0.86\textheight}
% AUTO-GENERATED -- do not hand-edit.  generated by scripts/make_paper_tables.py from the run started 2026-10-01T22:16:39 on code d82fd26 (full; 29,256 s; 7 workers)
% Classical no-op by noise class, core families, headline stratum: every classical variant's median improvement factor and win rate against the last value in each noise class
% source : results/phase1/phase1_aggregated.csv (per (method, regime, noise, g) cell; capped == 0)
% filter : is_holdout == 0, target_g == 0.1, capped == 0; noise class of a (family, noise) cell: noise-free = sigma 0 and no intrinsic noise; intrinsic noise only = sigma 0 and one of the intrinsic-noise families (noisy_plateau, heavy_noise_plat); otherwise sigma = the level; MI = median over the class's cells of med_improve, bold = inside [0.9, 1.1]; win vs last = mean over the class's cells of win_rate_vs_last, bold = inside [0.4, 0.6]; variants in the order of f04
\begin{tabular}{llrrrrrrrr}
\toprule
Family & Method & \multicolumn{2}{c}{noise-free} & \multicolumn{2}{c}{intrinsic noise only} & \multicolumn{2}{c}{$\sigma = 0.001$} & \multicolumn{2}{c}{$\sigma = 0.005$} \\
\cmidrule(lr){3-4}\cmidrule(lr){5-6}\cmidrule(lr){7-8}\cmidrule(lr){9-10}
 & & MI & win vs last & MI & win vs last & MI & win vs last & MI & win vs last \\
\midrule
Shanks & \meth{shanks\_1} & 1.908 & 0.857 & \textbf{1.028} & \textbf{0.433} & \textbf{1.000} & \textbf{0.519} & \textbf{1.000} & \textbf{0.517} \\
Shanks & \meth{shanks\_2} & 3.810 & 0.714 & \textbf{0.993} & \textbf{0.500} & \textbf{0.983} & 0.179 & \textbf{0.953} & 0.194 \\
Shanks & \meth{shanks\_3} & 7.804 & 0.714 & \textbf{1.038} & \textbf{0.500} & \textbf{0.984} & 0.244 & \textbf{0.966} & 0.273 \\
Shanks & \meth{shanks\_4} & 8.507 & 0.700 & \textbf{1.013} & \textbf{0.500} & \textbf{0.966} & 0.188 & \textbf{0.962} & 0.248 \\
\addlinespace[2pt]
Wynn $\varepsilon$ & \meth{wynn\_eps\_1} & 1.990 & 0.857 & \textbf{1.028} & \textbf{0.433} & \textbf{1.000} & \textbf{0.519} & \textbf{1.000} & \textbf{0.517} \\
Wynn $\varepsilon$ & \meth{wynn\_eps\_2} & 2.985 & 0.748 & 1.192 & 0.617 & \textbf{0.990} & 0.367 & \textbf{0.965} & 0.315 \\
Wynn $\varepsilon$ & \meth{wynn\_eps\_3} & 5.075 & 0.786 & 1.614 & \textbf{0.600} & \textbf{0.991} & 0.387 & \textbf{0.967} & 0.315 \\
\addlinespace[2pt]
Wynn $\rho$ & \meth{wynn\_rho\_1} & 1.615 & 0.643 & 1.132 & \textbf{0.550} & \textbf{0.998} & 0.340 & \textbf{0.993} & 0.342 \\
Wynn $\rho$ & \meth{wynn\_rho\_2} & \textbf{1.012} & \textbf{0.500} & \textbf{1.088} & \textbf{0.533} & \textbf{1.004} & \textbf{0.500} & \textbf{0.984} & 0.358 \\
Wynn $\rho$ & \meth{wynn\_rho\_3} & 7.335 & 0.707 & \textbf{0.929} & \textbf{0.433} & \textbf{0.995} & \textbf{0.423} & \textbf{0.988} & 0.385 \\
\addlinespace[2pt]
Levin & \meth{levin\_t1} & 1.990 & 0.857 & \textbf{1.028} & \textbf{0.433} & \textbf{1.000} & \textbf{0.519} & \textbf{1.000} & \textbf{0.517} \\
Levin & \meth{levin\_t2} & 2.777 & 0.857 & \textbf{1.013} & \textbf{0.517} & \textbf{0.989} & 0.312 & \textbf{0.990} & 0.381 \\
Levin & \meth{levin\_u1} & \textbf{1.039} & \textbf{0.500} & \textbf{1.029} & \textbf{0.433} & \textbf{1.000} & \textbf{0.513} & \textbf{1.000} & \textbf{0.506} \\
Levin & \meth{levin\_u2} & 7.885 & 0.714 & \textbf{1.010} & \textbf{0.517} & \textbf{0.989} & 0.312 & \textbf{0.990} & 0.383 \\
Levin & \meth{levin\_v1} & 7.330 & 0.714 & \textbf{0.942} & \textbf{0.483} & \textbf{0.984} & 0.185 & \textbf{0.954} & 0.298 \\
Levin & \meth{levin\_v2} & 8.827 & 0.786 & \textbf{0.942} & \textbf{0.467} & \textbf{0.981} & 0.212 & \textbf{0.966} & 0.302 \\
\addlinespace[2pt]
Brezinski $\theta$ & \meth{brezinski\_theta1} & 1.938 & 0.857 & 0.332 & 0.217 & \textbf{1.010} & \textbf{0.504} & \textbf{1.042} & \textbf{0.494} \\
Brezinski $\theta$ & \meth{brezinski\_theta2} & 3.888 & 0.714 & 0.193 & 0.117 & \textbf{0.954} & 0.194 & 0.871 & 0.254 \\
\addlinespace[2pt]
Anderson & \meth{anderson\_1} & \textbf{1.007} & 0.857 & \textbf{1.028} & \textbf{0.433} & \textbf{1.000} & \textbf{0.519} & \textbf{1.000} & \textbf{0.517} \\
Anderson & \meth{anderson\_2} & \textbf{1.014} & 0.857 & 1.239 & \textbf{0.567} & \textbf{1.000} & \textbf{0.450} & \textbf{0.999} & \textbf{0.442} \\
Anderson & \meth{anderson\_3} & \textbf{1.021} & 0.857 & 1.234 & \textbf{0.550} & \textbf{0.999} & 0.383 & \textbf{0.998} & \textbf{0.417} \\
\midrule
\multicolumn{2}{l}{in the $\pm 10\%$ band (MI in $[0.9, 1.1]$)} & 5/21 &  & 14/21 &  & 21/21 &  & 20/21 &  \\
\multicolumn{2}{l}{win rate vs last in $[0.4, 0.6]$ (coin flip)} &  & 2/21 &  & 18/21 &  & 9/21 &  & 8/21 \\
\bottomrule
\end{tabular}

\end{adjustbox}
\end{table}
\end{landscape}
\begin{landscape}
\begin{table}[p]
\centering
\caption{\texttt{f21b\_noop\_by\_noise\_variants\_holdout.tex}: Classical no-op by noise class, out-of-design families, headline stratum: every classical variant's median improvement factor and win rate against the last value in each noise class}
\label{tab:f21b_noop_by_noise_variants_holdout}
\begin{adjustbox}{max width=\linewidth, max totalheight=0.86\textheight}
% AUTO-GENERATED -- do not hand-edit.  generated by scripts/make_paper_tables.py from the run started 2026-10-01T22:16:39 on code d82fd26 (full; 29,256 s; 7 workers)
% Classical no-op by noise class, held-out families, headline stratum: every classical variant's median improvement factor and win rate against the last value in each noise class
% source : results/phase1/phase1_aggregated.csv (per (method, regime, noise, g) cell; capped == 0)
% filter : is_holdout == 1, target_g == 0.1, capped == 0; noise class of a (family, noise) cell: noise-free = sigma 0 and no intrinsic noise; intrinsic noise only = sigma 0 and one of the intrinsic-noise families (noisy_plateau, heavy_noise_plat); otherwise sigma = the level; MI = median over the class's cells of med_improve, bold = inside [0.9, 1.1]; win vs last = mean over the class's cells of win_rate_vs_last, bold = inside [0.4, 0.6]; variants in the order of f04
\begin{tabular}{llrrrrrr}
\toprule
Family & Method & \multicolumn{2}{c}{noise-free} & \multicolumn{2}{c}{$\sigma = 0.001$} & \multicolumn{2}{c}{$\sigma = 0.005$} \\
\cmidrule(lr){3-4}\cmidrule(lr){5-6}\cmidrule(lr){7-8}
 & & MI & win vs last & MI & win vs last & MI & win vs last \\
\midrule
Shanks & \meth{shanks\_1} & 2.314 & 1.000 & \textbf{1.001} & \textbf{0.567} & \textbf{1.001} & \textbf{0.567} \\
Shanks & \meth{shanks\_2} & 3.407 & 0.800 & \textbf{0.978} & 0.147 & \textbf{0.933} & 0.153 \\
Shanks & \meth{shanks\_3} & 7.376 & 1.000 & \textbf{0.978} & 0.233 & \textbf{0.949} & 0.240 \\
Shanks & \meth{shanks\_4} & 7.244 & 0.993 & \textbf{0.964} & 0.133 & \textbf{0.949} & 0.193 \\
\addlinespace[2pt]
Wynn $\varepsilon$ & \meth{wynn\_eps\_1} & 2.314 & 1.000 & \textbf{1.001} & \textbf{0.567} & \textbf{1.001} & \textbf{0.567} \\
Wynn $\varepsilon$ & \meth{wynn\_eps\_2} & 2.791 & 0.800 & \textbf{0.987} & 0.393 & \textbf{0.940} & 0.300 \\
Wynn $\varepsilon$ & \meth{wynn\_eps\_3} & 3.281 & 1.000 & \textbf{0.988} & \textbf{0.407} & \textbf{0.946} & 0.280 \\
\addlinespace[2pt]
Wynn $\rho$ & \meth{wynn\_rho\_1} & 6.932 & 0.800 & \textbf{0.998} & 0.380 & \textbf{0.994} & 0.367 \\
Wynn $\rho$ & \meth{wynn\_rho\_2} & 2.749 & \textbf{0.600} & \textbf{1.010} & \textbf{0.560} & \textbf{0.988} & 0.313 \\
Wynn $\rho$ & \meth{wynn\_rho\_3} & 2.696 & \textbf{0.600} & \textbf{0.987} & \textbf{0.453} & \textbf{0.988} & 0.347 \\
\addlinespace[2pt]
Levin & \meth{levin\_t1} & 2.314 & 1.000 & \textbf{1.001} & \textbf{0.567} & \textbf{1.001} & \textbf{0.567} \\
Levin & \meth{levin\_t2} & 2.568 & 1.000 & \textbf{0.988} & 0.300 & \textbf{0.991} & 0.340 \\
Levin & \meth{levin\_u1} & \textbf{0.969} & \textbf{0.400} & \textbf{1.001} & \textbf{0.560} & \textbf{1.001} & \textbf{0.553} \\
Levin & \meth{levin\_u2} & 2.671 & 1.000 & \textbf{0.988} & 0.300 & \textbf{0.990} & 0.340 \\
Levin & \meth{levin\_v1} & 2.696 & 1.000 & \textbf{0.979} & 0.140 & \textbf{0.928} & 0.253 \\
Levin & \meth{levin\_v2} & 4.118 & 1.000 & \textbf{0.977} & 0.120 & \textbf{0.944} & 0.260 \\
\addlinespace[2pt]
Brezinski $\theta$ & \meth{brezinski\_theta1} & 2.222 & 1.000 & \textbf{1.016} & \textbf{0.567} & \textbf{1.067} & \textbf{0.560} \\
Brezinski $\theta$ & \meth{brezinski\_theta2} & 3.732 & 0.800 & \textbf{0.937} & 0.180 & 0.842 & 0.287 \\
\addlinespace[2pt]
Anderson & \meth{anderson\_1} & \textbf{1.020} & 1.000 & \textbf{1.000} & \textbf{0.567} & \textbf{1.001} & \textbf{0.567} \\
Anderson & \meth{anderson\_2} & \textbf{1.040} & 1.000 & \textbf{1.000} & \textbf{0.453} & \textbf{0.999} & \textbf{0.440} \\
Anderson & \meth{anderson\_3} & \textbf{1.059} & 1.000 & \textbf{0.997} & 0.400 & \textbf{0.997} & 0.393 \\
\midrule
\multicolumn{2}{l}{in the $\pm 10\%$ band (MI in $[0.9, 1.1]$)} & 4/21 &  & 21/21 &  & 20/21 &  \\
\multicolumn{2}{l}{win rate vs last in $[0.4, 0.6]$ (coin flip)} &  & 3/21 &  & 10/21 &  & 7/21 \\
\bottomrule
\end{tabular}

\end{adjustbox}
\end{table}
\end{landscape}
\begin{landscape}
\begin{table}[p]
\centering
\caption{\texttt{f22\_per\_family\_g0.02.tex}: rational\_fit and single\_exp\_fit family by family at g = 0.02: the last value, rational\_fit (median error, win rate, validity, catastrophe rate), single\_exp\_fit (median error, win rate) and the best rank-eligible accelerator per curve family, with the families where each named method loses}
\label{tab:f22_per_family_g0.02}
\begin{adjustbox}{max width=\linewidth, max totalheight=0.86\textheight}
% AUTO-GENERATED -- do not hand-edit.  generated by scripts/make_paper_tables.py from the run started 2026-10-01T22:16:39 on code d82fd26 (full; 29,256 s; 7 workers)
% rational_fit and single_exp_fit family by family at g = 0.02: the last value, rational_fit (median error, win rate, validity, catastrophe rate), single_exp_fit (median error, win rate) and the best rank-eligible accelerator per curve family, with the families where each named method loses
% source : results/phase1/phase1_aggregated.csv (per (method, regime, noise, g) cell; capped == 0)
% source : results/phase1/phase1_global.csv, results/phase1/phase1_global_holdout.csv (rank_eligible == 1: the candidates of the last two columns)
% filter : target_g == 0.02; one row per curve family (core, then out-of-design); cells = the family's uncapped (noise) cells; family-level statistics over the family's uncapped (noise) cells: median error = median of the cell med_error, win rate = mean of win_rate_vs_last, validity and catastrophe rate = means over cells (the aggregation of the pooled tables, median of cell medians); single_exp_fit = the conservative alternative the paper names (median error and win rate vs the last value, the same statistics); best rank-eligible accelerator = the lowest family median error among the accelerators that are rank-eligible in the pooled table of that set and stratum (ties by method name); a family with no uncapped cell prints 'capped'; the footers give, per set and named method, the uncapped families where its win rate against the last value exceeds 0.5 (the family win rate = mean over the family's cells of the per-cell win rate)
\begin{tabular}{llrrrrrrrrlr}
\toprule
Set & curve family & cells & last value & \multicolumn{4}{c}{\meth{rational\_fit}} & \multicolumn{2}{c}{\meth{single\_exp\_fit}} & \multicolumn{2}{c}{best rank-eligible accelerator} \\
\cmidrule(lr){5-8}\cmidrule(lr){9-10}\cmidrule(lr){11-12}
 & & & med.\ err & med.\ err & win vs last & $\rhoV$ & $\rhoC$ & med.\ err & win vs last & method & med.\ err \\
\midrule
core & \texttt{single\_exp} & 3 & 0.0201 & 0.0227 & 0.433 & 1.000 & 0.000 & 0.0004 & 1.000 & \meth{single\_exp\_fit} & 0.0004 \\
core & \texttt{two\_exp} & 3 & 0.0970 & 0.0228 & 1.000 & 1.000 & 0.000 & 0.0051 & 1.000 & \meth{single\_exp\_fit} & 0.0051 \\
core & \texttt{three\_exp} & 3 & 0.1442 & 0.0312 & 1.000 & 1.000 & 0.000 & 0.0970 & 1.000 & \meth{pade\_23} & 0.0209 \\
core & \texttt{four\_exp} & 3 & 0.1574 & 0.0356 & 1.000 & 1.000 & 0.000 & 0.1106 & 1.000 & \meth{richardson\_2} & 0.0112 \\
core & \texttt{power\_law} & 3 & 0.0333 & 0.0084 & 1.000 & 1.000 & 0.000 & 0.0270 & 0.956 & \meth{richardson\_1} & 0.0029 \\
core & \texttt{rational\_decay} & 3 & 0.1704 & 0.0021 & 1.000 & 1.000 & 0.000 & 0.1179 & 1.000 & \meth{rational\_fit} & 0.0021 \\
core & \texttt{mixed\_pow\_rat} & 3 & 0.1461 & 0.0202 & 1.000 & 1.000 & 0.000 & 0.1036 & 1.000 & \meth{richardson\_a05} & 0.0028 \\
core & \texttt{multiphase} & 3 & 0.1951 & 0.1327 & 1.000 & 1.000 & 0.000 & 0.1867 & 0.989 & \meth{log\_linear} & 0.0119 \\
core & \texttt{osc\_exp} & 3 & 0.0507 & 0.0304 & 0.600 & 1.000 & 0.000 & 0.0368 & 1.000 & \meth{pade\_22} & 0.0049 \\
core & \texttt{damped\_osc\_pow} & 3 & 0.0546 & 0.0152 & 1.000 & 1.000 & 0.000 & 0.0442 & 0.989 & \meth{richardson\_1} & 0.0038 \\
core & \texttt{log\_slow} & 0 & capped & capped & capped & capped & capped & capped & capped & capped & capped \\
core & \texttt{delayed\_plateau} & 3 & 0.2391 & 0.0403 & 1.000 & 1.000 & 0.000 & 0.0264 & 1.000 & \meth{richardson\_a05} & 0.0205 \\
core & \texttt{staircase} & 3 & 0.1800 & 0.0519 & 1.000 & 1.000 & 0.000 & 0.0686 & 1.000 & \meth{double\_exp\_fit} & 0.0328 \\
core & \texttt{noisy\_plateau} & 3 & 0.0024 & 0.0010 & 0.611 & 1.000 & 0.111 & 0.0008 & 0.689 & \meth{richardson\_a20} & 0.0004 \\
core & \texttt{slow\_osc\_power} & 3 & 0.0275 & 0.0636 & 0.267 & 1.000 & 0.000 & 0.0066 & 0.944 & \meth{double\_exp\_fit} & 0.0066 \\
core & \texttt{log\_oscillatory} & 0 & capped & capped & capped & capped & capped & capped & capped & capped & capped \\
core & \texttt{heavy\_noise\_plat} & 3 & 0.0029 & 0.0013 & 0.733 & 1.000 & 0.067 & 0.0010 & 0.778 & \meth{richardson\_a20} & 0.0009 \\
core & \texttt{broken\_power\_law} & 0 & capped & capped & capped & capped & capped & capped & capped & capped & capped \\
\addlinespace[2pt]
out-of-design & \texttt{stretched\_exp} & 3 & 0.1531 & 0.0032 & 1.000 & 1.000 & 0.000 & 0.0933 & 1.000 & \meth{rational\_fit} & 0.0032 \\
out-of-design & \texttt{logistic\_tail} & 3 & 0.0935 & 0.0883 & 0.533 & 1.000 & 0.000 & 0.0159 & 0.867 & \meth{double\_exp\_fit} & 0.0159 \\
out-of-design & \texttt{inv\_sqrt\_log} & 3 & 0.0204 & 0.0056 & 0.967 & 1.000 & 0.011 & 0.0169 & 0.911 & \meth{pade\_22} & 0.0024 \\
out-of-design & \texttt{random\_knots} & 0 & capped & capped & capped & capped & capped & capped & capped & capped & capped \\
out-of-design & \texttt{real\_boot\_a} & 3 & 0.1634 & 0.0570 & 1.000 & 1.000 & 0.000 & 0.1194 & 1.000 & \meth{pade\_45} & 0.0315 \\
out-of-design & \texttt{real\_boot\_b} & 3 & 0.0501 & 0.0045 & 1.000 & 1.000 & 0.000 & 0.0383 & 1.000 & \meth{richardson\_a10} & 0.0005 \\
\midrule
\multicolumn{12}{l}{core, \meth{rational\_fit}: its win rate against the last value exceeds 0.5 in 13 of 15 uncapped families; not in \texttt{single\_exp} (0.433), \texttt{slow\_osc\_power} (0.267); its median error exceeds the last value's in \texttt{single\_exp} (0.0227 vs 0.0201), \texttt{slow\_osc\_power} (0.0636 vs 0.0275)} \\
\multicolumn{12}{l}{core, \meth{single\_exp\_fit}: its win rate against the last value exceeds 0.5 in 15 of 15 uncapped families; its median error exceeds the last value's in no family} \\
\multicolumn{12}{l}{out-of-design, \meth{rational\_fit}: its win rate against the last value exceeds 0.5 in 5 of 5 uncapped families; its median error exceeds the last value's in no family} \\
\multicolumn{12}{l}{out-of-design, \meth{single\_exp\_fit}: its win rate against the last value exceeds 0.5 in 5 of 5 uncapped families; its median error exceeds the last value's in no family} \\
\bottomrule
\end{tabular}

\end{adjustbox}
\end{table}
\end{landscape}
\begin{landscape}
\begin{table}[p]
\centering
\caption{\texttt{f22\_per\_family\_g0.1.tex}: rational\_fit and single\_exp\_fit family by family at g = 0.1: the last value, rational\_fit (median error, win rate, validity, catastrophe rate), single\_exp\_fit (median error, win rate) and the best rank-eligible accelerator per curve family, with the families where each named method loses}
\label{tab:f22_per_family_g0.1}
\begin{adjustbox}{max width=\linewidth, max totalheight=0.86\textheight}
% AUTO-GENERATED -- do not hand-edit.  generated by scripts/make_paper_tables.py from the run started 2026-10-01T22:16:39 on code d82fd26 (full; 29,256 s; 7 workers)
% rational_fit family by family at g = 0.1: the last value, rational_fit (median error, win rate, validity, catastrophe rate) and the best rank-eligible accelerator per curve family, with the families where the default loses
% source : results/phase1/phase1_aggregated.csv (per (method, regime, noise, g) cell; capped == 0)
% source : results/phase1/phase1_global.csv, results/phase1/phase1_global_holdout.csv (rank_eligible == 1: the candidates of the last two columns)
% filter : target_g == 0.1; one row per curve family (core, then held-out); cells = the family's uncapped (noise) cells; family-level statistics over the family's uncapped (noise) cells: median error = median of the cell med_error, win rate = mean of win_rate_vs_last, validity and catastrophe rate = means over cells (the aggregation of the pooled tables, median of cell medians); best rank-eligible accelerator = the lowest family median error among the accelerators that are rank-eligible in the pooled table of that set and stratum (ties by method name); a family with no uncapped cell prints 'capped'
\begin{tabular}{llrrrrrrlr}
\toprule
Set & curve family & cells & last value & \multicolumn{4}{c}{\meth{rational\_fit}} & \multicolumn{2}{c}{best rank-eligible accelerator} \\
\cmidrule(lr){5-8}\cmidrule(lr){9-10}
 & & & med.\ err & med.\ err & win vs last & $\rhoV$ & $\rhoC$ & method & med.\ err \\
\midrule
core & \texttt{single\_exp} & 3 & 0.0185 & 0.0275 & 0.333 & 1.000 & 0.000 & \meth{single\_exp\_fit} & 0.0003 \\
core & \texttt{two\_exp} & 3 & 0.0891 & 0.0275 & 1.000 & 1.000 & 0.000 & \meth{single\_exp\_fit} & 0.0042 \\
core & \texttt{three\_exp} & 3 & 0.1325 & 0.0339 & 1.000 & 1.000 & 0.000 & \meth{pade\_12} & 0.0219 \\
core & \texttt{four\_exp} & 3 & 0.1446 & 0.0329 & 1.000 & 1.000 & 0.000 & \meth{richardson\_2} & 0.0157 \\
core & \texttt{power\_law} & 3 & 0.0306 & 0.0066 & 1.000 & 1.000 & 0.000 & \meth{richardson\_1} & 0.0024 \\
core & \texttt{rational\_decay} & 3 & 0.1565 & 0.0019 & 1.000 & 1.000 & 0.000 & \meth{rational\_fit} & 0.0019 \\
core & \texttt{mixed\_pow\_rat} & 3 & 0.1342 & 0.0148 & 1.000 & 1.000 & 0.000 & \meth{richardson\_3} & 0.0079 \\
core & \texttt{multiphase} & 3 & 0.1792 & 0.1180 & 1.000 & 1.000 & 0.000 & \meth{log\_linear} & 0.0167 \\
core & \texttt{osc\_exp} & 3 & 0.0468 & 0.0363 & 0.989 & 1.000 & 0.000 & \meth{pade\_22} & 0.0026 \\
core & \texttt{damped\_osc\_pow} & 3 & 0.0501 & 0.0120 & 1.000 & 1.000 & 0.000 & \meth{richardson\_1} & 0.0026 \\
core & \texttt{log\_slow} & 0 & capped & capped & capped & capped & capped & capped & capped \\
core & \texttt{delayed\_plateau} & 3 & 0.2197 & 0.0481 & 1.000 & 1.000 & 0.000 & \meth{single\_exp\_fit} & 0.0142 \\
core & \texttt{staircase} & 3 & 0.1800 & 0.0519 & 1.000 & 1.000 & 0.000 & \meth{double\_exp\_fit} & 0.0328 \\
core & \texttt{noisy\_plateau} & 3 & 0.0024 & 0.0008 & 0.678 & 1.000 & 0.089 & \meth{richardson\_a20} & 0.0004 \\
core & \texttt{slow\_osc\_power} & 3 & 0.0252 & 0.0627 & 0.267 & 1.000 & 0.000 & \meth{double\_exp\_fit} & 0.0088 \\
core & \texttt{log\_oscillatory} & 0 & capped & capped & capped & capped & capped & capped & capped \\
core & \texttt{heavy\_noise\_plat} & 3 & 0.0029 & 0.0011 & 0.767 & 1.000 & 0.056 & \meth{double\_exp\_fit} & 0.0009 \\
core & \texttt{broken\_power\_law} & 3 & 0.1185 & 0.0729 & 1.000 & 1.000 & 0.000 & \meth{pade\_11} & 0.0601 \\
\addlinespace[2pt]
held-out & \texttt{stretched\_exp} & 3 & 0.1406 & 0.0074 & 1.000 & 1.000 & 0.000 & \meth{pade\_11} & 0.0045 \\
held-out & \texttt{logistic\_tail} & 3 & 0.0862 & 0.1076 & 0.300 & 1.000 & 0.000 & \meth{richardson\_a05} & 0.0065 \\
held-out & \texttt{inv\_sqrt\_log} & 3 & 0.0187 & 0.0044 & 0.967 & 1.000 & 0.011 & \meth{pade\_22} & 0.0021 \\
held-out & \texttt{random\_knots} & 0 & capped & capped & capped & capped & capped & capped & capped \\
held-out & \texttt{real\_boot\_a} & 3 & 0.1502 & 0.0488 & 1.000 & 1.000 & 0.000 & \meth{pade\_45} & 0.0253 \\
held-out & \texttt{real\_boot\_b} & 3 & 0.0460 & 0.0049 & 1.000 & 1.000 & 0.000 & \meth{richardson\_a10} & 0.0004 \\
\midrule
\multicolumn{10}{l}{core: \meth{rational\_fit} wins against the last value on more than half of the cells in 14 of 16 uncapped families; not in \texttt{single\_exp} (0.333), \texttt{slow\_osc\_power} (0.267); its median error exceeds the last value's in \texttt{single\_exp} (0.0275 vs 0.0185), \texttt{slow\_osc\_power} (0.0627 vs 0.0252)} \\
\multicolumn{10}{l}{held-out: \meth{rational\_fit} wins against the last value on more than half of the cells in 4 of 5 uncapped families; not in \texttt{logistic\_tail} (0.300); its median error exceeds the last value's in \texttt{logistic\_tail} (0.1076 vs 0.0862)} \\
\bottomrule
\end{tabular}

\end{adjustbox}
\end{table}
\end{landscape}
\begin{landscape}
\begin{table}[p]
\centering
\caption{\texttt{f22\_per\_family\_g0.5.tex}: rational\_fit and single\_exp\_fit family by family at g = 0.5: the last value, rational\_fit (median error, win rate, validity, catastrophe rate), single\_exp\_fit (median error, win rate) and the best rank-eligible accelerator per curve family, with the families where each named method loses}
\label{tab:f22_per_family_g0.5}
\begin{adjustbox}{max width=\linewidth, max totalheight=0.86\textheight}
% AUTO-GENERATED -- do not hand-edit.  generated by scripts/make_paper_tables.py from the run started 2026-10-01T22:16:39 on code d82fd26 (full; 29,256 s; 7 workers)
% rational_fit and single_exp_fit family by family at g = 0.5: the last value, rational_fit (median error, win rate, validity, catastrophe rate), single_exp_fit (median error, win rate) and the best rank-eligible accelerator per curve family, with the families where each named method loses
% source : results/phase1/phase1_aggregated.csv (per (method, regime, noise, g) cell; capped == 0)
% source : results/phase1/phase1_global.csv, results/phase1/phase1_global_holdout.csv (rank_eligible == 1: the candidates of the last two columns)
% filter : target_g == 0.5; one row per curve family (core, then out-of-design); cells = the family's uncapped (noise) cells; family-level statistics over the family's uncapped (noise) cells: median error = median of the cell med_error, win rate = mean of win_rate_vs_last, validity and catastrophe rate = means over cells (the aggregation of the pooled tables, median of cell medians); single_exp_fit = the conservative alternative the paper names (median error and win rate vs the last value, the same statistics); best rank-eligible accelerator = the lowest family median error among the accelerators that are rank-eligible in the pooled table of that set and stratum (ties by method name); a family with no uncapped cell prints 'capped'; the footers give, per set and named method, the uncapped families where its win rate against the last value exceeds 0.5 (the family win rate = mean over the family's cells of the per-cell win rate)
\begin{tabular}{llrrrrrrrrlr}
\toprule
Set & curve family & cells & last value & \multicolumn{4}{c}{\meth{rational\_fit}} & \multicolumn{2}{c}{\meth{single\_exp\_fit}} & \multicolumn{2}{c}{best rank-eligible accelerator} \\
\cmidrule(lr){5-8}\cmidrule(lr){9-10}\cmidrule(lr){11-12}
 & & & med.\ err & med.\ err & win vs last & $\rhoV$ & $\rhoC$ & med.\ err & win vs last & method & med.\ err \\
\midrule
core & \texttt{single\_exp} & 3 & 0.0105 & 0.0232 & 0.111 & 1.000 & 0.044 & 0.0002 & 1.000 & \meth{single\_exp\_fit} & 0.0002 \\
core & \texttt{two\_exp} & 3 & 0.0499 & 0.0148 & 1.000 & 1.000 & 0.000 & 0.0016 & 1.000 & \meth{single\_exp\_fit} & 0.0016 \\
core & \texttt{three\_exp} & 3 & 0.0738 & 0.0167 & 1.000 & 1.000 & 0.000 & 0.0343 & 1.000 & \meth{pade\_23} & 0.0104 \\
core & \texttt{four\_exp} & 3 & 0.0806 & 0.0088 & 1.000 & 1.000 & 0.000 & 0.0364 & 1.000 & \meth{richardson\_2} & 0.0055 \\
core & \texttt{power\_law} & 3 & 0.0170 & 0.0019 & 1.000 & 1.000 & 0.000 & 0.0107 & 0.944 & \meth{richardson\_1} & 0.0009 \\
core & \texttt{rational\_decay} & 3 & 0.0870 & 0.0008 & 1.000 & 1.000 & 0.000 & 0.0375 & 1.000 & \meth{rational\_fit} & 0.0008 \\
core & \texttt{mixed\_pow\_rat} & 3 & 0.0747 & 0.0039 & 1.000 & 1.000 & 0.000 & 0.0336 & 1.000 & \meth{pade\_22} & 0.0031 \\
core & \texttt{multiphase} & 3 & 0.0996 & 0.0481 & 1.000 & 1.000 & 0.000 & 0.0912 & 0.989 & \meth{log\_linear} & 0.0071 \\
core & \texttt{osc\_exp} & 3 & 0.0270 & 0.0158 & 0.578 & 1.000 & 0.000 & 0.0208 & 0.956 & \meth{richardson\_a10} & 0.0017 \\
core & \texttt{damped\_osc\_pow} & 3 & 0.0279 & 0.0033 & 1.000 & 1.000 & 0.000 & 0.0176 & 0.989 & \meth{richardson\_1} & 0.0009 \\
core & \texttt{log\_slow} & 3 & 0.0883 & 0.0446 & 1.000 & 1.000 & 0.000 & 0.0754 & 1.000 & \meth{log\_fit} & 0.0007 \\
core & \texttt{delayed\_plateau} & 3 & 0.1228 & 0.0283 & 1.000 & 1.000 & 0.000 & 0.0141 & 1.000 & \meth{double\_exp\_fit} & 0.0138 \\
core & \texttt{staircase} & 3 & 0.1300 & 0.0788 & 1.000 & 1.000 & 0.000 & 0.0608 & 1.000 & \meth{double\_exp\_fit} & 0.0136 \\
core & \texttt{noisy\_plateau} & 3 & 0.0025 & 0.0005 & 0.811 & 1.000 & 0.022 & 0.0006 & 0.833 & \meth{richardson\_a20} & 0.0004 \\
core & \texttt{slow\_osc\_power} & 3 & 0.0141 & 0.0514 & 0.078 & 1.000 & 0.100 & 0.0199 & 0.067 & \meth{richardson\_a20} & 0.0102 \\
core & \texttt{log\_oscillatory} & 3 & 0.1028 & 0.0506 & 0.922 & 1.000 & 0.000 & 0.0941 & 0.667 & \meth{log\_linear} & 0.0159 \\
core & \texttt{heavy\_noise\_plat} & 3 & 0.0029 & 0.0009 & 0.778 & 1.000 & 0.056 & 0.0010 & 0.800 & \meth{log\_fit} & 0.0008 \\
core & \texttt{broken\_power\_law} & 3 & 0.0658 & 0.0203 & 1.000 & 1.000 & 0.000 & 0.0203 & 1.000 & \meth{pade\_11} & 0.0091 \\
\addlinespace[2pt]
out-of-design & \texttt{stretched\_exp} & 3 & 0.0784 & 0.0048 & 1.000 & 1.000 & 0.000 & 0.0266 & 1.000 & \meth{pade\_11} & 0.0034 \\
out-of-design & \texttt{logistic\_tail} & 3 & 0.0498 & 0.1127 & 0.078 & 1.000 & 0.000 & 0.0364 & 0.989 & \meth{pade\_13} & 0.0072 \\
out-of-design & \texttt{inv\_sqrt\_log} & 3 & 0.0104 & 0.0013 & 0.922 & 1.000 & 0.000 & 0.0070 & 0.889 & \meth{richardson\_1} & 0.0009 \\
out-of-design & \texttt{random\_knots} & 3 & 0.0297 & 0.0048 & 0.978 & 1.000 & 0.000 & 0.0161 & 0.978 & \meth{richardson\_1} & 0.0043 \\
out-of-design & \texttt{real\_boot\_a} & 3 & 0.0834 & 0.0163 & 1.000 & 1.000 & 0.000 & 0.0427 & 1.000 & \meth{richardson\_3} & 0.0120 \\
out-of-design & \texttt{real\_boot\_b} & 3 & 0.0257 & 0.0067 & 1.000 & 1.000 & 0.000 & 0.0144 & 1.000 & \meth{richardson\_3} & 0.0033 \\
\midrule
\multicolumn{12}{l}{core, \meth{rational\_fit}: its win rate against the last value exceeds 0.5 in 16 of 18 uncapped families; not in \texttt{single\_exp} (0.111), \texttt{slow\_osc\_power} (0.078); its median error exceeds the last value's in \texttt{single\_exp} (0.0232 vs 0.0105), \texttt{slow\_osc\_power} (0.0514 vs 0.0141)} \\
\multicolumn{12}{l}{core, \meth{single\_exp\_fit}: its win rate against the last value exceeds 0.5 in 17 of 18 uncapped families; not in \texttt{slow\_osc\_power} (0.067); its median error exceeds the last value's in \texttt{slow\_osc\_power} (0.0199 vs 0.0141)} \\
\multicolumn{12}{l}{out-of-design, \meth{rational\_fit}: its win rate against the last value exceeds 0.5 in 5 of 6 uncapped families; not in \texttt{logistic\_tail} (0.078); its median error exceeds the last value's in \texttt{logistic\_tail} (0.1127 vs 0.0498)} \\
\multicolumn{12}{l}{out-of-design, \meth{single\_exp\_fit}: its win rate against the last value exceeds 0.5 in 6 of 6 uncapped families; its median error exceeds the last value's in no family} \\
\bottomrule
\end{tabular}

\end{adjustbox}
\end{table}
\end{landscape}
\begin{landscape}
\begin{table}[p]
\centering
\caption{\texttt{f23\_family\_head\_to\_head.tex}: Family head-to-head: curve families and cells where rational\_fit has the strictly lower median error than each other leading method and the last value, per stratum, core and out-of-design}
\label{tab:f23_family_head_to_head}
\begin{adjustbox}{max width=\linewidth, max totalheight=0.86\textheight}
% AUTO-GENERATED -- do not hand-edit.  generated by scripts/make_paper_tables.py from the run started 2026-10-01T22:16:39 on code d82fd26 (full; 29,256 s; 7 workers)
% Family head-to-head: curve families and cells where rational_fit has the strictly lower median error than each other leading method and the last value, per stratum, core and out-of-design
% source : results/phase1/phase1_aggregated.csv (per (method, regime, noise, g) cell; capped == 0)
% filter : capped == 0; comparators = every other member of LEAD (the union of the 10 lowest core median errors, the 5 lowest out-of-design median errors, the 3 highest core win rates against the last value (rank-eligible accelerators at g = 0.1) and the recorded-curve methods richardson_1, rational_fit; ordered by core rank, then out-of-design rank) and last_value; fam. = curve families where rational_fit's family median error (median over the family's cells of med_error) is strictly below the comparator's, k/n with n = families where both are finite; cells = the same over the (family, noise) cells with the cell med_error
\begin{tabular}{lrrrrrrrrrrrr}
\toprule
Comparator & \multicolumn{4}{c}{$g = 0.5$} & \multicolumn{4}{c}{$g = 0.1$} & \multicolumn{4}{c}{$g = 0.02$} \\
\cmidrule(lr){2-5}\cmidrule(lr){6-9}\cmidrule(lr){10-13}
 & core fam. & core cells & out-of-design fam. & out-of-design cells & core fam. & core cells & out-of-design fam. & out-of-design cells & core fam. & core cells & out-of-design fam. & out-of-design cells \\
\midrule
\meth{log\_linear} & 12/18 & 36/54 & 4/6 & 12/18 & 10/16 & 29/48 & 3/5 & 9/15 & 9/15 & 27/45 & 3/5 & 9/15 \\
\meth{richardson\_3} & 10/18 & 30/54 & 2/6 & 6/18 & 9/16 & 27/48 & 2/5 & 6/15 & 9/15 & 28/45 & 2/5 & 6/15 \\
\meth{richardson\_2} & 12/18 & 34/54 & 3/6 & 9/18 & 9/16 & 27/48 & 3/5 & 8/15 & 8/15 & 24/45 & 3/5 & 9/15 \\
\meth{double\_exp\_fit} & 10/18 & 28/54 & 5/6 & 14/18 & 7/16 & 20/48 & 4/5 & 11/15 & 8/15 & 22/45 & 4/5 & 11/15 \\
\meth{pade\_34} & 15/18 & 40/54 & 5/6 & 12/18 & 10/16 & 28/48 & 3/5 & 8/15 & 10/15 & 29/45 & 4/5 & 11/15 \\
\meth{richardson\_1} & 12/18 & 37/54 & 3/6 & 11/18 & 8/16 & 26/48 & 2/5 & 7/15 & 7/15 & 24/45 & 2/5 & 7/15 \\
\meth{pade\_23} & 14/18 & 38/54 & 5/6 & 11/18 & 11/16 & 31/48 & 3/5 & 8/15 & 11/15 & 33/45 & 4/5 & 11/15 \\
\meth{pade\_12} & 14/18 & 38/54 & 4/6 & 11/18 & 10/16 & 31/48 & 3/5 & 9/15 & 10/15 & 31/45 & 3/5 & 10/15 \\
\meth{pade\_45} & 15/18 & 39/54 & 4/6 & 10/18 & 11/16 & 32/48 & 3/5 & 9/15 & 11/15 & 32/45 & 4/5 & 11/15 \\
\meth{pade\_33} & 7/18 & 28/54 & 3/6 & 9/18 & 8/16 & 27/48 & 2/5 & 7/15 & 7/15 & 28/44 & 4/5 & 10/15 \\
\meth{pade\_22} & 10/18 & 30/54 & 3/6 & 9/18 & 10/16 & 26/46 & 3/5 & 8/15 & 10/15 & 25/43 & 4/5 & 9/15 \\
\meth{single\_exp\_fit} & 13/18 & 35/54 & 5/6 & 15/18 & 9/16 & 27/48 & 4/5 & 12/15 & 9/15 & 27/45 & 4/5 & 12/15 \\
\meth{last\_value} & 16/18 & 48/54 & 5/6 & 15/18 & 14/16 & 42/48 & 4/5 & 12/15 & 13/15 & 38/45 & 5/5 & 15/15 \\
\bottomrule
\end{tabular}

\end{adjustbox}
\end{table}
\end{landscape}
\begin{landscape}
\begin{table}[p]
\centering
\caption{\texttt{f24\_leading\_fits.tex}: The leading fits side by side at the headline stratum: core and out-of-design rank, median error, win rate vs the last value, catastrophe and validity rates, sensitivity to the assumed asymptote, and the recorded curves}
\label{tab:f24_leading_fits}
\begin{adjustbox}{max width=\linewidth, max totalheight=0.86\textheight}
% AUTO-GENERATED -- do not hand-edit.  generated by scripts/make_paper_tables.py from the run started 2026-10-01T22:16:39 on code d82fd26 (full; 29,256 s; 7 workers)
% The leading fits side by side at the headline stratum: core and held-out rank, median error, win rate vs the last value, catastrophe and validity rates, sensitivity to the assumed asymptote, and the recorded curves
% source : results/phase1/phase1_global.csv (core regimes, all three strata, capped cells excluded)
% source : results/phase1/phase1_global_holdout.csv (held-out regimes)
% source : results/phase5b/phase5b_sweep1_consumers.csv (Phase 5b sweep 1b; the assumed-asymptote sensitivity)
% source : results/real_data/real_data_roster_v2.csv (the recorded curves, every method under the benchmark protocol)
% roster  : results/real_data/real_data_roster_provenance.json (evaluation started 2026-10-03T13:02:45 on code 7518fad, 5040 rows; added after the full run)
% filter : target_g == 0.1; members of LEAD (the union of the 10 lowest core median errors, the 5 lowest held-out median errors, the 3 highest core win rates against the last value (rank-eligible accelerators at g = 0.1) and the recorded-curve methods richardson_1, rational_fit; ordered by core rank, then held-out rank); r = rank among the rank-eligible accelerators; L_hat sens. = max change of the core median error across the five assumed-asymptote modes as a percentage of the minimum (the sweep-1b invariance rows), 'indep.' for a method whose output does not depend on the assumed asymptote; recorded curves: wins vs the last value before and after the curve's minimum (k/n) and the median ratio of the method's error to the last value's over the valid pre-minimum cells
\begin{tabular}{lrrrrrrrrrrrrrr}
\toprule
Method & \multicolumn{5}{c}{core} & \multicolumn{5}{c}{held-out} & $\hat L$ sens. & \multicolumn{3}{c}{recorded curves} \\
\cmidrule(lr){2-6}\cmidrule(lr){7-11}\cmidrule(lr){13-15}
 & $r$ & med.\ err & win vs last & $\rhoC$ & $\rhoV$ & $r$ & med.\ err & win vs last & $\rhoC$ & $\rhoV$ & & wins pre & wins post & ratio pre \\
\midrule
\meth{log\_linear} & 1 & 0.0202 & 0.756 & 0.112 & 0.953 & 10 & 0.0305 & 0.847 & 0.000 & 1.000 & 362.8\% & 34/65 & 5/25 & 0.883 \\
\meth{richardson\_3} & 2 & 0.0222 & 0.774 & 0.097 & 0.980 & 8 & 0.0260 & 0.798 & 0.035 & 0.973 & 2.7\% & 45/65 & 5/25 & 0.310 \\
\meth{richardson\_2} & 3 & 0.0257 & 0.806 & 0.085 & 0.992 & 11 & 0.0316 & 0.820 & 0.013 & 0.987 & 0.9\% & 50/65 & 4/25 & 0.326 \\
\meth{double\_exp\_fit} & 4 & 0.0287 & 0.911 & 0.029 & 0.979 & 13 & 0.0322 & 0.962 & 0.002 & 1.000 & 29.7\% & 59/65 & 22/25 & 0.520 \\
\meth{pade\_34} & 5 & 0.0291 & 0.685 & 0.095 & 0.992 & 3 & 0.0211 & 0.722 & 0.087 & 0.991 & indep. & 18/65 & 6/25 & 1.787 \\
\meth{rational\_fit} & 6 & 0.0296 & 0.877 & 0.009 & 1.000 & 1 & 0.0074 & 0.853 & 0.002 & 1.000 & 0.0\% & 65/65 & 13/25 & 0.270 \\
\meth{richardson\_1} & 7 & 0.0301 & 0.811 & 0.088 & 1.000 & 12 & 0.0318 & 0.807 & 0.000 & 1.000 & 0.0\% & 63/65 & 10/25 & 0.268 \\
\meth{pade\_23} & 8 & 0.0302 & 0.701 & 0.093 & 0.996 & 6 & 0.0243 & 0.720 & 0.082 & 0.993 & indep. & 23/65 & 2/25 & 1.549 \\
\meth{pade\_12} & 9 & 0.0314 & 0.684 & 0.087 & 0.995 & 5 & 0.0225 & 0.729 & 0.084 & 0.991 & indep. & 30/65 & 3/25 & 1.030 \\
\meth{pade\_45} & 10 & 0.0318 & 0.700 & 0.085 & 0.994 & 7 & 0.0253 & 0.713 & 0.089 & 0.998 & indep. & 25/65 & 3/25 & 1.265 \\
\meth{pade\_33} & 11 & 0.0329 & 0.678 & 0.139 & 0.995 & 2 & 0.0153 & 0.858 & 0.080 & 0.960 & indep. & 65/65 & 15/25 & 0.282 \\
\meth{pade\_22} & 12 & 0.0333 & 0.657 & 0.163 & 0.924 & 4 & 0.0217 & 0.864 & 0.071 & 0.947 & indep. & 59/65 & 13/25 & 0.265 \\
\meth{single\_exp\_fit} & 13 & 0.0361 & 0.962 & 0.008 & 1.000 & 15 & 0.0342 & 0.973 & 0.002 & 1.000 & 1.7\% & 65/65 & 22/25 & 0.628 \\
\bottomrule
\end{tabular}

\end{adjustbox}
\end{table}
\end{landscape}
\begin{landscape}
\begin{table}[p]
\centering
\caption{\texttt{f25\_real\_roster.tex}: The recorded curves, every method: validity, wins against the last value, catastrophes and the median error ratio before and after the curve minimum, wins over all cells and the strict-validity count}
\label{tab:f25_real_roster}
\begin{adjustbox}{max width=\linewidth, max totalheight=0.86\textheight}
% AUTO-GENERATED -- do not hand-edit.  generated by scripts/make_paper_tables.py from the run started 2026-10-01T22:16:39 on code d82fd26 (full; 29,256 s; 7 workers)
% The recorded curves, every method: validity, wins against the last value, catastrophes and the median error ratio before and after the curve minimum, wins over all cells and the strict-validity count
% source : results/real_data/real_data_roster_v2.csv (the six recorded XGBoost curves on the (depth x target) grid; every accelerator and the four deployable trivial predictors under the benchmark's configuration and per-record scoring, src.evaluation)
% roster  : results/real_data/real_data_roster_provenance.json (evaluation started 2026-10-03T13:02:45 on code 7518fad, 5040 rows; added after the full run)
% filter : all rows; accelerators in the family order of f20, then the four trivial predictors as comparators; dagger = the excluded set; T = method type; before / after the minimum = post_min_target 0 / 1 (target round beyond the recorded curve's argmin); valid = benchmark validity (finite and in [-0.5, 500]), k/n; wins = cells with win_vs_last == 1 (an invalid record never wins); catastr. = invalid or error above 5 x the last value's; ratio = median over the valid cells of the error ratio to the last value's (skill_vs_last); strictly valid = the legacy rule (finite and in [0, 2 x window max]) over all cells
\begin{tabular}{lllrrrrrrrrrr}
\toprule
Method & family & T & \multicolumn{4}{c}{before the minimum} & \multicolumn{4}{c}{after the minimum} & \multicolumn{2}{c}{all cells} \\
\cmidrule(lr){4-7}\cmidrule(lr){8-11}\cmidrule(lr){12-13}
 & & & valid & wins & catastr. & ratio & valid & wins & catastr. & ratio & wins & strictly valid \\
\midrule
$\dagger$\meth{linear} & baseline & \TE & 44/65 & 1/65 & 45/65 & 6.156 & 17/25 & 0/25 & 25/25 & 16.212 & 1/90 & 47/90 \\
\meth{log\_linear} & baseline & \TE & 65/65 & 34/65 & 0/65 & 0.883 & 25/25 & 5/25 & 4/25 & 2.803 & 39/90 & 90/90 \\
$\dagger$\meth{geom\_avg\_diff} & baseline & \LE & 52/65 & 52/65 & 13/65 & 0.465 & 14/25 & 14/25 & 11/25 & 0.764 & 66/90 & 66/90 \\
\addlinespace[2pt]
\meth{richardson\_1} & Richardson & \TE & 65/65 & 63/65 & 0/65 & 0.268 & 25/25 & 10/25 & 2/25 & 1.312 & 73/90 & 90/90 \\
\meth{richardson\_2} & Richardson & \TE & 65/65 & 50/65 & 0/65 & 0.326 & 25/25 & 4/25 & 2/25 & 1.723 & 54/90 & 90/90 \\
\meth{richardson\_3} & Richardson & \TE & 56/65 & 45/65 & 9/65 & 0.310 & 25/25 & 5/25 & 2/25 & 1.414 & 50/90 & 81/90 \\
\meth{richardson\_a05} & Richardson & \TE & 65/65 & 54/65 & 0/65 & 0.847 & 25/25 & 11/25 & 3/25 & 1.487 & 65/90 & 90/90 \\
\meth{richardson\_a10} & Richardson & \TE & 65/65 & 39/65 & 0/65 & 0.515 & 25/25 & 7/25 & 1/25 & 1.469 & 46/90 & 90/90 \\
\meth{richardson\_a20} & Richardson & \TE & 65/65 & 39/65 & 0/65 & 0.877 & 25/25 & 13/25 & 1/25 & 0.833 & 52/90 & 90/90 \\
\addlinespace[2pt]
\meth{single\_exp\_fit} & parametric fit & \TE & 65/65 & 65/65 & 0/65 & 0.628 & 25/25 & 22/25 & 1/25 & 0.689 & 87/90 & 90/90 \\
\meth{double\_exp\_fit} & parametric fit & \TE & 65/65 & 59/65 & 1/65 & 0.520 & 25/25 & 22/25 & 1/25 & 0.532 & 81/90 & 90/90 \\
\meth{rational\_fit} & parametric fit & \TE & 65/65 & 65/65 & 0/65 & 0.270 & 25/25 & 13/25 & 2/25 & 0.783 & 78/90 & 90/90 \\
\meth{log\_fit} & parametric fit & \TE & 65/65 & 57/65 & 0/65 & 0.666 & 25/25 & 10/25 & 3/25 & 1.829 & 67/90 & 90/90 \\
\addlinespace[2pt]
\meth{shanks\_1} & Shanks & \LE & 65/65 & 29/65 & 0/65 & 1.049 & 25/25 & 9/25 & 0/25 & 1.080 & 38/90 & 90/90 \\
\meth{shanks\_2} & Shanks & \LE & 65/65 & 22/65 & 0/65 & 1.029 & 25/25 & 13/25 & 0/25 & 0.997 & 35/90 & 90/90 \\
\meth{shanks\_3} & Shanks & \LE & 65/65 & 16/65 & 0/65 & 1.039 & 25/25 & 3/25 & 0/25 & 1.101 & 19/90 & 90/90 \\
\meth{shanks\_4} & Shanks & \LE & 65/65 & 28/65 & 0/65 & 1.016 & 25/25 & 12/25 & 0/25 & 1.048 & 40/90 & 90/90 \\
\addlinespace[2pt]
\meth{wynn\_eps\_1} & Wynn $\varepsilon$ & \LE & 65/65 & 29/65 & 0/65 & 1.049 & 25/25 & 9/25 & 0/25 & 1.080 & 38/90 & 90/90 \\
\meth{wynn\_eps\_2} & Wynn $\varepsilon$ & \LE & 62/65 & 31/65 & 3/65 & 1.006 & 25/25 & 15/25 & 0/25 & 0.997 & 46/90 & 87/90 \\
\meth{wynn\_eps\_3} & Wynn $\varepsilon$ & \LE & 65/65 & 38/65 & 4/65 & 0.887 & 25/25 & 15/25 & 2/25 & 0.987 & 53/90 & 90/90 \\
\addlinespace[2pt]
\meth{wynn\_rho\_1} & Wynn $\rho$ & \LE & 65/65 & 29/65 & 3/65 & 1.081 & 25/25 & 8/25 & 0/25 & 1.132 & 37/90 & 87/90 \\
\meth{wynn\_rho\_2} & Wynn $\rho$ & \LE & 65/65 & 21/65 & 8/65 & 1.202 & 25/25 & 13/25 & 2/25 & 0.992 & 34/90 & 87/90 \\
\meth{wynn\_rho\_3} & Wynn $\rho$ & \LE & 62/65 & 33/65 & 13/65 & 0.996 & 25/25 & 11/25 & 3/25 & 1.186 & 44/90 & 84/90 \\
\addlinespace[2pt]
\meth{pade\_11} & Pad\'{e} & \TE & 65/65 & 65/65 & 0/65 & 0.235 & 25/25 & 16/25 & 2/25 & 0.818 & 81/90 & 90/90 \\
\meth{pade\_12} & Pad\'{e} & \TE & 64/65 & 30/65 & 10/65 & 1.030 & 23/25 & 3/25 & 12/25 & 4.305 & 33/90 & 82/90 \\
\meth{pade\_13} & Pad\'{e} & \TE & 62/65 & 0/65 & 36/65 & 5.743 & 24/25 & 0/25 & 22/25 & 17.143 & 0/90 & 41/90 \\
$\dagger$\meth{pade\_21} & Pad\'{e} & \TE & 64/65 & 30/65 & 4/65 & 1.231 & 25/25 & 2/25 & 13/25 & 5.280 & 32/90 & 88/90 \\
\meth{pade\_22} & Pad\'{e} & \TE & 65/65 & 59/65 & 1/65 & 0.265 & 25/25 & 13/25 & 2/25 & 0.712 & 72/90 & 90/90 \\
\meth{pade\_23} & Pad\'{e} & \TE & 65/65 & 23/65 & 13/65 & 1.549 & 25/25 & 2/25 & 11/25 & 3.690 & 25/90 & 84/90 \\
$\dagger$\meth{pade\_31} & Pad\'{e} & \TE & 55/65 & 1/65 & 54/65 & 13.640 & 17/25 & 0/25 & 25/25 & 41.170 & 1/90 & 37/90 \\
$\dagger$\meth{pade\_32} & Pad\'{e} & \TE & 63/65 & 23/65 & 7/65 & 1.535 & 25/25 & 3/25 & 10/25 & 3.235 & 26/90 & 85/90 \\
\meth{pade\_33} & Pad\'{e} & \TE & 65/65 & 65/65 & 0/65 & 0.282 & 25/25 & 15/25 & 2/25 & 0.762 & 80/90 & 90/90 \\
\meth{pade\_34} & Pad\'{e} & \TE & 63/65 & 18/65 & 14/65 & 1.787 & 23/25 & 6/25 & 9/25 & 1.986 & 24/90 & 78/90 \\
\meth{pade\_44} & Pad\'{e} & \TE & 65/65 & 60/65 & 0/65 & 0.280 & 25/25 & 15/25 & 1/25 & 0.774 & 75/90 & 90/90 \\
\meth{pade\_45} & Pad\'{e} & \TE & 63/65 & 25/65 & 16/65 & 1.265 & 24/25 & 3/25 & 11/25 & 3.317 & 28/90 & 77/90 \\
\addlinespace[2pt]
\meth{levin\_t1} & Levin & \LE & 65/65 & 29/65 & 0/65 & 1.049 & 25/25 & 9/25 & 0/25 & 1.080 & 38/90 & 90/90 \\
\meth{levin\_t2} & Levin & \LE & 65/65 & 0/65 & 0/65 & 1.026 & 25/25 & 5/25 & 0/25 & 1.041 & 5/90 & 90/90 \\
\meth{levin\_u1} & Levin & \LE & 65/65 & 26/65 & 0/65 & 1.056 & 25/25 & 9/25 & 0/25 & 1.079 & 35/90 & 90/90 \\
\meth{levin\_u2} & Levin & \LE & 65/65 & 0/65 & 0/65 & 1.027 & 25/25 & 5/25 & 0/25 & 1.039 & 5/90 & 90/90 \\
\meth{levin\_v1} & Levin & \LE & 65/65 & 0/65 & 0/65 & 1.045 & 25/25 & 4/25 & 0/25 & 1.065 & 4/90 & 90/90 \\
\meth{levin\_v2} & Levin & \LE & 65/65 & 3/65 & 0/65 & 1.060 & 25/25 & 0/25 & 0/25 & 1.078 & 3/90 & 90/90 \\
\addlinespace[2pt]
\meth{brezinski\_theta1} & Brezinski $\theta$ & \LE & 65/65 & 29/65 & 0/65 & 1.040 & 25/25 & 9/25 & 0/25 & 1.066 & 38/90 & 90/90 \\
\meth{brezinski\_theta2} & Brezinski $\theta$ & \LE & 65/65 & 29/65 & 3/65 & 1.059 & 25/25 & 15/25 & 1/25 & 0.904 & 44/90 & 90/90 \\
\addlinespace[2pt]
$\dagger$\meth{neville\_2} & Neville & \TE & 65/65 & 18/65 & 20/65 & 2.278 & 25/25 & 0/25 & 20/25 & 12.238 & 18/90 & 70/90 \\
$\dagger$\meth{neville\_3} & Neville & \TE & 32/65 & 0/65 & 65/65 & 44.120 & 10/25 & 0/25 & 25/25 & 255.031 & 0/90 & 3/90 \\
$\dagger$\meth{neville\_4} & Neville & \TE & 28/65 & 0/65 & 65/65 & 353.112 & 5/25 & 0/25 & 25/25 & 26100.700 & 0/90 & 0/90 \\
\addlinespace[2pt]
\meth{anderson\_1} & Anderson & \LE & 65/65 & 32/65 & 0/65 & 1.000 & 25/25 & 9/25 & 0/25 & 1.030 & 41/90 & 90/90 \\
\meth{anderson\_2} & Anderson & \LE & 65/65 & 33/65 & 0/65 & 0.998 & 25/25 & 14/25 & 0/25 & 0.999 & 47/90 & 90/90 \\
\meth{anderson\_3} & Anderson & \LE & 65/65 & 27/65 & 0/65 & 1.002 & 25/25 & 11/25 & 0/25 & 1.001 & 38/90 & 90/90 \\
\addlinespace[2pt]
\meth{median\_ensemble} & ensemble & \TE & 65/65 & 33/65 & 0/65 & 0.991 & 25/25 & 9/25 & 0/25 & 1.053 & 42/90 & 90/90 \\
\meth{stability\_weighted} & ensemble & \TE & 65/65 & 32/65 & 2/65 & 1.000 & 25/25 & 4/25 & 2/25 & 1.110 & 36/90 & 90/90 \\
\meth{best\_shanks\_wynn} & ensemble & \LE & 65/65 & 29/65 & 0/65 & 1.049 & 25/25 & 9/25 & 0/25 & 1.080 & 38/90 & 90/90 \\
\midrule
\meth{constant\_assumed} & trivial & \LE & 65/65 & 0/65 & 43/65 & 9.119 & 25/25 & 0/25 & 21/25 & 20.945 & 0/90 & 90/90 \\
\meth{last\_value} & trivial & \LE & 65/65 & 0/65 & 0/65 & 1.000 & 25/25 & 0/25 & 0/25 & 1.000 & 0/90 & 90/90 \\
\meth{window\_mean} & trivial & \LE & 65/65 & 0/65 & 0/65 & 1.710 & 25/25 & 0/25 & 1/25 & 2.190 & 0/90 & 90/90 \\
\meth{window\_min} & trivial & \LE & 65/65 & 0/65 & 0/65 & 1.000 & 25/25 & 2/25 & 0/25 & 1.000 & 2/90 & 90/90 \\
\bottomrule
\end{tabular}

\end{adjustbox}
\end{table}
\end{landscape}
\begin{landscape}
\begin{table}[p]
\centering
\caption{\texttt{f26\_real\_classical.tex}: The classical variants on the recorded curves: median improvement factor, wins against the last value and validity before and after the curve minimum and over all cells}
\label{tab:f26_real_classical}
\begin{adjustbox}{max width=\linewidth, max totalheight=0.86\textheight}
% AUTO-GENERATED -- do not hand-edit.  generated by scripts/make_paper_tables.py from the run started 2026-10-01T22:16:39 on code d82fd26 (full; 29,256 s; 7 workers)
% The classical variants on the recorded curves: median improvement factor, wins against the last value and validity before and after the curve minimum and over all cells
% source : results/real_data/real_data_roster_v2.csv (the recorded curves, every method under the benchmark protocol)
% roster  : results/real_data/real_data_roster_provenance.json (evaluation started 2026-10-03T13:02:45 on code 7518fad, 5040 rows; added after the full run)
% filter : family in ('shanks', 'wynn_eps', 'wynn_rho', 'levin', 'brezinski', 'anderson') (the 21 classical variants); pre-minimum / post-minimum / all cells (post_min_target 0 / 1 / any); MI = median over the valid rows of improve_ratio (last-value error / method error), bold = inside [0.9, 1.1]; wins = cells with win_vs_last == 1 (k/n); valid = benchmark validity (k/n); footers: variants in the MI band, variants with a win rate inside the coin-flip band, the mean win rate, and the in-band count per dataset over the pre-minimum cells
\begin{tabular}{llrrrrrrrrr}
\toprule
Family & Method & \multicolumn{3}{c}{pre-minimum cells} & \multicolumn{3}{c}{post-minimum cells} & \multicolumn{3}{c}{all cells} \\
\cmidrule(lr){3-5}\cmidrule(lr){6-8}\cmidrule(lr){9-11}
 & & MI & wins & valid & MI & wins & valid & MI & wins & valid \\
\midrule
Shanks & \meth{shanks\_1} & \textbf{0.953} & 29/65 & 65/65 & \textbf{0.926} & 9/25 & 25/25 & \textbf{0.948} & 38/90 & 90/90 \\
Shanks & \meth{shanks\_2} & \textbf{0.972} & 22/65 & 65/65 & \textbf{1.003} & 13/25 & 25/25 & \textbf{0.973} & 35/90 & 90/90 \\
Shanks & \meth{shanks\_3} & \textbf{0.962} & 16/65 & 65/65 & \textbf{0.908} & 3/25 & 25/25 & \textbf{0.950} & 19/90 & 90/90 \\
Shanks & \meth{shanks\_4} & \textbf{0.984} & 28/65 & 65/65 & \textbf{0.954} & 12/25 & 25/25 & \textbf{0.982} & 40/90 & 90/90 \\
\addlinespace[2pt]
Wynn $\varepsilon$ & \meth{wynn\_eps\_1} & \textbf{0.953} & 29/65 & 65/65 & \textbf{0.926} & 9/25 & 25/25 & \textbf{0.948} & 38/90 & 90/90 \\
Wynn $\varepsilon$ & \meth{wynn\_eps\_2} & \textbf{0.994} & 31/65 & 62/65 & \textbf{1.003} & 15/25 & 25/25 & \textbf{1.003} & 46/90 & 87/90 \\
Wynn $\varepsilon$ & \meth{wynn\_eps\_3} & 1.128 & 38/65 & 65/65 & \textbf{1.013} & 15/25 & 25/25 & 1.110 & 53/90 & 90/90 \\
\addlinespace[2pt]
Wynn $\rho$ & \meth{wynn\_rho\_1} & \textbf{0.925} & 29/65 & 65/65 & 0.884 & 8/25 & 25/25 & \textbf{0.920} & 37/90 & 90/90 \\
Wynn $\rho$ & \meth{wynn\_rho\_2} & 0.832 & 21/65 & 65/65 & \textbf{1.008} & 13/25 & 25/25 & \textbf{0.910} & 34/90 & 90/90 \\
Wynn $\rho$ & \meth{wynn\_rho\_3} & \textbf{1.004} & 33/65 & 62/65 & 0.843 & 11/25 & 25/25 & \textbf{1.002} & 44/90 & 87/90 \\
\addlinespace[2pt]
Levin & \meth{levin\_t1} & \textbf{0.953} & 29/65 & 65/65 & \textbf{0.926} & 9/25 & 25/25 & \textbf{0.948} & 38/90 & 90/90 \\
Levin & \meth{levin\_t2} & \textbf{0.975} & 0/65 & 65/65 & \textbf{0.960} & 5/25 & 25/25 & \textbf{0.971} & 5/90 & 90/90 \\
Levin & \meth{levin\_u1} & \textbf{0.947} & 26/65 & 65/65 & \textbf{0.927} & 9/25 & 25/25 & \textbf{0.943} & 35/90 & 90/90 \\
Levin & \meth{levin\_u2} & \textbf{0.974} & 0/65 & 65/65 & \textbf{0.962} & 5/25 & 25/25 & \textbf{0.967} & 5/90 & 90/90 \\
Levin & \meth{levin\_v1} & \textbf{0.957} & 0/65 & 65/65 & \textbf{0.939} & 4/25 & 25/25 & \textbf{0.952} & 4/90 & 90/90 \\
Levin & \meth{levin\_v2} & \textbf{0.943} & 3/65 & 65/65 & \textbf{0.927} & 0/25 & 25/25 & \textbf{0.938} & 3/90 & 90/90 \\
\addlinespace[2pt]
Brezinski $\theta$ & \meth{brezinski\_theta1} & \textbf{0.962} & 29/65 & 65/65 & \textbf{0.938} & 9/25 & 25/25 & \textbf{0.958} & 38/90 & 90/90 \\
Brezinski $\theta$ & \meth{brezinski\_theta2} & \textbf{0.944} & 29/65 & 65/65 & 1.106 & 15/25 & 25/25 & \textbf{0.974} & 44/90 & 90/90 \\
\addlinespace[2pt]
Anderson & \meth{anderson\_1} & \textbf{1.000} & 32/65 & 65/65 & \textbf{0.971} & 9/25 & 25/25 & \textbf{0.999} & 41/90 & 90/90 \\
Anderson & \meth{anderson\_2} & \textbf{1.002} & 33/65 & 65/65 & \textbf{1.001} & 14/25 & 25/25 & \textbf{1.001} & 47/90 & 90/90 \\
Anderson & \meth{anderson\_3} & \textbf{0.998} & 27/65 & 65/65 & \textbf{0.999} & 11/25 & 25/25 & \textbf{0.998} & 38/90 & 90/90 \\
\midrule
\multicolumn{2}{l}{in the $\pm 10\%$ band (MI in $[0.9, 1.1]$)} & 19/21 & &  & 18/21 & &  & 20/21 & &  \\
\multicolumn{2}{l}{win rate vs last in $[0.4, 0.6]$ (coin flip)} &  & 14/21 &  &  & 9/21 &  &  & 13/21 &  \\
\multicolumn{2}{l}{mean win rate vs last} &  & 0.355 &  &  & 0.377 &  &  & 0.361 &  \\
\multicolumn{11}{l}{pre-minimum cells, variants in the band per dataset: \texttt{adult} 15/21 (5 cells); \texttt{covertype} 19/21 (15 cells); \texttt{higgs} 20/21 (15 cells); \texttt{jannis} 11/21 (15 cells); \texttt{miniboone} 11/21 (15 cells)} \\
\bottomrule
\end{tabular}

\end{adjustbox}
\end{table}
\end{landscape}
\begin{landscape}
\begin{table}[p]
\centering
\caption{\texttt{f27\_selection\_one\_pilot.tex}: Choosing a method by trial on the synthetic families, design one pilot (one\_pilot), headline stratum: per set, noise class and candidate pool, how often the method chosen on the pilot run(s) is valid on the final run, its catastrophe rate there and the default's, how often it beats the default and the last value, and the median errors of the chosen method, the default and the last value}
\label{tab:f27_selection_one_pilot}
\begin{adjustbox}{max width=\linewidth, max totalheight=0.86\textheight}
% AUTO-GENERATED -- do not hand-edit.  generated by scripts/make_paper_tables.py from the run started 2026-10-01T22:16:39 on code d82fd26 (full; 29,256 s; 7 workers)
% Choosing a method by trial on the synthetic families, design one pilot (one_pilot), headline stratum: per set, noise class and candidate pool, how often the method chosen on the pilot run(s) is valid on the final run, its catastrophe rate there and the default's, how often it beats the default and the last value, and the median errors of the chosen method, the default and the last value
% source : results/phase1/phase1_selection_global.csv (scripts/derive_selection.py on the git-ignored phase1_records.csv: the (family, noise, g) cells with capped == 0, all seeds; one row per set, stratum, noise class, candidate pool and design)
% selection: results/phase1/phase1_selection_provenance.json (analysis started 2026-10-03T18:14:05 on code b4d5518 from 369,360 records; cells per design {'one_pilot': 195, 'many_pilots': 195}, trials per design {'one_pilot': 169650, 'many_pilots': 5850}; added after the full run, evaluates no method)
% filter : design == one_pilot (every ordered pair of distinct seeds (pilot, final); chosen = the pool member with the lowest error on the pilot seed among those valid there (ties: registry order); its record on the final seed is taken as it is (an invalid record stays invalid, no fallback)); target_g == 0.1; noise class of a (family, noise) cell: noise-free = sigma 0 and no intrinsic noise; intrinsic noise only = sigma 0 and one of the intrinsic-noise families (noisy_plateau, heavy_noise_plat); otherwise sigma = the level; 'all cells' pools the classes; candidate pools: the leading fits = LEAD (the union of the 10 lowest core median errors, the 5 lowest out-of-design median errors, the 3 highest core win rates against the last value (rank-eligible accelerators at g = 0.1) and the recorded-curve methods richardson_1, rational_fit; ordered by core rank, then out-of-design rank: log_linear, richardson_3, richardson_2, double_exp_fit, pade_34, rational_fit, richardson_1, pade_23, pade_12, pade_45, pade_33, pade_22, single_exp_fit), the two named methods = rational_fit (the default) and single_exp_fit, the classical variants = the 21 members of the families ('shanks', 'wynn_eps', 'wynn_rho', 'levin', 'brezinski', 'anderson'), all accelerators = the 52 accelerators of the registry; a trial = one (pilot, final) choice on one cell; chosen valid = mean over the class's cells of the share of trials whose chosen record is valid on the final run; win vs default (vs last) = mean over cells of the share of trials where the chosen record is valid and its error is strictly below rational_fit's (the last value's) on the same final run (an invalid chosen record never wins; a valid one wins against an invalid reference); default chosen = mean over cells of the share of trials where the chosen method is rational_fit itself (never a strict win); families with win rate > 0.5 = the families whose mean over their cells of the win share exceeds 0.5, k/n; median error = median over the class's cells of the cell median over the trials (chosen: over the trials where it is valid; default and last value: over the final runs of the trials, so the 'all cells' values reproduce the pooled Phase 1 medians); ratio to default = median over cells of the cell median of chosen error / rational_fit error over the trials where both are valid; catastrophe rate (chosen, default) = mean over the cells of the share of trials whose chosen record (rational_fit's record) on the final run is catastrophic, the Phase 1 flag as recorded (the estimate is invalid, or its error exceeds CAT_MULT = 5 times the last value's on the same run); a trial without a choice counts as catastrophic, as it counts as invalid
\begin{tabular}{lllrrrrrrrrrrrrr}
\toprule
Set & noise class & candidate pool & families & chosen valid & \multicolumn{2}{c}{catastrophe rate} & win vs default & default chosen & win vs last & \multicolumn{2}{c}{families with win rate $> 0.5$} & \multicolumn{3}{c}{median error} & ratio to default \\
\cmidrule(lr){6-7}\cmidrule(lr){11-12}\cmidrule(lr){13-15}
 & & & & & chosen & default & & & & vs default & vs last & chosen & default & last value & \\
\midrule
core & noise-free & the leading fits (13) & 14 & 0.997 & 0.008 & 0.000 & 0.828 & 0.024 & 0.970 & 12/14 & 14/14 & 0.0057 & 0.0334 & 0.1255 & 0.148 \\
 &  & the two named methods & 14 & 1.000 & 0.000 & 0.000 & 0.306 & 0.631 & 0.998 & 4/14 & 14/14 & 0.0143 & 0.0334 & 0.1255 & 1.000 \\
 &  & the classical variants (21) & 14 & 1.000 & 0.000 & 0.000 & 0.802 & 0.000 & 0.929 & 12/14 & 13/14 & 0.0066 & 0.0334 & 0.1255 & 0.194 \\
 &  & all accelerators (52) & 14 & 1.000 & 0.006 & 0.000 & 0.880 & 0.017 & 0.975 & 13/14 & 14/14 & 0.0003 & 0.0334 & 0.1255 & 0.017 \\
\addlinespace[2pt]
 & intrinsic noise only & the leading fits (13) & 2 & 0.994 & 0.109 & 0.067 & 0.386 & 0.133 & 0.652 & 0/2 & 2/2 & 0.0011 & 0.0009 & 0.0023 & 1.005 \\
 &  & the two named methods & 2 & 1.000 & 0.047 & 0.067 & 0.356 & 0.400 & 0.767 & 0/2 & 2/2 & 0.0008 & 0.0009 & 0.0023 & 1.000 \\
 &  & the classical variants (21) & 2 & 0.998 & 0.153 & 0.067 & 0.289 & 0.000 & 0.487 & 0/2 & 1/2 & 0.0023 & 0.0009 & 0.0023 & 2.057 \\
 &  & all accelerators (52) & 2 & 0.994 & 0.119 & 0.067 & 0.416 & 0.017 & 0.619 & 0/2 & 2/2 & 0.0013 & 0.0009 & 0.0023 & 1.178 \\
\addlinespace[2pt]
 & $\sigma = 0.001$ & the leading fits (13) & 16 & 0.995 & 0.026 & 0.010 & 0.664 & 0.065 & 0.922 & 11/16 & 16/16 & 0.0073 & 0.0303 & 0.1038 & 0.621 \\
 &  & the two named methods & 16 & 1.000 & 0.010 & 0.010 & 0.297 & 0.613 & 0.967 & 4/16 & 16/16 & 0.0134 & 0.0303 & 0.1038 & 1.000 \\
 &  & the classical variants (21) & 16 & 0.997 & 0.024 & 0.010 & 0.126 & 0.000 & 0.383 & 2/16 & 2/16 & 0.1059 & 0.0303 & 0.1038 & 3.811 \\
 &  & all accelerators (52) & 16 & 0.997 & 0.024 & 0.010 & 0.640 & 0.035 & 0.903 & 11/16 & 16/16 & 0.0080 & 0.0303 & 0.1038 & 0.742 \\
\addlinespace[2pt]
 & $\sigma = 0.005$ & the leading fits (13) & 16 & 0.995 & 0.025 & 0.008 & 0.523 & 0.125 & 0.882 & 8/16 & 15/16 & 0.0140 & 0.0296 & 0.1022 & 0.948 \\
 &  & the two named methods & 16 & 1.000 & 0.008 & 0.008 & 0.279 & 0.608 & 0.944 & 3/16 & 16/16 & 0.0152 & 0.0296 & 0.1022 & 1.000 \\
 &  & the classical variants (21) & 16 & 0.998 & 0.023 & 0.008 & 0.131 & 0.000 & 0.368 & 2/16 & 1/16 & 0.1053 & 0.0296 & 0.1022 & 3.836 \\
 &  & all accelerators (52) & 16 & 0.992 & 0.028 & 0.008 & 0.503 & 0.079 & 0.846 & 8/16 & 15/16 & 0.0189 & 0.0296 & 0.1022 & 0.986 \\
\addlinespace[2pt]
 & all cells & the leading fits (13) & 16 & 0.996 & 0.024 & 0.009 & 0.653 & 0.076 & 0.911 & 12/16 & 16/16 & 0.0104 & 0.0296 & 0.1030 & 0.619 \\
 &  & the two named methods & 16 & 1.000 & 0.008 & 0.009 & 0.296 & 0.608 & 0.960 & 4/16 & 16/16 & 0.0139 & 0.0296 & 0.1030 & 1.000 \\
 &  & the classical variants (21) & 16 & 0.998 & 0.022 & 0.009 & 0.332 & 0.000 & 0.542 & 2/16 & 14/16 & 0.0311 & 0.0296 & 0.1030 & 2.014 \\
 &  & all accelerators (52) & 16 & 0.996 & 0.024 & 0.009 & 0.655 & 0.044 & 0.893 & 12/16 & 16/16 & 0.0071 & 0.0296 & 0.1030 & 0.492 \\
\addlinespace[2pt]
\midrule
out-of-design & noise-free & the leading fits (13) & 5 & 0.997 & 0.005 & 0.000 & 0.816 & 0.020 & 0.983 & 5/5 & 5/5 & 0.0007 & 0.0053 & 0.0862 & 0.143 \\
 &  & the two named methods & 5 & 1.000 & 0.000 & 0.000 & 0.149 & 0.827 & 0.975 & 1/5 & 5/5 & 0.0053 & 0.0053 & 0.0862 & 1.000 \\
 &  & the classical variants (21) & 5 & 1.000 & 0.000 & 0.000 & 0.847 & 0.000 & 1.000 & 4/5 & 5/5 & 0.0014 & 0.0053 & 0.0862 & 0.145 \\
 &  & all accelerators (52) & 5 & 0.997 & 0.005 & 0.000 & 0.913 & 0.000 & 0.987 & 5/5 & 5/5 & 0.0004 & 0.0053 & 0.0862 & 0.075 \\
\addlinespace[2pt]
 & $\sigma = 0.001$ & the leading fits (13) & 5 & 0.998 & 0.008 & 0.000 & 0.674 & 0.040 & 0.971 & 4/5 & 5/5 & 0.0083 & 0.0074 & 0.0862 & 0.666 \\
 &  & the two named methods & 5 & 1.000 & 0.000 & 0.000 & 0.149 & 0.827 & 0.975 & 1/5 & 5/5 & 0.0074 & 0.0074 & 0.0862 & 1.000 \\
 &  & the classical variants (21) & 5 & 0.997 & 0.021 & 0.000 & 0.100 & 0.000 & 0.405 & 0/5 & 0/5 & 0.0971 & 0.0074 & 0.0862 & 4.527 \\
 &  & all accelerators (52) & 5 & 0.995 & 0.013 & 0.000 & 0.720 & 0.020 & 0.950 & 5/5 & 5/5 & 0.0054 & 0.0074 & 0.0862 & 0.734 \\
\addlinespace[2pt]
 & $\sigma = 0.005$ & the leading fits (13) & 5 & 1.000 & 0.006 & 0.007 & 0.468 & 0.180 & 0.931 & 2/5 & 5/5 & 0.0204 & 0.0102 & 0.0846 & 1.000 \\
 &  & the two named methods & 5 & 1.000 & 0.007 & 0.007 & 0.154 & 0.787 & 0.948 & 1/5 & 5/5 & 0.0102 & 0.0102 & 0.0846 & 1.000 \\
 &  & the classical variants (21) & 5 & 0.999 & 0.007 & 0.007 & 0.137 & 0.000 & 0.359 & 1/5 & 0/5 & 0.0969 & 0.0102 & 0.0846 & 3.269 \\
 &  & all accelerators (52) & 5 & 0.999 & 0.006 & 0.007 & 0.517 & 0.107 & 0.910 & 3/5 & 5/5 & 0.0096 & 0.0102 & 0.0846 & 0.672 \\
\addlinespace[2pt]
 & all cells & the leading fits (13) & 5 & 0.998 & 0.007 & 0.002 & 0.653 & 0.080 & 0.962 & 4/5 & 5/5 & 0.0073 & 0.0074 & 0.0862 & 0.666 \\
 &  & the two named methods & 5 & 1.000 & 0.002 & 0.002 & 0.151 & 0.813 & 0.966 & 1/5 & 5/5 & 0.0074 & 0.0074 & 0.0862 & 1.000 \\
 &  & the classical variants (21) & 5 & 0.999 & 0.009 & 0.002 & 0.361 & 0.000 & 0.588 & 1/5 & 5/5 & 0.0465 & 0.0074 & 0.0862 & 3.081 \\
 &  & all accelerators (52) & 5 & 0.997 & 0.008 & 0.002 & 0.716 & 0.042 & 0.949 & 5/5 & 5/5 & 0.0054 & 0.0074 & 0.0862 & 0.460 \\
\bottomrule
\end{tabular}

\end{adjustbox}
\end{table}
\end{landscape}
\begin{landscape}
\begin{table}[p]
\centering
\caption{\texttt{f27b\_selection\_many\_pilots.tex}: Choosing a method by trial on the synthetic families, design 29 pilots (many\_pilots), headline stratum: per set, noise class and candidate pool, how often the method chosen on the pilot run(s) is valid on the final run, its catastrophe rate there and the default's, how often it beats the default and the last value, and the median errors of the chosen method, the default and the last value}
\label{tab:f27b_selection_many_pilots}
\begin{adjustbox}{max width=\linewidth, max totalheight=0.86\textheight}
% AUTO-GENERATED -- do not hand-edit.  generated by scripts/make_paper_tables.py from the run started 2026-10-01T22:16:39 on code d82fd26 (full; 29,256 s; 7 workers)
% Choosing a method by trial on the synthetic families, design many pilot runs (many_pilots), headline stratum: per set, noise class and candidate pool, how often the method chosen on the pilot run(s) is valid on the final run, beats the default and the last value there, and the median errors of the chosen method, the default and the last value
% source : results/phase1/phase1_selection_global.csv (scripts/derive_selection.py on the git-ignored phase1_records.csv: the (family, noise, g) cells with capped == 0, all seeds; one row per set, stratum, noise class, candidate pool and design)
% selection: results/phase1/phase1_selection_provenance.json (analysis started 2026-10-03T15:58:20 on code 08d5116 from 369,360 records; cells per design {'one_pilot': 195, 'many_pilots': 195}, trials per design {'one_pilot': 169650, 'many_pilots': 5850}; added after the full run, evaluates no method)
% filter : design == many_pilots (every seed in turn is the final run and the other seeds are the pilots; chosen = the pool member with the lowest median error over its valid pilot records among the members valid on at least RANK_MIN_VALID of the pilot seeds (0.9; ties: registry order)); target_g == 0.1; noise class of a (family, noise) cell: noise-free = sigma 0 and no intrinsic noise; intrinsic noise only = sigma 0 and one of the intrinsic-noise families (noisy_plateau, heavy_noise_plat); otherwise sigma = the level; 'all cells' pools the classes; candidate pools: the leading fits = LEAD (the union of the 10 lowest core median errors, the 5 lowest held-out median errors, the 3 highest core win rates against the last value (rank-eligible accelerators at g = 0.1) and the recorded-curve methods richardson_1, rational_fit; ordered by core rank, then held-out rank: log_linear, richardson_3, richardson_2, double_exp_fit, pade_34, rational_fit, richardson_1, pade_23, pade_12, pade_45, pade_33, pade_22, single_exp_fit), the two named methods = rational_fit (the default) and single_exp_fit, the classical variants = the 21 members of the families ('shanks', 'wynn_eps', 'wynn_rho', 'levin', 'brezinski', 'anderson'), all accelerators = the 52 accelerators of the registry; a trial = one (pilot, final) choice on one cell; chosen valid = mean over the class's cells of the share of trials whose chosen record is valid on the final run; win vs default (vs last) = mean over cells of the share of trials where the chosen record is valid and its error is strictly below rational_fit's (the last value's) on the same final run (an invalid chosen record never wins; a valid one wins against an invalid reference); default chosen = mean over cells of the share of trials where the chosen method is rational_fit itself (never a strict win); families with win rate > 0.5 = the families whose mean over their cells of the win share exceeds 0.5, k/n; median error = median over the class's cells of the cell median over the trials (chosen: over the trials where it is valid; default and last value: over the final runs of the trials, so the 'all cells' values reproduce the pooled Phase 1 medians); ratio to default = median over cells of the cell median of chosen error / rational_fit error over the trials where both are valid
\begin{tabular}{lllrrrrrrrrrrr}
\toprule
Set & noise class & candidate pool & families & chosen valid & win vs default & default chosen & win vs last & \multicolumn{2}{c}{families with win rate $> 0.5$} & \multicolumn{3}{c}{median error} & ratio to default \\
\cmidrule(lr){9-10}\cmidrule(lr){11-13}
 & & & & & & & & vs default & vs last & chosen & default & last value & \\
\midrule
core & noise-free & the leading fits (13) & 14 & 1.000 & 0.905 & 0.000 & 1.000 & 14/14 & 14/14 & 0.0027 & 0.0334 & 0.1255 & 0.105 \\
 &  & the two named methods & 14 & 1.000 & 0.314 & 0.643 & 1.000 & 5/14 & 14/14 & 0.0143 & 0.0334 & 0.1255 & 1.000 \\
 &  & the classical variants (21) & 14 & 1.000 & 0.802 & 0.000 & 0.929 & 12/14 & 13/14 & 0.0066 & 0.0334 & 0.1255 & 0.194 \\
 &  & all accelerators (52) & 14 & 1.000 & 0.950 & 0.000 & 1.000 & 14/14 & 14/14 & 0.0003 & 0.0334 & 0.1255 & 0.017 \\
\addlinespace[2pt]
 & intrinsic noise only & the leading fits (13) & 2 & 1.000 & 0.583 & 0.033 & 0.800 & 2/2 & 2/2 & 0.0008 & 0.0009 & 0.0023 & 0.758 \\
 &  & the two named methods & 2 & 1.000 & 0.400 & 0.250 & 0.750 & 1/2 & 2/2 & 0.0009 & 0.0009 & 0.0023 & 0.884 \\
 &  & the classical variants (21) & 2 & 1.000 & 0.383 & 0.000 & 0.583 & 0/2 & 2/2 & 0.0016 & 0.0009 & 0.0023 & 1.592 \\
 &  & all accelerators (52) & 2 & 1.000 & 0.650 & 0.000 & 0.817 & 2/2 & 2/2 & 0.0006 & 0.0009 & 0.0023 & 0.641 \\
\addlinespace[2pt]
 & $\sigma = 0.001$ & the leading fits (13) & 16 & 0.998 & 0.762 & 0.062 & 0.956 & 13/16 & 16/16 & 0.0064 & 0.0303 & 0.1038 & 0.484 \\
 &  & the two named methods & 16 & 1.000 & 0.358 & 0.531 & 0.971 & 7/16 & 16/16 & 0.0131 & 0.0303 & 0.1038 & 1.000 \\
 &  & the classical variants (21) & 16 & 0.998 & 0.127 & 0.000 & 0.471 & 3/16 & 8/16 & 0.1026 & 0.0303 & 0.1038 & 3.729 \\
 &  & all accelerators (52) & 16 & 0.998 & 0.800 & 0.062 & 0.965 & 14/16 & 16/16 & 0.0064 & 0.0303 & 0.1038 & 0.399 \\
\addlinespace[2pt]
 & $\sigma = 0.005$ & the leading fits (13) & 16 & 1.000 & 0.679 & 0.125 & 0.952 & 13/16 & 16/16 & 0.0092 & 0.0296 & 0.1022 & 0.585 \\
 &  & the two named methods & 16 & 1.000 & 0.346 & 0.531 & 0.954 & 7/16 & 16/16 & 0.0152 & 0.0296 & 0.1022 & 1.000 \\
 &  & the classical variants (21) & 16 & 1.000 & 0.146 & 0.000 & 0.535 & 2/16 & 15/16 & 0.1001 & 0.0296 & 0.1022 & 3.536 \\
 &  & all accelerators (52) & 16 & 1.000 & 0.700 & 0.125 & 0.977 & 13/16 & 16/16 & 0.0092 & 0.0296 & 0.1022 & 0.565 \\
\addlinespace[2pt]
 & all cells & the leading fits (13) & 16 & 0.999 & 0.769 & 0.064 & 0.961 & 15/16 & 16/16 & 0.0077 & 0.0296 & 0.1030 & 0.391 \\
 &  & the two named methods & 16 & 1.000 & 0.343 & 0.552 & 0.965 & 6/16 & 16/16 & 0.0139 & 0.0296 & 0.1030 & 1.000 \\
 &  & the classical variants (21) & 16 & 0.999 & 0.341 & 0.000 & 0.631 & 2/16 & 15/16 & 0.0286 & 0.0296 & 0.1030 & 1.606 \\
 &  & all accelerators (52) & 16 & 0.999 & 0.804 & 0.062 & 0.973 & 15/16 & 16/16 & 0.0048 & 0.0296 & 0.1030 & 0.268 \\
\addlinespace[2pt]
\midrule
held-out & noise-free & the leading fits (13) & 5 & 1.000 & 0.933 & 0.000 & 1.000 & 5/5 & 5/5 & 0.0007 & 0.0053 & 0.0862 & 0.143 \\
 &  & the two named methods & 5 & 1.000 & 0.173 & 0.800 & 0.993 & 1/5 & 5/5 & 0.0053 & 0.0053 & 0.0862 & 1.000 \\
 &  & the classical variants (21) & 5 & 1.000 & 0.847 & 0.000 & 1.000 & 4/5 & 5/5 & 0.0014 & 0.0053 & 0.0862 & 0.145 \\
 &  & all accelerators (52) & 5 & 1.000 & 0.993 & 0.000 & 1.000 & 5/5 & 5/5 & 0.0004 & 0.0053 & 0.0862 & 0.075 \\
\addlinespace[2pt]
 & $\sigma = 0.001$ & the leading fits (13) & 5 & 0.993 & 0.773 & 0.000 & 0.953 & 5/5 & 5/5 & 0.0053 & 0.0074 & 0.0862 & 0.517 \\
 &  & the two named methods & 5 & 1.000 & 0.173 & 0.800 & 0.993 & 1/5 & 5/5 & 0.0074 & 0.0074 & 0.0862 & 1.000 \\
 &  & the classical variants (21) & 5 & 1.000 & 0.087 & 0.000 & 0.433 & 0/5 & 1/5 & 0.1023 & 0.0074 & 0.0862 & 4.257 \\
 &  & all accelerators (52) & 5 & 0.993 & 0.867 & 0.000 & 0.960 & 5/5 & 5/5 & 0.0026 & 0.0074 & 0.0862 & 0.517 \\
\addlinespace[2pt]
 & $\sigma = 0.005$ & the leading fits (13) & 5 & 1.000 & 0.527 & 0.400 & 0.973 & 3/5 & 5/5 & 0.0102 & 0.0102 & 0.0846 & 0.736 \\
 &  & the two named methods & 5 & 1.000 & 0.173 & 0.800 & 0.973 & 1/5 & 5/5 & 0.0102 & 0.0102 & 0.0846 & 1.000 \\
 &  & the classical variants (21) & 5 & 1.000 & 0.160 & 0.000 & 0.560 & 1/5 & 5/5 & 0.0847 & 0.0102 & 0.0846 & 3.107 \\
 &  & all accelerators (52) & 5 & 1.000 & 0.540 & 0.400 & 0.980 & 3/5 & 5/5 & 0.0061 & 0.0102 & 0.0846 & 0.357 \\
\addlinespace[2pt]
 & all cells & the leading fits (13) & 5 & 0.998 & 0.744 & 0.133 & 0.976 & 4/5 & 5/5 & 0.0053 & 0.0074 & 0.0862 & 0.511 \\
 &  & the two named methods & 5 & 1.000 & 0.173 & 0.800 & 0.987 & 1/5 & 5/5 & 0.0074 & 0.0074 & 0.0862 & 1.000 \\
 &  & the classical variants (21) & 5 & 1.000 & 0.364 & 0.000 & 0.664 & 1/5 & 5/5 & 0.0419 & 0.0074 & 0.0862 & 2.565 \\
 &  & all accelerators (52) & 5 & 0.998 & 0.800 & 0.133 & 0.980 & 5/5 & 5/5 & 0.0026 & 0.0074 & 0.0862 & 0.225 \\
\bottomrule
\end{tabular}

\end{adjustbox}
\end{table}
\end{landscape}
\begin{landscape}
\begin{table}[p]
\centering
\caption{\texttt{f27c\_selection\_by\_family.tex}: Choosing a method by trial family by family on the synthetic families (the leading fits (13)), headline stratum, both designs: per curve family the chosen method's win rate against the default and the last value, how often the default itself is chosen, the catastrophe rate and median error of the chosen method, next to the default's catastrophe rate and median error and the last value's median error}
\label{tab:f27c_selection_by_family}
\begin{adjustbox}{max width=\linewidth, max totalheight=0.86\textheight}
% AUTO-GENERATED -- do not hand-edit.  generated by scripts/make_paper_tables.py from the run started 2026-10-01T22:16:39 on code d82fd26 (full; 29,256 s; 7 workers)
% Choosing a method by trial family by family on the synthetic families (the leading fits (13)), headline stratum, both designs: per curve family the chosen method's win rate against the default and the last value, how often the default itself is chosen, the catastrophe rate and median error of the chosen method, next to the default's catastrophe rate and median error and the last value's median error
% source : results/phase1/phase1_selection_cells.csv (scripts/derive_selection.py on the git-ignored phase1_records.csv: one row per uncapped (set, family, noise, g) cell, candidate pool and design)
% source : results/phase1/phase1_aggregated.csv (internal check only: the default's family values reproduce f22's)
% selection: results/phase1/phase1_selection_provenance.json (analysis started 2026-10-03T18:14:05 on code b4d5518 from 369,360 records; cells per design {'one_pilot': 195, 'many_pilots': 195}, trials per design {'one_pilot': 169650, 'many_pilots': 5850}; added after the full run, evaluates no method)
% filter : target_g == 0.1; pool == lead (the leading fits (13): log_linear, richardson_3, richardson_2, double_exp_fit, pade_34, rational_fit, richardson_1, pade_23, pade_12, pade_45, pade_33, pade_22, single_exp_fit); one row per curve family (core, then held-out, the order of f22), cells = the family's uncapped (noise) cells; a family's rate = the mean over its cells of the per-cell share of trials (win vs default, default chosen, win vs last and the catastrophe rate of the chosen method, defined as in f27 / f27b), its median error = the median over its cells of the cell median (the aggregation of f22); designs: one pilot run = one_pilot (every ordered pair of distinct seeds (pilot, final); chosen = the pool member with the lowest error on the pilot seed among those valid there (ties: registry order); its record on the final seed is taken as it is (an invalid record stays invalid, no fallback)); all other seeds as pilots = many_pilots (every seed in turn is the final run and the other seeds are the pilots; chosen = the pool member with the lowest median error over its valid pilot records among the members valid on at least RANK_MIN_VALID of the pilot seeds (0.9; ties: registry order)); rational_fit = the default's catastrophe rate (mean over the family's cells of the share of the trials' final runs where its record is catastrophic) and its median error over the same final runs; last value = its median error over the same final runs; catastrophe rate (chosen, default) = mean over the cells of the share of trials whose chosen record (rational_fit's record) on the final run is catastrophic, the Phase 1 flag as recorded (the estimate is invalid, or its error exceeds CAT_MULT = 5 times the last value's on the same run); a trial without a choice counts as catastrophic, as it counts as invalid; a family with no uncapped cell prints 'capped'; the footers give, per set and design, the uncapped families where the win rate against the default exceeds 0.5 (the family win rate = mean over the family's cells of the per-cell share)
\begin{tabular}{llrrrrrrrrrrrrrr}
\toprule
Set & curve family & cells & \multicolumn{5}{c}{one pilot run} & \multicolumn{5}{c}{all other seeds as pilots} & \multicolumn{2}{c}{\meth{rational\_fit}} & last value \\
\cmidrule(lr){4-8}\cmidrule(lr){9-13}\cmidrule(lr){14-15}
 & & & win vs default & default chosen & win vs last & $\rhoC$ & med.\ err & win vs default & default chosen & win vs last & $\rhoC$ & med.\ err & $\rhoC$ & med.\ err & med.\ err \\
\midrule
core & \texttt{single\_exp} & 3 & 0.970 & 0.000 & 0.965 & 0.001 & 0.0004 & 1.000 & 0.000 & 1.000 & 0.000 & 0.0003 & 0.000 & 0.0275 & 0.0185 \\
core & \texttt{two\_exp} & 3 & 0.763 & 0.044 & 0.990 & 0.001 & 0.0057 & 0.933 & 0.000 & 1.000 & 0.000 & 0.0042 & 0.000 & 0.0275 & 0.0891 \\
core & \texttt{three\_exp} & 3 & 0.630 & 0.000 & 0.961 & 0.009 & 0.0270 & 0.622 & 0.000 & 0.933 & 0.011 & 0.0159 & 0.000 & 0.0339 & 0.1325 \\
core & \texttt{four\_exp} & 3 & 0.676 & 0.011 & 0.939 & 0.013 & 0.0220 & 0.900 & 0.000 & 1.000 & 0.000 & 0.0155 & 0.000 & 0.0329 & 0.1446 \\
core & \texttt{power\_law} & 3 & 0.674 & 0.133 & 0.969 & 0.004 & 0.0026 & 0.589 & 0.333 & 1.000 & 0.000 & 0.0032 & 0.000 & 0.0066 & 0.0306 \\
core & \texttt{rational\_decay} & 3 & 0.261 & 0.411 & 1.000 & 0.000 & 0.0020 & 0.233 & 0.667 & 1.000 & 0.000 & 0.0019 & 0.000 & 0.0019 & 0.1565 \\
core & \texttt{mixed\_pow\_rat} & 3 & 0.515 & 0.167 & 0.984 & 0.005 & 0.0142 & 0.644 & 0.000 & 1.000 & 0.000 & 0.0066 & 0.000 & 0.0148 & 0.1342 \\
core & \texttt{multiphase} & 3 & 0.967 & 0.000 & 0.993 & 0.000 & 0.0188 & 1.000 & 0.000 & 1.000 & 0.000 & 0.0167 & 0.000 & 0.1180 & 0.1792 \\
core & \texttt{osc\_exp} & 3 & 0.803 & 0.000 & 0.924 & 0.000 & 0.0030 & 0.956 & 0.000 & 1.000 & 0.000 & 0.0026 & 0.000 & 0.0363 & 0.0468 \\
core & \texttt{damped\_osc\_pow} & 3 & 0.739 & 0.111 & 0.960 & 0.007 & 0.0031 & 0.833 & 0.000 & 0.967 & 0.000 & 0.0026 & 0.000 & 0.0120 & 0.0501 \\
core & \texttt{log\_slow} & 0 & capped & capped & capped & capped & capped & capped & capped & capped & capped & capped & capped & capped & capped \\
core & \texttt{delayed\_plateau} & 3 & 0.589 & 0.111 & 0.977 & 0.023 & 0.0160 & 0.800 & 0.000 & 1.000 & 0.000 & 0.0142 & 0.000 & 0.0481 & 0.2197 \\
core & \texttt{staircase} & 3 & 0.470 & 0.000 & 0.999 & 0.001 & 0.0527 & 0.600 & 0.000 & 1.000 & 0.000 & 0.0328 & 0.000 & 0.0519 & 0.1800 \\
core & \texttt{noisy\_plateau} & 3 & 0.410 & 0.122 & 0.616 & 0.106 & 0.0010 & 0.667 & 0.022 & 0.744 & 0.044 & 0.0009 & 0.089 & 0.0008 & 0.0024 \\
core & \texttt{slow\_osc\_power} & 3 & 0.842 & 0.000 & 0.845 & 0.065 & 0.0088 & 0.911 & 0.000 & 0.944 & 0.000 & 0.0088 & 0.000 & 0.0627 & 0.0252 \\
core & \texttt{log\_oscillatory} & 0 & capped & capped & capped & capped & capped & capped & capped & capped & capped & capped & capped & capped & capped \\
core & \texttt{heavy\_noise\_plat} & 3 & 0.343 & 0.100 & 0.618 & 0.150 & 0.0015 & 0.611 & 0.000 & 0.789 & 0.056 & 0.0010 & 0.056 & 0.0011 & 0.0029 \\
core & \texttt{broken\_power\_law} & 3 & 0.796 & 0.000 & 0.840 & 0.000 & 0.0605 & 1.000 & 0.000 & 1.000 & 0.000 & 0.0604 & 0.000 & 0.0729 & 0.1185 \\
\addlinespace[2pt]
held-out & \texttt{stretched\_exp} & 3 & 0.347 & 0.178 & 0.968 & 0.002 & 0.0083 & 0.456 & 0.333 & 1.000 & 0.000 & 0.0053 & 0.000 & 0.0074 & 0.1406 \\
held-out & \texttt{logistic\_tail} & 3 & 0.923 & 0.000 & 0.980 & 0.000 & 0.0287 & 0.911 & 0.000 & 0.978 & 0.000 & 0.0265 & 0.000 & 0.1076 & 0.0862 \\
held-out & \texttt{inv\_sqrt\_log} & 3 & 0.626 & 0.178 & 0.941 & 0.010 & 0.0025 & 0.600 & 0.333 & 0.967 & 0.011 & 0.0021 & 0.011 & 0.0044 & 0.0187 \\
held-out & \texttt{random\_knots} & 0 & capped & capped & capped & capped & capped & capped & capped & capped & capped & capped & capped & capped & capped \\
held-out & \texttt{real\_boot\_a} & 3 & 0.681 & 0.000 & 0.973 & 0.011 & 0.0323 & 0.867 & 0.000 & 0.933 & 0.033 & 0.0221 & 0.000 & 0.0488 & 0.1502 \\
held-out & \texttt{real\_boot\_b} & 3 & 0.688 & 0.044 & 0.946 & 0.010 & 0.0016 & 0.889 & 0.000 & 1.000 & 0.000 & 0.0010 & 0.000 & 0.0049 & 0.0460 \\
\midrule
\multicolumn{16}{l}{core, one pilot run: the win rate against the default exceeds 0.5 in 12 of 16 uncapped families; not in \texttt{rational\_decay} (0.261), \texttt{staircase} (0.470), \texttt{noisy\_plateau} (0.410), \texttt{heavy\_noise\_plat} (0.343)} \\
\multicolumn{16}{l}{core, all other seeds as pilots: the win rate against the default exceeds 0.5 in 15 of 16 uncapped families; not in \texttt{rational\_decay} (0.233)} \\
\multicolumn{16}{l}{held-out, one pilot run: the win rate against the default exceeds 0.5 in 4 of 5 uncapped families; not in \texttt{stretched\_exp} (0.347)} \\
\multicolumn{16}{l}{held-out, all other seeds as pilots: the win rate against the default exceeds 0.5 in 4 of 5 uncapped families; not in \texttt{stretched\_exp} (0.456)} \\
\bottomrule
\end{tabular}

\end{adjustbox}
\end{table}
\end{landscape}
\begin{landscape}
\begin{table}[p]
\centering
\caption{\texttt{f28\_selection\_recorded.tex}: Choosing a method by trial on the recorded curves, cells before the minimum: leave one depth out and leave one dataset out, per candidate pool, the chosen method's score against the default's and against the last value}
\label{tab:f28_selection_recorded}
\begin{adjustbox}{max width=\linewidth, max totalheight=0.86\textheight}
% AUTO-GENERATED -- do not hand-edit.  generated by scripts/make_paper_tables.py from the run started 2026-10-01T22:16:39 on code d82fd26 (full; 29,256 s; 7 workers)
% Choosing a method by trial on the recorded curves, cells before the minimum: leave one depth out and leave one dataset out, per candidate pool, the chosen method's score against the default's and against the last value
% source : results/real_data/real_data_roster_v2.csv (the recorded curves, every method under the benchmark protocol)
% roster  : results/real_data/real_data_roster_provenance.json (evaluation started 2026-10-03T13:02:45 on code 7518fad, 5040 rows; added after the full run)
% filter : post_min_target == 0 (the 65 cells before the curve minimum, datasets adult, covertype, higgs, jannis, miniboone); a method's score on a set of cells = the median of skill_vs_last (error / the last value's error) over those cells with an invalid record counted as +infinity; same curve, leave one depth out = for each dataset and observation depth, the chosen method is the pool member with the lowest score on the pre-minimum cells of the other depths of that dataset (ties: registry order), evaluated by its score on the left-out depth's pre-minimum cells next to rational_fit's score there; another dataset, leave one dataset out = chosen on the pre-minimum cells of the other datasets, evaluated on the left-out dataset; below / equal / above = folds where the chosen score is below / equal to / above the default's; chosen < 1 = folds where the chosen method beats the last value at the median; candidate pools as in f27 (lead: log_linear, richardson_3, richardson_2, double_exp_fit, pade_34, rational_fit, richardson_1, pade_23, pade_12, pade_45, pade_33, pade_22, single_exp_fit; named: rational_fit, single_exp_fit; classical: the 21 variants; all: 52 accelerators)
\begin{tabular}{llrrrrrrrl}
\toprule
Design & candidate pool & folds & \multicolumn{3}{c}{chosen score vs default's} & chosen $< 1$ & \multicolumn{2}{c}{median score over folds} & methods chosen \\
\cmidrule(lr){4-6}\cmidrule(lr){8-9}
 & & & below & equal & above & & chosen & default & \\
\midrule
same curve, leave one depth out & the leading fits (13) & 25 & 14 & 2 & 9 & 24 & 0.216 & 0.333 & \meth{pade\_33} (6), \meth{log\_linear} (5), \meth{richardson\_1} (5), \meth{double\_exp\_fit} (5), \meth{rational\_fit} (2), \meth{pade\_22} (2) \\
 & the two named methods & 25 & 4 & 20 & 1 & 25 & 0.317 & 0.333 & \meth{rational\_fit} (20), \meth{single\_exp\_fit} (5) \\
 & the classical variants (21) & 25 & 4 & 0 & 21 & 11 & 1.062 & 0.333 & \meth{wynn\_eps\_3} (5), \meth{wynn\_rho\_2} (5), \meth{wynn\_rho\_3} (4), \meth{brezinski\_theta2} (4), \meth{wynn\_rho\_1} (3), \meth{shanks\_4} (2), \meth{brezinski\_theta1} (2) \\
 & all accelerators (52) & 25 & 16 & 0 & 9 & 23 & 0.200 & 0.333 & \meth{log\_linear} (5), \meth{richardson\_1} (5), \meth{double\_exp\_fit} (5), \meth{pade\_33} (4), \meth{pade\_44} (4), \meth{richardson\_a10} (2) \\
\midrule
another dataset, leave one dataset out & the leading fits (13) & 5 & 1 & 2 & 2 & 5 & 0.344 & 0.270 & \meth{rational\_fit} (2), \meth{richardson\_1} (1), \meth{pade\_22} (1), \meth{pade\_33} (1) \\
 & the two named methods & 5 & 0 & 5 & 0 & 5 & 0.270 & 0.270 & \meth{rational\_fit} (5) \\
 & the classical variants (21) & 5 & 0 & 0 & 5 & 1 & 1.220 & 0.270 & \meth{wynn\_eps\_3} (4), \meth{wynn\_rho\_1} (1) \\
 & all accelerators (52) & 5 & 3 & 1 & 1 & 5 & 0.263 & 0.270 & \meth{pade\_11} (3), \meth{richardson\_1} (1), \meth{rational\_fit} (1) \\
\bottomrule
\end{tabular}

\end{adjustbox}
\end{table}
\end{landscape}
\begin{landscape}
\begin{table}[p]
\centering
\caption{\texttt{f29\_real\_named\_by\_dataset.tex}: The two named methods on the recorded curves, per dataset and pooled: cells, wins against the last value and the median error ratio before and after the curve minimum}
\label{tab:f29_real_named_by_dataset}
\begin{adjustbox}{max width=\linewidth, max totalheight=0.86\textheight}
% AUTO-GENERATED -- do not hand-edit.  generated by scripts/make_paper_tables.py from the run started 2026-10-01T22:16:39 on code d82fd26 (full; 29,256 s; 7 workers)
% The two named methods on the recorded curves, per dataset and pooled: cells, wins against the last value and the median error ratio before and after the curve minimum
% source : results/real_data/real_data_roster_v2.csv (the recorded curves, every method under the benchmark protocol)
% roster  : results/real_data/real_data_roster_provenance.json (evaluation started 2026-10-03T13:02:45 on code 7518fad, 5040 rows; added after the full run)
% filter : method in ('rational_fit', 'single_exp_fit'); one row per dataset and a pooled row; before / after = post_min_target 0 / 1 (target round beyond the recorded curve's argmin); cells = (depth, target) cells of the split; wins = cells with win_vs_last == 1 (k/n; an invalid record never wins); ratio = median over the valid cells of skill_vs_last (the method's error / the last value's); a split with no cell prints --
\begin{tabular}{lrrrrrrrrrr}
\toprule
Dataset & \multicolumn{2}{c}{cells} & \multicolumn{4}{c}{\meth{rational\_fit}} & \multicolumn{4}{c}{\meth{single\_exp\_fit}} \\
\cmidrule(lr){2-3}\cmidrule(lr){4-7}\cmidrule(lr){8-11}
 & before & after & wins before & ratio before & wins after & ratio after & wins before & ratio before & wins after & ratio after \\
\midrule
\texttt{adult} & 5 & 10 & 5/5 & 0.721 & 4/10 & 1.077 & 5/5 & 0.585 & 10/10 & 0.580 \\
\texttt{bank\_marketing} & 0 & 15 & -- & -- & 9/15 & 0.758 & -- & -- & 12/15 & 0.777 \\
\texttt{covertype} & 15 & 0 & 15/15 & 0.270 & -- & -- & 15/15 & 0.753 & -- & -- \\
\texttt{higgs} & 15 & 0 & 15/15 & 0.199 & -- & -- & 15/15 & 0.614 & -- & -- \\
\texttt{jannis} & 15 & 0 & 15/15 & 0.369 & -- & -- & 15/15 & 0.553 & -- & -- \\
\texttt{miniboone} & 15 & 0 & 15/15 & 0.233 & -- & -- & 15/15 & 0.635 & -- & -- \\
\midrule
\textbf{pooled} & 65 & 25 & 65/65 & 0.270 & 13/25 & 0.783 & 65/65 & 0.628 & 22/25 & 0.689 \\
\bottomrule
\end{tabular}

\end{adjustbox}
\end{table}
\end{landscape}
\begin{landscape}
\begin{table}[p]
\centering
\caption{\texttt{f30\_real\_final\_loss\_by\_depth.tex}: Final loss by observation depth on the 4 recorded curves whose minimum is at the final round (covertype, higgs, jannis, miniboone): the error of the last value, rational\_fit and single\_exp\_fit as a percentage of the final loss, median and maximum over the datasets}
\label{tab:f30_real_final_loss_by_depth}
\begin{adjustbox}{max width=\linewidth, max totalheight=0.86\textheight}
% AUTO-GENERATED -- do not hand-edit.  generated by scripts/make_paper_tables.py from the run started 2026-10-01T22:16:39 on code d82fd26 (full; 29,256 s; 7 workers)
% Final loss by observation depth on the 4 recorded curves whose minimum is at the final round (covertype, higgs, jannis, miniboone): the error of the last value, rational_fit and single_exp_fit as a percentage of the final loss, median and maximum over the datasets
% source : results/real_data/real_data_roster_v2.csv (the recorded curves, every method under the benchmark protocol)
% source : results/real_data/real_data_curves.csv (the recorded curves: the final round)
% roster  : results/real_data/real_data_roster_provenance.json (evaluation started 2026-10-03T13:02:45 on code 7518fad, 5040 rows; added after the full run)
% filter : datasets whose curve minimum is at the final recorded round (covertype, higgs, jannis, miniboone: argmin_round == 500 == the largest target round == the last recorded round); target_round == 500; one row per observation depth; share of rounds observed = depth / 500; per method the median and the maximum over those 4 datasets of 100 x error / true_val (the recorded loss at the target round), over the valid records; when a record is invalid the count of valid ones is printed beside the value
\begin{tabular}{lrrrrrrr}
\toprule
$\nobs$ & share of rounds observed & \multicolumn{2}{c}{\meth{last\_value}} & \multicolumn{2}{c}{\meth{rational\_fit}} & \multicolumn{2}{c}{\meth{single\_exp\_fit}} \\
\cmidrule(lr){3-4}\cmidrule(lr){5-6}\cmidrule(lr){7-8}
 & & median & max & median & max & median & max \\
\midrule
30 & 6\% & 52.21\% & 94.60\% & 13.80\% & 39.41\% & 32.91\% & 74.82\% \\
60 & 12\% & 24.42\% & 60.77\% & 10.33\% & 26.34\% & 21.49\% & 56.19\% \\
90 & 18\% & 14.38\% & 46.77\% & 3.27\% & 18.12\% & 9.59\% & 37.30\% \\
120 & 24\% & 9.64\% & 38.84\% & 2.65\% & 18.39\% & 5.23\% & 29.27\% \\
150 & 30\% & 6.95\% & 33.51\% & 1.10\% & 15.90\% & 3.87\% & 22.95\% \\
\midrule
\multicolumn{8}{l}{datasets whose minimum is at the final recorded round (500): \texttt{covertype}, \texttt{higgs}, \texttt{jannis}, \texttt{miniboone}; target = round 500; share of rounds observed = $\nobs / 500$} \\
\bottomrule
\end{tabular}

\end{adjustbox}
\end{table}
\end{landscape}
\end{document}
